% Les activités artisanales — Géographie 1re Tech — Cours signé OnBuch+
% Programme officiel MINESEC. Compiler avec : tectonic N05-activites-artisanales.tex
\newcommand{\DOCMATIERE}{Géographie}
\newcommand{\DOCNIVEAU}{1re Tech}
\newcommand{\DOCMODULE}{Géographie – classe de Première technique}
\newcommand{\DOCLECON}{N05}
\newcommand{\DOCTITRE}{Les activités artisanales}
% OnBuch+ — préambule « cours signés » ENRICHI (compilé avec tectonic / XeLaTeX)
% Système visuel « Pop » d'OnBuch+ : fond crème (jamais blanc), bordures encre
% épaisses, ombres « dures » décalées, titres Archivo Black en capitales.
% Version MATHÉMATIQUES : graphes (pgfplots/tikz), boîtes pédagogiques
% (définition, propriété, méthode, attention, le savais-tu, exercice résolu,
% exercice, TP, bilan), page de garde et signature OnBuch+ ancrée en bas.
%
% Macros attendues dans main.tex AVANT \input{preamble} :
%   \DOCMATIERE  \DOCNIVEAU  \DOCMODULE  \DOCLECON (numéro)  \DOCTITRE
\documentclass[11pt]{article}

\usepackage{fontspec}
\usepackage[french]{babel}
\usepackage[a4paper,margin=2cm,top=3.4cm,bottom=2.8cm,headheight=1.4cm,footskip=1.3cm]{geometry}
\usepackage{amsmath,amssymb,amsfonts}
\usepackage{mathtools} % \xmapsto, \coloneqq... souvent utilisés par les modèles
\usepackage{cancel} % simplifications barrées (bilans de forces)
% \abs{z} (module, valeur absolue) et \norm{u} (norme), utilisés par les modèles
\providecommand{\abs}{}\let\abs\relax\DeclarePairedDelimiter{\abs}{\lvert}{\rvert}
\providecommand{\norm}{}\let\norm\relax\DeclarePairedDelimiter{\norm}{\lVert}{\rVert}
\usepackage[table]{xcolor} % [table] : fournit aussi \rowcolors (alternance de lignes)
\usepackage{pagecolor}
\usepackage{tikz}
\usetikzlibrary{arrows.meta,positioning,calc,shapes.geometric,shapes.misc,shapes.arrows,decorations.pathmorphing,decorations.pathreplacing,decorations.markings,patterns,shadows,mindmap,trees,fit,backgrounds,matrix,intersections,angles,quotes}
\usepackage{pgfplots}
\usepackage[version=4]{mhchem} % équations nucléaires \ce{^{235}_{92}U}
\usepackage[european,siunitx]{circuitikz} % schémas électriques
\pgfplotsset{compat=1.18}
\usepgfplotslibrary{fillbetween,groupplots} % aires entre courbes (calcul intégral)
\usepackage{enumitem}
\usepackage{fancyhdr}
% [most] charge toutes les bibliothèques tcolorbox (listings, minted, raster,
% documentation...) et gonfle inutilement la mémoire TeX sur un document long
% (~300 boîtes) : on ne charge que ce qui est réellement utilisé.
\usepackage[skins,breakable]{tcolorbox}
\usepackage{graphicx}
\usepackage{colortbl}
\usepackage{booktabs}
\usepackage{tabularx}
\usepackage{multirow}
\usepackage{diagbox} % case d'en-tête diagonale des tableaux à double entrée
% Colonnes extensibles centrée / gauche / droite (C, L, R), souvent utilisées par les modèles
\newcolumntype{C}{>{\centering\arraybackslash}X}
\newcolumntype{L}{>{\raggedright\arraybackslash}X}
\newcolumntype{R}{>{\raggedleft\arraybackslash}X}
\usepackage{array}
\usepackage{multicol}
\usepackage{siunitx}
% parse-numbers=false : accepte aussi \SI{2,5 \times 10^{-3}}{...}
\sisetup{locale=FR,output-decimal-marker={,},group-minimum-digits=4,parse-numbers=false}
\DeclareSIUnit\ohm{\ensuremath{\symup{\Omega}}} % U+2126 absent de la police
\DeclareSIUnit\division{div} % réglages d'oscilloscope (ms/div, V/div)
\DeclareSIUnit\tr{tr} % tours (vitesse de rotation, ex. \SI{1500}{\tr\per\minute})
\DeclareSIUnit\molar{M} % concentration molaire (ex. \SI{3}{\molar}, \SI{1}{\micro\molar})
\DeclareSIUnit\an{an} % année (le modèle confond parfois avec \year, un compteur TeX) — ex. \SI{5}{\kilo\meter\per\an}
\DeclareSIUnit\FCFA{FCFA} % franc CFA (ex. \SI{500}{\FCFA\per\kilo\gram})
\DeclareSIUnit\degreeBrix{\textdegree Brix} % teneur en sucre d'un sirop/jus (réfractométrie)
\DeclareSIUnit\month{mois} % le modèle utilise parfois \month, un compteur TeX, comme unité de durée
\DeclareSIUnit\microliter{\micro\liter} % le modèle écrit parfois \microliter en un seul token
\DeclareSIUnit\million{millions} % dénombrement (ex. \SI{15}{\million\per\mL} spermatozoïdes)
\usepackage{qrcode}
% \degree : commande fréquemment utilisée par les modèles pour un angle ou une
% température en dehors de \SI{}{} (non fournie par les paquets chargés).
\providecommand{\degree}{^{\circ}}
% \qed (fin de démonstration), utilisé par les modèles sans amsthm
\providecommand{\qed}{\hfill$\square$}
% \barre : double barre des valeurs interdites dans un tableau de variations (tkz-tab)
\providecommand{\barre}{\ensuremath{\|}}
% environnement proof (amsthm non chargé) utilisé par les modèles
\ifdefined\proof\else\newenvironment{proof}[1][Démonstration]{\par\noindent\textit{#1.}\ }{\qed\par}\fi
\usepackage{lastpage}
\usepackage{hyperref}
\hypersetup{hidelinks}

% ============================================================
% Palette « Pop » OnBuch+ — mêmes hex que la classe P (plus_ui.dart)
% ============================================================
\definecolor{poporange}{HTML}{F59321}
\definecolor{popdark}{HTML}{E07A0C}
\definecolor{popcream}{HTML}{FFF6E8}
\definecolor{popink}{HTML}{1C1714}
\definecolor{poppurple}{HTML}{7A5AE0}
\definecolor{popgreen}{HTML}{2E9E6A}
\definecolor{poppink}{HTML}{E0457B}
\definecolor{popblue}{HTML}{2F7DE1}
\definecolor{popgold}{HTML}{C98A12}
\definecolor{popmuted}{HTML}{8A7F76}
\definecolor{popline}{HTML}{EADFD2}
\definecolor{popgray}{HTML}{8A7F76}
\colorlet{popgrey}{popgray}
% Alias tolérants (le modèle nomme parfois les couleurs différemment)
\colorlet{popwhite}{white}
\colorlet{popblack}{popink}
\colorlet{popred}{poppink}
\colorlet{popyellow}{popgold}
% Teintes claires (fonds de tableaux/encadrés) — le modèle nomme parfois
% les couleurs avec un suffixe L (ex. popblueL) sans les avoir définies.
\colorlet{poporangeL}{poporange!20}
\colorlet{popdarkL}{popdark!20}
\colorlet{poppurpleL}{poppurple!20}
\colorlet{popgreenL}{popgreen!20}
\colorlet{poppinkL}{poppink!20}
\colorlet{popblueL}{popblue!20}
\colorlet{popgoldL}{popgold!20}
\colorlet{popredL}{popred!20}
\colorlet{popgrayL}{popgray!20}
% Teintes claires pour fonds de tableaux / zones
\colorlet{popblueL}{popblue!10!white}
\colorlet{poporangeL}{poporange!14!white}
\colorlet{popgreenL}{popgreen!12!white}
\colorlet{poppurpleL}{poppurple!12!white}
\colorlet{poppinkL}{poppink!12!white}
\colorlet{popgoldL}{popgold!14!white}

\pagecolor{popcream}

% ============================================================
% Typographie
% ============================================================
\defaultfontfeatures{Ligatures=TeX}
\setmainfont{PlusJakartaSans-Medium.ttf}[Path=../../../fonts/,BoldFont=PlusJakartaSans-Bold.ttf]
% Maths en sans-serif (Fira Math) avec chiffres et lettres droites de Plus
% Jakarta, pour que les formules se fondent dans le texte.
\usepackage[math-style=ISO,bold-style=ISO]{unicode-math}
\setmathfont{FiraMath-Regular.otf}[Path=../../../fonts/]
\setmathfont{PlusJakartaSans-Medium.ttf}[Path=../../../fonts/,range={up/num,up/{Latin,latin},\mathup}]
\usepackage{icomma} % virgule décimale sans espace en mode math
% Caractères mathématiques Unicode absents des polices de texte (Plus Jakarta,
% Archivo Black) : rendus par leur équivalent mathématique, en texte comme en
% formule — sinon ils s'affichent en carré vide (ex. « Arithmétique dans ℤ »).
\usepackage{newunicodechar}
\newunicodechar{ℝ}{\ensuremath{\mathbb{R}}}
\newunicodechar{ℤ}{\ensuremath{\mathbb{Z}}}
\newunicodechar{ℂ}{\ensuremath{\mathbb{C}}}
\newunicodechar{ℕ}{\ensuremath{\mathbb{N}}}
\newunicodechar{ℚ}{\ensuremath{\mathbb{Q}}}
\newunicodechar{𝔻}{\ensuremath{\mathbb{D}}}
\newunicodechar{α}{\ensuremath{\alpha}}
\newunicodechar{β}{\ensuremath{\beta}}
\newunicodechar{γ}{\ensuremath{\gamma}}
\newunicodechar{δ}{\ensuremath{\delta}}
\newunicodechar{ε}{\ensuremath{\varepsilon}}
\newunicodechar{θ}{\ensuremath{\theta}}
\newunicodechar{λ}{\ensuremath{\lambda}}
\newunicodechar{μ}{\ensuremath{\mu}}
\newunicodechar{σ}{\ensuremath{\sigma}}
\newunicodechar{τ}{\ensuremath{\tau}}
\newunicodechar{φ}{\ensuremath{\varphi}}
\newunicodechar{ψ}{\ensuremath{\psi}}
\newunicodechar{ω}{\ensuremath{\omega}}
\newunicodechar{Σ}{\ensuremath{\Sigma}}
% majuscules produites par \MakeUppercase dans les titres (α-aminés -> Α)
\newunicodechar{Α}{\ensuremath{\alpha}}
\newunicodechar{Β}{\ensuremath{\beta}}
\newunicodechar{Γ}{\ensuremath{\Gamma}}
\newunicodechar{Δ}{\ensuremath{\Delta}}
\newunicodechar{Ε}{\ensuremath{\varepsilon}}
\newunicodechar{Θ}{\ensuremath{\Theta}}
\newunicodechar{Λ}{\ensuremath{\Lambda}}
\newunicodechar{Μ}{\ensuremath{\mu}}
\newunicodechar{Τ}{\ensuremath{\tau}}
\newunicodechar{Φ}{\ensuremath{\Phi}}
\newunicodechar{Ψ}{\ensuremath{\Psi}}
\newunicodechar{Ω}{\ensuremath{\Omega}}
\newunicodechar{↦}{\ensuremath{\mapsto}}
\newunicodechar{∈}{\ensuremath{\in}}
\newunicodechar{∉}{\ensuremath{\notin}}
\newunicodechar{∧}{\ensuremath{\wedge}}
\newunicodechar{⇔}{\ensuremath{\Leftrightarrow}}
\newunicodechar{⟺}{\ensuremath{\Longleftrightarrow}}
\newunicodechar{⇒}{\ensuremath{\Rightarrow}}
\newunicodechar{⟹}{\ensuremath{\Longrightarrow}}
\newunicodechar{∘}{\ensuremath{\circ}}
\newunicodechar{∪}{\ensuremath{\cup}}
\newunicodechar{∩}{\ensuremath{\cap}}
\newunicodechar{≡}{\ensuremath{\equiv}}
\newunicodechar{⊂}{\ensuremath{\subset}}
\newunicodechar{⊥}{\ensuremath{\perp}}
\newunicodechar{∃}{\ensuremath{\exists}}
\newunicodechar{∀}{\ensuremath{\forall}}
\newunicodechar{∛}{\ensuremath{\sqrt[3]{\,}}}
\newunicodechar{∜}{\ensuremath{\sqrt[4]{\,}}}
\newunicodechar{⃗}{\ensuremath{{}^{\to}}}
\newunicodechar{ⁿ}{\ifmmode^{n}\else\textsuperscript{\ensuremath{n}}\fi}
\newunicodechar{ˣ}{\ifmmode^{x}\else\textsuperscript{\ensuremath{x}}\fi}
\newunicodechar{ᵇ}{\ifmmode^{b}\else\textsuperscript{\ensuremath{b}}\fi}
\newunicodechar{ʳ}{\ifmmode^{r}\else\textsuperscript{\ensuremath{r}}\fi}
\newunicodechar{ᵘ}{\ifmmode^{u}\else\textsuperscript{\ensuremath{u}}\fi}
\newunicodechar{ʸ}{\ifmmode^{y}\else\textsuperscript{\ensuremath{y}}\fi}
\newunicodechar{ᵃ}{\ifmmode^{a}\else\textsuperscript{\ensuremath{a}}\fi}
\newunicodechar{ᵉ}{\ifmmode^{e}\else\textsuperscript{\ensuremath{e}}\fi}
\newunicodechar{ᵏ}{\ifmmode^{k}\else\textsuperscript{\ensuremath{k}}\fi}
\newunicodechar{ᵖ}{\ifmmode^{p}\else\textsuperscript{\ensuremath{p}}\fi}
\newunicodechar{ᴺ}{\ifmmode^{N}\else\textsuperscript{\ensuremath{N}}\fi}
\newunicodechar{ᶜ}{\ifmmode^{c}\else\textsuperscript{\ensuremath{c}}\fi}
\newunicodechar{ʰ}{\ifmmode^{h}\else\textsuperscript{\ensuremath{h}}\fi}
\newunicodechar{ᵠ}{\ifmmode^{q}\else\textsuperscript{\ensuremath{q}}\fi}
\newunicodechar{ₙ}{\ifmmode_{n}\else\textsubscript{\ensuremath{n}}\fi}
\newunicodechar{ₐ}{\ifmmode_{a}\else\textsubscript{\ensuremath{a}}\fi}
\newunicodechar{ᵢ}{\ifmmode_{i}\else\textsubscript{\ensuremath{i}}\fi}
\newunicodechar{ᵣ}{\ifmmode_{r}\else\textsubscript{\ensuremath{r}}\fi}
\newunicodechar{ₖ}{\ifmmode_{k}\else\textsubscript{\ensuremath{k}}\fi}
\newunicodechar{ᵦ}{\ifmmode_{\beta}\else\textsubscript{\ensuremath{\beta}}\fi}
\newunicodechar{ₓ}{\ifmmode_{x}\else\textsubscript{\ensuremath{x}}\fi}
\newfontfamily\popdisplay{ArchivoBlack-Regular.ttf}[Path=../../../fonts/]
\newfontfamily\poplogo{SpaceGrotesk-Bold.ttf}[Path=../../../fonts/]
\newfontfamily\popsemi{PlusJakartaSans-SemiBold.ttf}[Path=../../../fonts/]
\newfontfamily\popxbold{PlusJakartaSans-ExtraBold.ttf}[Path=../../../fonts/]
\newfontfamily\popxboldls{PlusJakartaSans-ExtraBold.ttf}[Path=../../../fonts/,LetterSpace=3.0]

\setlist[enumerate]{leftmargin=1.7em,itemsep=5pt,topsep=4pt}
\setlist[itemize]{leftmargin=1.7em,itemsep=4pt,topsep=4pt,label={\color{poporange}\textbullet}}
\setlength{\parindent}{0pt}
\setlength{\parskip}{5pt}
\renewcommand{\arraystretch}{1.25}

% pgfplots : style Pop par défaut
\pgfplotsset{
  every axis/.append style={axis line style={popink,line width=1pt},tick style={popink},
    grid style={popline,line width=0.5pt},label style={font=\small},tick label style={font=\footnotesize}},
  popcurve/.style={poporange,line width=1.8pt,no markers,smooth},
}
% Même style disponible dans un \draw TikZ simple (hors pgfplots)
\tikzset{popcurve/.style={poporange,line width=1.8pt}}
\tikzset{popgrid/.style={popline,very thin}}
\colorlet{popcurve}{poporange} % le nom du style est parfois utilisé comme couleur

% ============================================================
% Pill / squiggle
% ============================================================
\newcommand{\poppill}[3]{%
  \tikz[baseline=-0.55ex]\node[draw=popink,line width=1.2pt,rounded corners=2.6pt,%
    fill=#2,inner xsep=6.5pt,inner ysep=3pt]{%
    \popxboldls\fontsize{7}{7}\selectfont\color{#3}\MakeUppercase{#1}};%
}
\newcommand{\popsquiggle}[1]{%
  \tikz[baseline=-0.4ex]{
    \draw[#1,line width=2.6pt,line cap=round]
      plot [smooth] coordinates {(0,0.09) (0.34,0.01) (0.68,0.21) (1.02,0.01) (1.36,0.10)};
  }%
}
% Mot-clé surligné (marqueur orange sous le texte)
\newcommand{\cle}[1]{{\popxbold\color{popdark}#1}}

% ============================================================
% En-tête / pied
% ============================================================
\pagestyle{fancy}
\fancyhf{}
\renewcommand{\headrulewidth}{0pt}
\renewcommand{\footrulewidth}{0pt}
\lhead{\raisebox{-0.15ex}{\poplogo\fontsize{13}{13}\selectfont\color{popink}OnBuch\color{poporange}+}}
\rhead{\poppill{\DOCMATIERE\ \textbullet\ \DOCNIVEAU\ \textbullet\ Leçon \DOCLECON}{white}{popink}}
\lfoot{\popxbold\fontsize{7.6}{7.6}\selectfont\color{popmuted}\MakeUppercase{Cours signé OnBuch+}\ \textbullet\ {\popsemi\DOCTITRE}}
\rfoot{\tikz[baseline=-0.6ex]\node[rounded corners=8.5pt,fill=popink,minimum height=17pt,minimum width=17pt,inner xsep=5pt,inner sep=0]{\popxbold\fontsize{8}{8}\selectfont\color{popcream}\thepage\,/\,\pageref*{LastPage}};}

% ============================================================
% Titres
% ============================================================
\newcommand{\chaptitre}[2]{%
  \begin{tcolorbox}[enhanced,colback=poporange,colframe=popink,boxrule=2.6pt,arc=11pt,
      drop shadow=popink,left=15pt,right=15pt,top=12pt,bottom=11pt,before skip=3pt,after skip=18pt]
    \poppill{#2}{popink}{popcream}\par\vspace{9pt}
    {\raggedright\hyphenpenalty=10000\exhyphenpenalty=10000\popdisplay\fontsize{22}{24}\selectfont\color{popcream}\MakeUppercase{#1}\par}
  \end{tcolorbox}
}
% Section numérotée : pastille orange avec le numéro + titre Archivo
\newcounter{popsec}
\newcommand{\coursec}[1]{%
  \stepcounter{popsec}\setcounter{popsub}{0}%
  \par\vspace{16pt}\noindent
  \begin{minipage}{\linewidth}
  \tikz[baseline=-0.7ex]\node[draw=popink,line width=1.6pt,fill=poporange,rounded corners=4pt,minimum size=22pt,inner sep=0]{\popdisplay\fontsize{12}{12}\selectfont\color{popcream}\thepopsec};%
  \hspace{8pt}{\popdisplay\fontsize{15}{17}\selectfont\color{popink}\MakeUppercase{#1}}\par
  \vspace{4pt}\noindent{\color{poporange}\rule{36pt}{3.6pt}}
  \end{minipage}\par\nopagebreak\vspace{8pt}
}
\newcounter{popsub}
\newcommand{\courssub}[1]{%
  \stepcounter{popsub}%
  \par\vspace{9pt}\noindent{\popxbold\fontsize{11.5}{13}\selectfont\color{popink}{\color{poporange}\thepopsec.\thepopsub}\hspace{6pt}#1}\par\nopagebreak\vspace{4pt}
}

% ============================================================
% Boîtes pédagogiques — même recette : fond blanc, bordure épaisse colorée,
% ombre dure encre, badge plein. Titre optionnel [#1].
% ============================================================
\tcbset{popbase/.style={breakable,enhanced,colback=white,boxrule=2.3pt,arc=9pt,
  shadow={3pt}{-3pt}{0pt}{popink},left=14pt,right=14pt,top=11pt,bottom=11pt,before skip=11pt,after skip=15pt,
  fontupper=\popsemi\fontsize{10.6}{14.5}\selectfont\color{popink},
  fontlower=\popsemi\fontsize{10.6}{14.5}\selectfont\color{popink}}}
% Coche dessinée (le glyphe ✓ n'existe pas dans les polices du document)
\newcommand{\popcheck}[1]{\tikz[baseline=-0.5ex]\draw[#1,line width=1.6pt,line cap=round,line join=round](-0.13,0.0)--(-0.04,-0.09)--(0.13,0.1);}
\providecommand{\coche}{\popcheck{popgreen}}
\providecommand{\resultat}[1]{\par\smallskip\noindent\fcolorbox{popgreen}{popgreenL}{\parbox{\dimexpr\linewidth-2\fboxsep-2\fboxrule\relax}{\textbf{Résultat.} #1}}\par\smallskip}
\newcommand{\popboxhead}[3]{\poppill{#1}{#2}{white}\ifx\relax#3\relax\else\hspace{6pt}{\popxbold\fontsize{10.6}{12}\selectfont\color{popink}#3}\fi\par\vspace{7pt}}

\newtcolorbox{aretenir}[1][]{popbase,colframe=popgreen,before upper={\popboxhead{À retenir}{popgreen}{#1}}}
\newtcolorbox{exemplebox}[1][]{popbase,colframe=poporange,before upper={\popboxhead{Exemple}{poporange}{#1}}}
\newtcolorbox{definition}[1][]{popbase,colframe=popblue,before upper={\popboxhead{Définition}{popblue}{#1}}}
\newtcolorbox{propriete}[1][]{popbase,colframe=poppurple,before upper={\popboxhead{Propriété}{poppurple}{#1}}}
\newtcolorbox{methode}[1][]{popbase,colframe=popgold,before upper={\popboxhead{Méthode}{popgold}{#1}}}
\newtcolorbox{attention}[1][]{popbase,colframe=poppink,before upper={\popboxhead{Attention}{poppink}{#1}}}
\newtcolorbox{experience}[1][]{popbase,colframe=popblue,colback=popblueL,before upper={\popboxhead{Expérience}{popblue}{#1}}}
% « Le savais-tu ? » : carte violette pleine (contexte, culture, Cameroun)
\newtcolorbox{savaistu}[1][]{popbase,colback=poppurple,colframe=popink,
  fontupper=\popsemi\fontsize{10.4}{14.2}\selectfont\color{white},
  before upper={\poppill{Le savais-tu\,?}{white}{poppurple}\ifx\relax#1\relax\else\hspace{6pt}{\popxbold\color{white}#1}\fi\par\vspace{7pt}}}
% Exercice résolu : énoncé en haut, \tcblower puis la solution sur fond crème
\newcounter{popexo}\newcounter{popexr}
\newtcolorbox{exoresolu}[1][]{popbase,colframe=poporange,colbacklower=popcream,
  segmentation style={popink,line width=1.2pt,dashed},
  before upper={\stepcounter{popexr}\popboxhead{Exercice résolu \thepopexr}{poporange}{#1}},
  before lower={\poppill{Solution}{popink}{popcream}\par\vspace{6pt}}}
% Exercice (non résolu en place) — la correction est dans la partie Corrigés
\newtcolorbox{exercice}[1][]{popbase,colframe=popink,
  before upper={\stepcounter{popexo}\popboxhead{Exercice \thepopexo}{popink}{#1}}}
% Corrigé
\newtcolorbox{corrige}[1][]{popbase,colframe=popgreen,colback=popgreenL,
  before upper={\popboxhead{Corrigé}{popgreen}{#1}}}
% Objectifs / prérequis / bilan
\newtcolorbox{objectifs}{popbase,colframe=popink,colback=poporangeL,
  before upper={\popboxhead{Objectifs de la leçon}{popink}{}}}
\newtcolorbox{prerequis}{popbase,colframe=popink,colback=popblueL,
  before upper={\popboxhead{Prérequis}{popblue}{}}}
\newtcolorbox{bilan}{popbase,colframe=popink,colback=white,boxrule=2.8pt,
  before upper={\popboxhead{Fiche bilan}{poporange}{L'essentiel à connaître pour l'examen}}}
% Figure encadrée avec légende
\newtcolorbox{popfigure}[1][]{enhanced,breakable=false,colback=white,colframe=popline,boxrule=1.4pt,arc=8pt,
  left=10pt,right=10pt,top=10pt,bottom=8pt,before skip=10pt,after skip=12pt,halign=center,
  lower separated=false,halign lower=center,
  fontlower=\popsemi\fontsize{9}{11.5}\selectfont\color{popmuted},
  #1}
\newcommand{\legende}[1]{\par\vspace{6pt}{\popsemi\fontsize{9}{11.5}\selectfont\color{popmuted}#1}}

% ============================================================
% Signature OnBuch+
% ============================================================
\newcommand{\poplogoinline}[1]{{\poplogo\fontsize{#1}{#1}\selectfont\color{popink}OnBuch\color{poporange}+}}
\newcommand{\signatureonbuch}{%
  \begin{tcolorbox}[enhanced,colback=popink,colframe=popink,boxrule=2.2pt,arc=10pt,
      drop shadow=poporange,left=14pt,right=14pt,top=10pt,bottom=10pt,before skip=0pt,after skip=0pt]
    \tikz[baseline=-0.7ex]\node[circle,fill=popgreen,draw=popcream,line width=1.3pt,minimum size=22pt,inner sep=0]{\popcheck{white}};%
    \hspace{9pt}\begin{minipage}[c]{0.62\linewidth}
      {\popxbold\fontsize{12}{13}\selectfont\color{popcream}Signé\ \textbullet\ {\poplogo OnBuch\color{poporange}+}}\par\vspace{2pt}
      {\raggedright\popsemi\fontsize{8}{10.5}\selectfont\color{popline}\MakeUppercase{Équipe pédagogique OnBuch+}\\\MakeUppercase{Conforme au programme officiel MINESEC}\par}
    \end{minipage}\hfill
    {\poplogo\fontsize{22}{22}\selectfont\color{popcream}OnBuch\color{poporange}+}
  \end{tcolorbox}%
}

% ============================================================
% Page de garde
% #1 titre · #2 matière · #3 niveau · #4 module · #5 n° leçon · #6 durée ·
% #7 description / objectifs (texte ou itemize)
% ============================================================
\newcommand{\pagegarde}[7]{%
  \thispagestyle{empty}
  \newgeometry{a4paper,margin=1.7cm,top=1.6cm,bottom=1.4cm}%
  \begin{tikzpicture}[remember picture,overlay]
    \fill[poporange!18] ([xshift=-3.2cm,yshift=-3.6cm]current page.north east) circle (4.6cm);
    \fill[poppurple!14] ([xshift=2.4cm,yshift=7.6cm]current page.south west) circle (3.8cm);
  \end{tikzpicture}%
  {\poplogo\fontsize{30}{30}\selectfont\color{popink}OnBuch\color{poporange}+}\hfill
  \raisebox{0.6ex}{\poppill{Cours signé \textbullet\ Premium}{popink}{popcream}}\par
  \vspace{1pt}\popsquiggle{poppurple}\par
  \vspace{8pt}
  \begin{tcolorbox}[enhanced,colback=poporange,colframe=popink,boxrule=2.8pt,arc=13pt,
      drop shadow=popink,left=18pt,right=18pt,top=12pt,bottom=13pt,after skip=10pt]
    \poppill{#2 \textbullet\ #3}{popink}{popcream}\hspace{5pt}\poppill{#4}{white}{popink}\par\vspace{8pt}
    {\popdisplay\fontsize{14}{14}\selectfont\color{popink}LEÇON #5}\par\vspace{4pt}
    {\raggedright\hyphenpenalty=10000\exhyphenpenalty=10000\popdisplay\fontsize{25}{27}\selectfont\color{popcream}\MakeUppercase{#1}\par}
    \vspace{8pt}
    {\popxbold\fontsize{9.5}{9.5}\selectfont\color{popink}\MakeUppercase{Durée conseillée : #6}}
  \end{tcolorbox}
  \begin{tcolorbox}[enhanced,colback=white,colframe=popink,boxrule=2.2pt,arc=10pt,
      drop shadow=popink,left=16pt,right=16pt,top=10pt,bottom=10pt,after skip=8pt]
    \poppill{Dans cette leçon}{poppurple}{white}\par\vspace{5pt}
    \popsemi\fontsize{9.3}{12.6}\selectfont\color{popink}#7
  \end{tcolorbox}
  \noindent\begin{minipage}[t]{0.487\linewidth}
    \begin{tcolorbox}[enhanced,colback=popgreen,colframe=popink,boxrule=2.2pt,arc=10pt,
        drop shadow=popink,left=11pt,right=11pt,top=10pt,bottom=10pt,height=3.1cm,valign=center]
      \tcbox[colback=white,colframe=popink,boxrule=1.3pt,arc=5pt,boxsep=0pt,nobeforeafter,
        left=3pt,right=3pt,top=3pt,bottom=3pt,valign=center]{\qrcode[height=1.9cm]{https://play.google.com/store/apps/details?id=com.luvvix.onbuch&hl=fr}}%
      \hspace{8pt}\begin{minipage}[c]{0.5\linewidth}
        {\popdisplay\fontsize{11}{12.5}\selectfont\color{white}\MakeUppercase{Télécharge OnBuch}}\par\vspace{4pt}
        {\raggedright\popsemi\fontsize{8.4}{11}\selectfont\color{white}Gratuit \textbullet\ Google Play\\Tuteur IA, annales, concours, résultats\par}
      \end{minipage}
    \end{tcolorbox}
  \end{minipage}\hfill
  \begin{minipage}[t]{0.487\linewidth}
    \begin{tcolorbox}[enhanced,colback=poppurple,colframe=popink,boxrule=2.2pt,arc=10pt,
        drop shadow=popink,left=11pt,right=11pt,top=10pt,bottom=10pt,height=3.1cm,valign=center]
      \tikz[baseline=-0.7ex]\node[circle,fill=white,draw=popink,line width=1.3pt,minimum size=1.6cm,inner sep=0]{\popdisplay\fontsize{24}{24}\selectfont\color{poppurple}+};%
      \hspace{8pt}\begin{minipage}[c]{0.6\linewidth}
        {\popdisplay\fontsize{11}{12.5}\selectfont\color{white}\MakeUppercase{Passe à OnBuch+}}\par\vspace{4pt}
        {\raggedright\popsemi\fontsize{8.4}{11}\selectfont\color{white}Tous les cours signés, Tuteur IA illimité, fiches PDF — onglet OnBuch+ de l'app\par}
      \end{minipage}
    \end{tcolorbox}
  \end{minipage}
  \vspace{8pt}
  \popcontenu
  \vfill
  \signatureonbuch
  \restoregeometry
  \newpage
}

% Rangée « Ce cours contient » (page de garde)
\newcommand{\poptile}[3]{% #1 couleur #2 gros texte #3 légende
  \begin{tcolorbox}[enhanced,colback=#1,colframe=popink,boxrule=2pt,arc=9pt,shadow={2.5pt}{-2.5pt}{0pt}{popink},
      left=6pt,right=6pt,top=9pt,bottom=9pt,halign=center,valign=center,height=1.75cm,width=\linewidth,nobeforeafter]
    {\popdisplay\fontsize{12.5}{13}\selectfont\color{white}\MakeUppercase{#2}}\par\vspace{3pt}
    {\centering\popsemi\fontsize{7.8}{9.5}\selectfont\color{white}#3\par}
  \end{tcolorbox}}
\newcommand{\popcontenu}{%
  \par\noindent{\popxboldls\fontsize{7.5}{7.5}\selectfont\color{popmuted}CE COURS SIGNÉ CONTIENT}\par\vspace{7pt}
  \noindent\begin{minipage}[t]{0.235\linewidth}\poptile{popblue}{Cours}{complet, illustré, pas à pas}\end{minipage}\hfill
  \begin{minipage}[t]{0.235\linewidth}\poptile{poporange}{Méthodes}{et exercices résolus}\end{minipage}\hfill
  \begin{minipage}[t]{0.235\linewidth}\poptile{poppink}{Exercices}{type examen + corrigés}\end{minipage}\hfill
  \begin{minipage}[t]{0.235\linewidth}\poptile{popgreen}{Fiche bilan}{l'essentiel à retenir}\end{minipage}\par}

% Sommaire
\newtcolorbox{sommaire}{enhanced,colback=white,colframe=popink,boxrule=2.2pt,arc=10pt,
  drop shadow=popink,left=18pt,right=18pt,top=13pt,bottom=13pt,before skip=6pt,after skip=20pt,
  before upper={\poppill{Sommaire}{poppurple}{white}\par\vspace{9pt}\popsemi\fontsize{10.6}{14.5}\selectfont\color{popink}}}

% Signature de fin de cours — poussée tout en bas de la dernière page
\newcommand{\findecours}{%
  \par\vfill\nopagebreak
  \begin{center}\popsquiggle{poporange}\end{center}\vspace{4pt}
  \signatureonbuch
}
\AtBeginDocument{\def\llbracket{\mathopen{[\mkern-3.5mu[}}\def\rrbracket{\mathclose{]\mkern-3.5mu]}}} % intervalles d'entiers (glyphe ⟦ absent de la police)
% Glyphes absents de Fira Math : redéfinis à partir de glyphes disponibles
\AtBeginDocument{%
  \def\setminus{\mathbin{\backslash}}\let\smallsetminus\setminus
  \def\vdots{\mathord{\vbox{\baselineskip3.2pt\lineskiplimit0pt\kern4pt\hbox{.}\hbox{.}\hbox{.}}}}%
  \def\ddots{\mathinner{\mkern1mu\raise6.5pt\hbox{.}\mkern2mu\raise3.6pt\hbox{.}\mkern2mu\raise0.7pt\hbox{.}\mkern1mu}}%
  \def\triangle{\mathord{\scalebox{0.9}{$\symup{\Delta}$}}}%
  \def\nabla{\mathord{\raisebox{\depth}{\scalebox{1}[-1]{$\symup{\Delta}$}}}}%
  \def\checkmark{\popcheck{popdark}}%
  \def\star{\mathbin{\ast}}\def\bigstar{\mathord{\ast}}%
  \def\diamond{\mathbin{\scalebox{0.6}{$\Diamond$}}}%
  \def\bigcup{\mathop{\mathchoice{\raisebox{-0.3em}{\scalebox{1.7}{$\cup$}}}{\scalebox{1.25}{$\cup$}}{\cup}{\cup}}\displaylimits}%
  \def\bigcap{\mathop{\mathchoice{\raisebox{-0.3em}{\scalebox{1.7}{$\cap$}}}{\scalebox{1.25}{$\cap$}}{\cap}{\cap}}\displaylimits}%
  \def\lesseqqgtr{\mathrel{\lessgtr}}\def\bigoplus{\mathop{\oplus}\displaylimits}%
}
\providecommand{\ssi}{si et seulement si} % abréviation fréquente
\AtBeginDocument{\providecommand{\wideparen}{\overparen}} % arc de cercle

\begin{document}

\pagegarde{Les activités artisanales}{Géographie}{1re Tech}{Géographie – classe de Première technique}{N05}{3 séances (fiche de progression harmonisée)}{Cette leçon étudie l'artisanat camerounais : sa définition, sa diversité, son importance économique et sociale, ainsi que les problèmes qui freinent son développement. Elle conduit l'élève à proposer et justifier des solutions concrètes à partir de documents et de situations de la vie courante.
\vspace{4pt}
\begin{multicols}{2}\small
\begin{itemize}
\item À la fin de cette leçon, je sais définir l'artisanat et le distinguer de l'industrie
\item À la fin de cette leçon, je sais citer et classer les principales branches de l'artisanat camerounais
\item À la fin de cette leçon, je sais localiser sur une carte les grands pôles artisanaux du Cameroun
\item À la fin de cette leçon, je sais expliquer l'importance économique et sociale de l'artisanat à partir de données chiffrées
\item À la fin de cette leçon, je sais identifier et hiérarchiser les problèmes de l'artisanat camerounais
\item À la fin de cette leçon, je sais proposer et justifier des solutions aux problèmes de l'artisanat
\end{itemize}\end{multicols}}

\begin{sommaire}
\begin{multicols}{2}
\begin{enumerate}
\item Définition et caractéristiques de l'artisanat
\item Les branches et la répartition spatiale de l'artisanat camerounais
\item L'importance de l'artisanat au Cameroun
\item Les problèmes de l'artisanat camerounais
\item TD2 — Les solutions aux problèmes de l'artisanat
\item TP — Cartographie et représentation graphique de l'artisanat
\item Activité d'intégration
\item Méthodes et exercices résolus
\item Exercices
\item Corrigés des exercices
\item Fiche bilan
\end{enumerate}
\end{multicols}
\end{sommaire}

\begin{objectifs}
\begin{itemize}
\item À la fin de cette leçon, je sais définir l'artisanat et le distinguer de l'industrie
\item À la fin de cette leçon, je sais citer et classer les principales branches de l'artisanat camerounais
\item À la fin de cette leçon, je sais localiser sur une carte les grands pôles artisanaux du Cameroun
\item À la fin de cette leçon, je sais expliquer l'importance économique et sociale de l'artisanat à partir de données chiffrées
\item À la fin de cette leçon, je sais identifier et hiérarchiser les problèmes de l'artisanat camerounais
\item À la fin de cette leçon, je sais proposer et justifier des solutions aux problèmes de l'artisanat
\item À la fin de cette leçon, je sais construire un croquis ou un diagramme à partir de données sur l'artisanat
\item À la fin de cette leçon, je sais rédiger et justifier une conclusion argumentée sur une situation artisanale concrète
\end{itemize}
\end{objectifs}

\begin{prerequis}
\begin{itemize}
\item Les secteurs d'activités économiques (primaire, secondaire, tertiaire) vus en classe de 4e/3e
\item La notion de secteur informel abordée en début d'année
\item La lecture de cartes et de données statistiques simples
\item Les grandes régions et villes du Cameroun (repères de localisation)
\end{itemize}
\end{prerequis}

\newpage

\coursec{Définition et caractéristiques de l'artisanat}

\textbf{Situation d'accroche.} À Foumban, dans la région de l'Ouest, un sculpteur sur bois formé par son père peine à écouler ses œuvres faute de matériel moderne, de crédit bancaire et d'accès aux marchés extérieurs, alors que les touristes affluent dans la cité des arts. À Douala-Bonabéri, une couturière qui emploie trois apprentis ne parvient pas à formaliser son atelier. Pourtant, l'artisanat camerounais fait vivre des millions de personnes et constitue le premier employeur du secteur informel. Pourquoi un secteur aussi riche reste-t-il fragile ? Quelles solutions l'État et les artisans eux-mêmes peuvent-ils apporter ? C'est à ces questions que cette leçon répond.

\textbf{Problématique.} Qu'est-ce que l'artisanat, quelles sont ses caractéristiques propres, et en quoi se distingue-t-il de l'industrie ?

\courssub{Définition de l'artisanat}

\textbf{Intuition.} Observez le quartier où vous vivez : le menuisier qui fabrique les portes et les lits, la couturière qui confectionne les pagnes, le forgeron qui répare les houes et les machettes, le cordonnier qui ressemelle les chaussures. Tous travaillent avec leurs mains et des outils simples, dans de petits ateliers, souvent à ciel ouvert ou sous un hangar. Tous produisent ou réparent des biens utiles à la vie quotidienne. Ces activités ont un nom : l'\cle{artisanat}.

\begin{definition}[L'artisanat]
L'\cle{artisanat} est une activité économique de \textbf{transformation} (production de biens) ou de \textbf{service} (réparation, entretien) réalisée à \textbf{petite échelle}, avec des \textbf{outils simples} et un \textbf{savoir-faire souvent traditionnel}, transmis le plus souvent par apprentissage. L'artisan, qui est à la fois le patron et l'ouvrier de son atelier, y joue un rôle central.
\end{definition}

Décomposons cette définition pour bien la comprendre :
\begin{itemize}
\item \textbf{Activité de transformation ou de service} : l'artisan transforme une matière première (bois, fer, cuir, tissu, argile) en produit fini (meuble, marmite, chaussure, pagne, poterie), ou bien il rend un service (réparation de cycles, coiffure, soudure).
\item \textbf{Petite échelle} : l'atelier est de taille réduite, la production est limitée, souvent réalisée à la commande.
\item \textbf{Outils simples} : marteau, scie, ciseau à bois, machine à coudre, enclume, tour de potier. L'énergie utilisée est surtout la force humaine.
\item \textbf{Savoir-faire traditionnel} : les techniques sont apprises par la pratique, auprès d'un maître artisan, et non dans une école d'ingénieurs.
\end{itemize}

\begin{exemplebox}[Un atelier de menuiserie de quartier]
Un atelier de menuiserie de quartier à Yaoundé ou à Bafoussam emploie en général \textbf{moins de 10 personnes} : le patron-menuisier, un ou deux ouvriers qualifiés et quelques apprentis. Il produit quelques meubles par semaine, sur commande. À l'opposé, une usine de meubles industrielle emploie \textbf{plusieurs centaines de salariés} et produit des milliers d'articles standardisés par mois à l'aide de machines automatisées.
\end{exemplebox}

\courssub{Distinction entre artisanat et industrie}

Il est essentiel, pour l'examen, de ne pas confondre artisanat et industrie. Les deux transforment des matières premières, mais ils diffèrent par la \textbf{taille de l'unité}, l'\textbf{outillage}, le \textbf{volume de production} et la \textbf{main-d'œuvre}.

\begin{center}
\begin{tabularx}{\linewidth}{|l|X|X|}\hline
\rowcolor{popblueL} \textbf{Critère} & \textbf{Artisanat} & \textbf{Industrie} \\ \hline
Taille de l'unité & Petit atelier (moins de 10 personnes en général) & Usine de grande taille (dizaines à centaines de salariés) \\ \hline
Capital & Faible (quelques centaines de milliers de FCFA) & Élevé (millions à milliards de FCFA) \\ \hline
Outillage & Outils simples, force humaine, peu de machines & Machines complexes, automatisées, énergie abondante \\ \hline
Volume de production & Faible, souvent à la commande, produits variés & Élevé, production de masse standardisée \\ \hline
Main-d'œuvre & Familiale et apprentis, polyvalente & Salariée, nombreuse et spécialisée \\ \hline
Organisation & Le patron est aussi ouvrier ; gestion simple & Hiérarchie complexe, gestion comptable \\ \hline
\end{tabularx}
\end{center}

\begin{popfigure}
\begin{tikzpicture}
% Colonne artisanat
\draw[fill=popgreenL,draw=popgreen,thick,rounded corners=4pt] (0,0) rectangle (5.6,6.4);
\node[font=\bfseries\large,text=popdark] at (2.8,5.9) {ARTISANAT};
% petit atelier
\draw[fill=popcream,draw=popdark] (0.6,3.6) rectangle (2.6,4.6);
\draw[draw=popdark] (0.5,4.6) -- (1.6,5.3) -- (2.7,4.6);
\node[font=\scriptsize] at (1.6,4.1) {atelier};
\draw[fill=popdark] (1.2,3.35) circle (0.09);
\draw[fill=popdark] (1.7,3.35) circle (0.09);
\draw[fill=popdark] (2.2,3.35) circle (0.09);
\node[font=\scriptsize] at (1.7,3.0) {moins de 10 pers.};
% outils
\draw[draw=popdark,thick] (3.4,4.3) -- (4.4,3.7);
\draw[fill=popdark] (3.3,4.35) rectangle (3.55,4.6);
\node[font=\scriptsize] at (4.3,4.35) {outils simples};
\node[font=\scriptsize,text=popdark,align=center] at (2.8,2.2) {Capital faible\\Production faible\\Savoir-faire manuel};
% Colonne industrie
\draw[fill=poporangeL,draw=poporange,thick,rounded corners=4pt] (6.4,0) rectangle (12,6.4);
\node[font=\bfseries\large,text=popdark] at (9.2,5.9) {INDUSTRIE};
% usine
\draw[fill=popcream,draw=popdark] (7.0,3.4) rectangle (10.6,4.4);
\draw[draw=popdark] (7.0,4.4) -- (7.6,5.0) -- (8.2,4.4) -- (8.8,5.0) -- (9.4,4.4) -- (10.0,5.0) -- (10.6,4.4);
\draw[fill=popcream,draw=popdark] (10.2,4.4) rectangle (10.5,5.5);
\node[font=\scriptsize] at (8.8,3.9) {usine};
\foreach \x in {7.2,7.7,8.2,8.7,9.2,9.7,10.2}{\draw[fill=popdark] (\x,3.1) circle (0.07);}
\node[font=\scriptsize] at (8.7,2.75) {centaines de salariés};
% engrenage simplifié
\draw[draw=popdark,thick] (11.3,4.6) circle (0.35);
\foreach \a in {0,45,...,315}{\draw[draw=popdark,thick] (11.3,4.6) ++(\a:0.35) -- ++(\a:0.18);}
\node[font=\scriptsize] at (11.3,3.95) {machines};
\node[font=\scriptsize,text=popdark,align=center] at (9.2,2.2) {Capital élevé\\Production de masse\\Techniques avancées};
% flèche de comparaison
\draw[<->,thick,poppurple] (5.7,3.2) -- (6.3,3.2);
\end{tikzpicture}
\legende{Schéma comparatif artisanat -- industrie : taille de l'unité, main-d'œuvre, outillage, capital et volume de production.}
\end{popfigure}

\begin{popfigure}
\begin{tikzpicture}
% Photo 1 : forge de Maroua (croquis)
\draw[fill=popcream,draw=popdark,thick] (0,0) rectangle (5.8,4.2);
\draw[fill=popgoldL,draw=popdark] (0,2.6) rectangle (5.8,4.2);
\node[font=\scriptsize\itshape] at (2.9,3.9) {Ciel du Sahel -- Maroua (Extrême-Nord)};
\draw[fill=popmuted!40,draw=popdark] (0.4,0.4) rectangle (3.2,2.6);
\draw[draw=popdark] (0.3,2.6) -- (1.8,3.3) -- (3.3,2.6);
% forgeron et enclume
\draw[fill=popdark] (1.3,1.6) circle (0.14);
\draw[draw=popdark,thick] (1.3,1.45) -- (1.3,0.8);
\draw[draw=popdark,thick] (1.3,1.2) -- (1.9,1.5);
\draw[draw=popdark,thick] (1.3,0.8) -- (1.0,0.45);
\draw[draw=popdark,thick] (1.3,0.8) -- (1.6,0.45);
\draw[fill=popmuted,draw=popdark] (2.1,0.9) rectangle (2.7,1.15);
\draw[fill=poporange,draw=poporange] (4.2,0.5) circle (0.28);
\draw[fill=poppink,draw=poppink] (4.2,0.5) circle (0.13);
\node[font=\scriptsize] at (4.2,0.15) {foyer (charbon)};
\node[font=\scriptsize\bfseries] at (2.9,-0.45) {1. Atelier de forge traditionnelle à Maroua};
% Photo 2 : ALUCAM Edéa
\draw[fill=popcream,draw=popdark,thick] (6.4,0) rectangle (12.2,4.2);
\draw[fill=popblueL,draw=popdark] (6.4,2.6) rectangle (12.2,4.2);
\node[font=\scriptsize\itshape] at (9.3,3.9) {Bord de la Sanaga -- Édéa (Littoral)};
\draw[fill=popmuted!40,draw=popdark] (6.8,0.6) rectangle (10.4,2.4);
\draw[draw=popdark] (6.8,2.4) -- (7.4,2.9) -- (8.0,2.4) -- (8.6,2.9) -- (9.2,2.4) -- (9.8,2.9) -- (10.4,2.4);
\draw[fill=popmuted!60,draw=popdark] (10.7,0.6) rectangle (11.1,3.3);
\draw[fill=popmuted!60,draw=popdark] (11.4,0.6) rectangle (11.8,3.0);
\node[font=\scriptsize] at (8.6,1.5) {halles de production};
\node[font=\scriptsize] at (11.25,0.25) {cheminées};
\node[font=\scriptsize\bfseries] at (9.3,-0.45) {2. Usine ALUCAM (aluminium) à Édéa};
\end{tikzpicture}
\legende{Deux mondes de la production au Cameroun : (1) la forge traditionnelle de Maroua, petite unité à outils simples ; (2) l'usine ALUCAM d'Édéa, grande industrie lourde alimentée par l'énergie hydroélectrique de la Sanaga (croquis d'après photographies).}
\end{popfigure}

\begin{attention}[Artisanat et secteur informel : ne pas confondre]
Erreur fréquente à l'examen : affirmer que << tout l'artisanat est informel >>. C'est faux. Le \cle{secteur informel} regroupe les activités non déclarées, non enregistrées et échappant à la fiscalité. Or il existe des artisans \textbf{formalisés} : immatriculés au registre de commerce, payant leurs impôts, organisés en coopératives ou en groupements d'initiative commune (GIC). La majorité des artisans camerounais travaille certes dans l'informel, mais l'artisanat et l'informel ne sont pas synonymes.
\end{attention}

\courssub{Les caractéristiques de l'artisanat}

L'observation des ateliers camerounais permet de dégager quatre caractéristiques essentielles.

\textbf{1. De petites unités de production.} L'atelier artisanal occupe un espace réduit : une échoppe, un hangar, une cour, parfois le bord de la route. Le nombre de travailleurs est faible, généralement inférieur à dix.

\textbf{2. Une main-d'œuvre familiale et des apprentis.} L'artisan travaille avec ses enfants, ses frères ou des apprentis venus apprendre le métier. Les salaires sont faibles ou remplacés par des avantages en nature (nourriture, logement, formation).

\textbf{3. Un faible capital.} L'équipement coûte peu : quelques outils, une machine à coudre, une forge. L'artisan n'a généralement pas accès au crédit bancaire, ce qui l'empêche de moderniser son atelier.

\textbf{4. Une forte présence dans le secteur informel.} Beaucoup d'ateliers ne sont pas déclarés : ils échappent aux statistiques officielles, à l'impôt, mais aussi à la protection sociale et aux aides de l'État.

\begin{aretenir}[Les quatre caractéristiques de l'artisanat]
Petites unités ; main-d'œuvre familiale et apprentis ; faible capital ; forte présence dans le secteur informel. Retenez la formule : << \textbf{petit, familial, pauvre, souvent informel} >>.
\end{aretenir}

\courssub{L'apprentissage et la transmission du savoir-faire}

Comment devient-on artisan au Cameroun ? Le plus souvent, non pas à l'école, mais \textbf{sur le tas}. Le jeune entre comme apprenti dans l'atelier d'un maître artisan (parent, voisin ou patron du quartier). Pendant plusieurs années, il observe, exécute les tâches simples, puis les gestes complexes, jusqu'à maîtriser le métier. C'est le système de l'\cle{apprentissage sur le tas}, parfois surnommé, par analogie avec les revendeuses qui apprennent le commerce en pratiquant, le système << \cle{bayam-sellam} >> de la formation : on apprend en faisant, sans diplôme ni manuel.

Ce mode de formation présente des avantages : il est gratuit ou peu coûteux, concret, adapté au rythme de l'apprenti, et il garantit la \textbf{transmission du savoir-faire de génération en génération}. À Foumban, les sculpteurs sur bois apprennent auprès de leur père, qui tenait lui-même son art de son propre père : certaines familles perpétuent ainsi des techniques séculaires (bronze à la cire perdue, broderie, vannerie). Mais cet apprentissage a aussi des limites : absence de formation théorique, de gestion et de comptabilité, techniques parfois dépassées, aucun diplôme reconnu pour valoriser la qualification.

\courssub{Artisanat d'art et artisanat de production}

On distingue classiquement deux grandes catégories d'artisanat.

\begin{definition}[Artisanat d'art et artisanat de production]
L'\cle{artisanat d'art} produit des objets à visée \textbf{culturelle et touristique} : sculptures, masques, statuettes, bijoux, poteries décoratives, objets en bronze de Foumban, vannerie fine. Sa valeur est esthétique et symbolique. \par
L'\cle{artisanat de production} fabrique des \textbf{biens de consommation courante} : meubles, vêtements, chaussures, marmites en aluminium, houes, briques, pains. Il répond aux besoins quotidiens de la population.
\end{definition}

\begin{exemplebox}[Deux exemples camerounais]
\textbf{Artisanat d'art} : le quartier des artisans de Foumban (Ouest) vend masques bamoun, trônes perlés et statuettes aux touristes et aux collectionneurs. \textbf{Artisanat de production} : les fondeurs d'aluminium de Douala recyclent les casseroles usagées pour fabriquer les marmites vendues sur tous les marchés du pays.
\end{exemplebox}

\begin{exoresolu}[Identifier et justifier une activité artisanale]
Énoncé. Madame Aïssatou fabrique et vend des beignets-haricots (<< beignets-bouillons >>) devant sa concession à Ngaoundéré, aidée de ses deux filles, avec une marmite, un foyer et une bassine. (a) Cette activité est-elle artisanale ? Justifiez avec deux caractéristiques. (b) S'agit-il d'artisanat d'art ou de production ? (c) Pourquoi dit-on qu'elle relève du secteur informel ?
\tcblower
Solution. (a) Oui, c'est une activité artisanale : c'est une \textbf{petite unité} (une vendeuse et deux aides familiales) utilisant un \textbf{outillage simple} (marmite, foyer) et un savoir-faire pratique, avec un \textbf{faible capital}. (b) C'est un artisanat \textbf{de production}, car elle fabrique un bien de consommation courante (nourriture), et non un objet culturel ou touristique. (c) L'activité est \textbf{informelle} car elle n'est ni déclarée ni immatriculée : elle échappe à la fiscalité et aux statistiques officielles, ce qui est le cas de la majorité des artisans camerounais.
\end{exoresolu}

\begin{savaistu}[L'artisanat, premier employeur du Cameroun]
Le Cameroun compte des \textbf{centaines de milliers d'artisans} -- forgerons, menuisiers, couturiers, coiffeurs, cordonniers, mécaniciens, potiers, vanniers -- ce qui fait de l'artisanat l'un des \textbf{premiers pourvoyeurs d'emplois} du pays et le premier employeur du secteur informel. Dans de nombreuses villes, il fournit plus d'emplois que l'industrie moderne. La ville de Foumban est d'ailleurs surnommée la << cité des arts >> pour la richesse de son artisanat sculptural, hérité du royaume bamoun.
\end{savaistu}

\begin{exercice}[À vous de jouer]
1. Donnez la définition de l'artisanat. 2. Citez trois différences entre l'artisanat et l'industrie. 3. Expliquez ce qu'est l'apprentissage << sur le tas >> et donnez un avantage et une limite de ce système. 4. Classez les activités suivantes en artisanat d'art ou de production : fabrication de marmites, sculpture de masques, couture de vêtements, fabrication de bijoux en perles pour touristes.
\end{exercice}

\begin{corrige}[Exercice]
1. L'artisanat est une activité économique de transformation ou de service réalisée à petite échelle, avec des outils simples et un savoir-faire souvent traditionnel. 2. Taille (petit atelier contre grande usine), outillage (outils simples contre machines automatisées), volume de production (faible, à la commande, contre production de masse) ; on peut ajouter le capital et la main-d'œuvre. 3. L'apprentissage sur le tas consiste à apprendre le métier par la pratique auprès d'un maître artisan, sans école ni diplôme. Avantage : gratuité, formation concrète, transmission du savoir-faire. Limite : absence de formation théorique et de diplôme reconnu. 4. Artisanat de production : marmites, couture de vêtements. Artisanat d'art : masques sculptés, bijoux en perles pour touristes.
\end{corrige}

\textbf{Transition.} Maintenant que l'artisanat est défini et caractérisé, la suite de la leçon étudiera ses branches et sa répartition spatiale au Cameroun, avant d'en mesurer l'importance économique et d'en analyser les problèmes.

\coursec{Les branches et la répartition spatiale de l'artisanat camerounais}

Dans la section précédente, nous avons défini l'artisanat et dégagé ses caractéristiques : petites unités de production, travail essentiellement manuel ou avec des outils simples, faible capital, savoir-faire transmis par apprentissage. Mais dire « artisanat » ne suffit pas : réparer une moto à Douala, sculpter un masque à Foumban ou tresser un panier à Bamenda sont des métiers très différents. L'artisanat camerounais se divise en plusieurs \cle{branches} d'activité, et ces branches ne sont pas réparties au hasard sur le territoire : elles forment des \cle{pôles} régionaux, liés aux ressources locales et aux traditions culturelles. C'est cette géographie de l'artisanat que nous allons maintenant étudier.

\courssub{Les grandes branches de l'artisanat camerounais}

\textbf{Intuition.}\par
Observez votre quartier ou votre village : le tailleur, le maçon, le mécanicien, la coiffeuse, le sculpteur n'appartiennent pas au même métier. Classer l'artisanat en branches, c'est regrouper les métiers qui fabriquent ou réparent des produits de même nature.

On distingue six grandes branches de l'artisanat au Cameroun :

\begin{enumerate}
\item \textbf{L'artisanat d'art} : sculpture sur bois (masques, statuettes), fonderie de bronze, vannerie (paniers, nattes), poterie, tissage, bijouterie (perles, colliers). Ce sont des productions à forte valeur culturelle.
\item \textbf{Le bâtiment et les travaux publics (BTP)} : maçonnerie, menuiserie (meubles, portes, fenêtres), plomberie, électricité du bâtiment, peinture. C'est la branche la plus répandue dans les villes en pleine croissance.
\item \textbf{La mécanique et la réparation} : garages automobiles, réparation de motos (les « benskin »), soudure, réparation d'appareils électroménagers et de téléphones.
\item \textbf{L'habillement} : couture, broderie, teinture, fabrication de chaussures (cordonnerie).
\item \textbf{L'agroalimentaire} : transformation artisanale des produits agricoles (huile de palme, gari, farine de maïs, jus de fruits, fumage du poisson), boulangerie, boucherie.
\item \textbf{Les services} : coiffure, réparation de cycles, photographie, cordonnerie de dépannage.
\end{enumerate}

\begin{definition}[L'artisanat d'art]
L'\cle{artisanat d'art} regroupe les productions manuelles à forte valeur culturelle destinées notamment au tourisme : masques, statuettes, tissus traditionnels, objets en bronze, poteries décoratives, bijoux. Il fait vivre la culture matérielle des peuples et constitue une vitrine du pays à l'étranger.
\end{definition}

\begin{exemplebox}[L'artisanat d'art au quotidien]
Un masque bamiléké sculpté dans le bois, un panier tressé en raphia du Nord-Ouest, une jarre en argile de l'Extrême-Nord ou un collier de perles bamoun sont des produits de l'artisanat d'art : ils exigent un savoir-faire précis, portent une identité culturelle et trouvent des acheteurs parmi les touristes et les collectionneurs.
\end{exemplebox}

\begin{attention}[Ne pas confondre les branches]
Une erreur fréquente consiste à classer la couture ou la coiffure dans l'artisanat d'art parce qu'elles sont « manuelles ». L'artisanat d'art se définit par la \textbf{valeur culturelle et artistique} du produit, pas seulement par le travail manuel. La couture relève de l'habillement, la coiffure des services. De même, la mécanique est une branche à part entière, distincte du bâtiment.
\end{attention}

\courssub{La répartition spatiale : les grands pôles artisanaux du Cameroun}

\textbf{Intuition.}\par
Pourquoi trouve-t-on des sculpteurs surtout à Foumban, des tanneurs de cuir surtout à Maroua et des garages surtout à Douala ? Parce que chaque métier a besoin de matières premières proches (bois, argile, cuir) et s'enracine dans une tradition locale transmise de génération en génération. L'artisanat est donc fortement \cle{localisé}.

On distingue quatre grands pôles artisanaux :

\begin{itemize}
\item \textbf{Foumban (Ouest)} : la « cité des arts ». Capitale du royaume bamoun, elle est le premier centre national de sculpture sur bois et de fonderie de bronze. La « Rue des Artisans », près du Palais royal des Bamoun, aligne les ateliers et les boutiques.
\item \textbf{Bamenda et le Nord-Ouest} : vannerie (paniers, nattes en raphia), tissage (pagne \emph{ndop}), sculpture et broderie. Les hautes terres fournissent le raphia et le bois.
\item \textbf{Maroua et l'Extrême-Nord} : travail du cuir (tannerie, sandales, sacs), poterie, tissage (coton), bijouterie. L'élevage y fournit les peaux, et l'argile des plaines alimente la poterie.
\item \textbf{Douala et Yaoundé} : les deux grandes métropoles concentrent l'artisanat urbain moderne : mécanique automobile, menuiserie, couture, réparation d'appareils, coiffure. La demande urbaine massive explique cette concentration.
\end{itemize}

\begin{exemplebox}[Foumban, capitale de l'artisanat d'art]
Foumban concentre \textbf{plusieurs centaines d'ateliers} de sculpture et de fonderie autour du Palais royal des Bamoun. Les artisans y produisent masques, statues, tabourets royaux et objets en bronze selon des techniques héritées des ateliers royaux. Le Musée du Palais et le Musée des Arts et Traditions bamoun attirent chaque année des visiteurs camerounais et étrangers, faisant de l'artisanat un pilier de l'économie locale.
\end{exemplebox}

\begin{popfigure}
\begin{center}
\begin{tikzpicture}[scale=0.92, every node/.style={transform shape}]
% Contour très simplifié du Cameroun
\draw[thick, popdark, fill=popcream] (1.2,0) -- (3.4,0) -- (4.6,0.8) -- (5.4,2.2) -- (5.0,3.4) -- (4.2,4.4) -- (3.6,5.6) -- (3.2,7.0) -- (2.6,8.2) -- (1.8,8.6) -- (1.2,7.6) -- (1.0,6.2) -- (0.6,5.0) -- (0.2,4.0) -- (0.0,3.0) -- (0.4,1.8) -- cycle;
% Nord
\draw[->, thick, popdark] (6.0,7.6) -- (6.0,8.6) node[above]{\textbf{N}};
% Pôles artisanaux
\fill[poporange] (2.0,3.6) circle (0.12) node[right=2pt, popdark]{\footnotesize \textbf{Foumban}};
\fill[poppurple] (1.2,4.6) circle (0.12) node[left=2pt, popdark]{\footnotesize \textbf{Bamenda}};
\fill[popgreen] (2.4,7.6) circle (0.12) node[right=2pt, popdark]{\footnotesize \textbf{Maroua}};
\fill[popblue] (0.7,1.2) circle (0.12) node[left=2pt, popdark]{\footnotesize \textbf{Douala}};
\fill[popblue] (2.2,1.6) circle (0.12) node[right=2pt, popdark]{\footnotesize \textbf{Yaoundé}};
% Légende des figurés
\node[anchor=west] at (4.6,6.6) {\footnotesize \textcolor{poporange}{$\bullet$} Artisanat d'art (sculpture, bronze)};
\node[anchor=west] at (4.6,6.1) {\footnotesize \textcolor{poppurple}{$\bullet$} Vannerie, tissage (ndop)};
\node[anchor=west] at (4.6,5.6) {\footnotesize \textcolor{popgreen}{$\bullet$} Cuir, poterie, tissage coton};
\node[anchor=west] at (4.6,5.1) {\footnotesize \textcolor{popblue}{$\bullet$} Mécanique, menuiserie, couture};
% Echelle indicative
\draw[thick, popdark] (4.7,4.4) -- (5.7,4.4);
\node[below, popdark] at (5.2,4.4) {\footnotesize $\approx$ 200 km};
\end{tikzpicture}
\end{center}
\legende{Carte schématique de la répartition des principaux pôles artisanaux du Cameroun. Chaque couleur de point correspond à la branche dominante du pôle. Contour simplifié, échelle indicative.}
\end{popfigure}

\courssub{Le lien entre spécialités artisanales, ressources locales et traditions culturelles}

\textbf{Observation.}\par
Confrontons la carte des pôles artisanaux à la carte des ressources naturelles et des peuples du Cameroun. On constate une correspondance frappante : là où il y a du bois, on sculpte ; là où il y a de l'argile, on fait de la poterie ; là où l'on élève des bœufs, on travaille le cuir.

\begin{propriete}[Facteurs de localisation de l'artisanat]
La répartition spatiale de l'artisanat s'explique par deux grands facteurs :
\begin{enumerate}
\item les \cle{ressources naturelles locales} : bois des forêts de l'Ouest et du Sud (sculpture, menuiserie), argile des plaines du Nord (poterie), peaux issues de l'élevage de l'Adamaoua et de l'Extrême-Nord (cuir), raphia des zones humides de l'Ouest et du Nord-Ouest (vannerie), coton du Nord (tissage) ;
\item les \cle{traditions culturelles} : les royaumes et chefferies (Bamoun, Bamiléké, peuples du Nord) ont entretenu des ateliers royaux et des corporations d'artisans, transmettant les techniques de parents à enfants depuis des siècles.
\end{enumerate}
\end{propriete}

\begin{center}
\begin{tabularx}{\linewidth}{|l|X|X|}
\hline
\rowcolor{popblueL} \textbf{Pôle} & \textbf{Ressource locale} & \textbf{Tradition / branche dominante} \\
\hline
Foumban (Ouest) & Bois, minerais (bronze) & Ateliers royaux bamoun : sculpture, fonderie \\
\hline
Bamenda (Nord-Ouest) & Raphia, bois & Vannerie, tissage du pagne \emph{ndop} \\
\hline
Maroua (Extrême-Nord) & Peaux (élevage), argile, coton & Tannerie, poterie, tissage \\
\hline
Douala -- Yaoundé & Demande urbaine, matériaux importés & Mécanique, menuiserie, couture \\
\hline
\end{tabularx}
\end{center}

\begin{savaistu}[Des tissus camerounais de renommée internationale]
Le pagne \emph{ndop} du Nord-Ouest et le \emph{toghu} (costume brodé de velours noir) sont des tissus artisanaux camerounais reconnus dans le monde entier. Portés lors des grandes cérémonies traditionnelles, ils ont été exhibés lors d'événements internationaux, notamment par des sportifs et des artistes camerounais, et inspirent aujourd'hui des créateurs de mode en Afrique et en Europe.
\end{savaistu}

\courssub{Lire et interpréter une carte de répartition des activités artisanales}

La lecture d'une carte de répartition est une compétence exigible au Probatoire (et ultérieurement au Baccalauréat technique). Elle obéit à une démarche précise.

\begin{methode}[Lire une carte de répartition artisanale]
\begin{enumerate}
\item \textbf{Identifier les figurés} : lire d'abord le titre puis la légende (couleurs, symboles) pour savoir quelles branches sont représentées.
\item \textbf{Repérer les concentrations} : localiser les zones où les figurés se regroupent (pôles) et celles qui en sont dépourvues.
\item \textbf{Décrire} : nommer les pôles du plus important au moins important, en citant les régions et les villes.
\item \textbf{Expliquer} : relier chaque concentration aux ressources locales (bois, argile, cuir, raphia) et à l'histoire culturelle des régions (royaumes, chefferies), ou à la demande urbaine pour les métropoles.
\item \textbf{Conclure} : dégager le déséquilibre général (artisanat d'art rural et traditionnel au Nord et à l'Ouest ; artisanat moderne concentré dans les grandes villes du Sud).
\end{enumerate}
\end{methode}

\begin{exoresolu}[Interpréter la carte des pôles artisanaux]
Énoncé : à partir de la carte schématique ci-dessus, (a) citez deux pôles de l'artisanat d'art et leur spécialité ; (b) expliquez la localisation de l'artisanat du cuir à Maroua ; (c) pourquoi Douala et Yaoundé concentrent-elles la mécanique et la couture ?
\tcblower
Solution détaillée :
\begin{enumerate}
\item[(a)] Foumban (Ouest) est spécialisé dans la sculpture sur bois et la fonderie de bronze ; Bamenda (Nord-Ouest) dans la vannerie en raphia et le tissage du pagne \emph{ndop}.
\item[(b)] Maroua se situe dans l'Extrême-Nord, une grande région d'élevage bovin : les peaux y sont abondantes et bon marché. De plus, la tradition du tannage et du travail du cuir y est ancienne, liée aux chefferies peules. Ressource et tradition se combinent donc pour expliquer ce pôle.
\item[(c)] Douala et Yaoundé sont les deux plus grandes villes du pays : elles concentrent une forte demande (véhicules à réparer, meubles, vêtements) et une main-d'œuvre nombreuse. Leur artisanat est donc dominé par les branches urbaines modernes (mécanique, menuiserie, couture), indépendantes des ressources naturelles locales.
\end{enumerate}
Conclusion : la carte montre un artisanat d'art enraciné dans les ressources et les traditions des régions de l'Ouest et du Nord, et un artisanat de services concentré dans les métropoles du Sud.
\end{exoresolu}

\begin{exercice}[Répartition par branche]
Le diagramme circulaire ci-dessous donne une répartition indicative des artisans camerounais par branche. 1) Quelle est la branche la plus importante ? 2) Calculez la part cumulée des branches « BTP » et « habillement ». 3) Expliquez pourquoi l'artisanat d'art, bien que minoritaire en effectifs, reste stratégique pour le Cameroun.
\end{exercice}

\begin{popfigure}
\begin{center}
\begin{tikzpicture}[scale=1.0]
% Diagramme circulaire : BTP 30%, habillement 20%, mécanique 15%, agroalimentaire 15%, services 12%, art 8%
\def\r{2.2}
\fill[popblue] (0,0) -- (90:\r) arc (90:90-108:\r) -- cycle;
\fill[poporange] (0,0) -- (90-108:\r) arc (90-108:90-108-72:\r) -- cycle;
\fill[popgreen] (0,0) -- (90-180:\r) arc (90-180:90-180-54:\r) -- cycle;
\fill[poppurple] (0,0) -- (90-234:\r) arc (90-234:90-234-54:\r) -- cycle;
\fill[popgold] (0,0) -- (90-288:\r) arc (90-288:90-288-43.2:\r) -- cycle;
\fill[poppink] (0,0) -- (90-331.2:\r) arc (90-331.2:90-360:\r) -- cycle;
% Etiquettes
\node at (90-54:1.5) {\footnotesize 30\,\%};
\node at (90-144:1.5) {\footnotesize 20\,\%};
\node at (90-207:1.5) {\footnotesize 15\,\%};
\node at (90-261:1.5) {\footnotesize 15\,\%};
\node at (90-309.6:1.55) {\footnotesize 12\,\%};
\node at (90-345.6:1.6) {\footnotesize 8\,\%};
% Légende
\node[anchor=west] at (2.8,1.5) {\footnotesize \textcolor{popblue}{$\bullet$} Bâtiment et travaux publics};
\node[anchor=west] at (2.8,1.0) {\footnotesize \textcolor{poporange}{$\bullet$} Habillement (couture, teinture)};
\node[anchor=west] at (2.8,0.5) {\footnotesize \textcolor{popgreen}{$\bullet$} Mécanique et réparation};
\node[anchor=west] at (2.8,0.0) {\footnotesize \textcolor{poppurple}{$\bullet$} Agroalimentaire};
\node[anchor=west] at (2.8,-0.5) {\footnotesize \textcolor{popgold}{$\bullet$} Services (coiffure, etc.)};
\node[anchor=west] at (2.8,-1.0) {\footnotesize \textcolor{poppink}{$\bullet$} Artisanat d'art};
\end{tikzpicture}
\end{center}
\legende{Répartition indicative des artisans camerounais par branche d'activité (en \% des effectifs, valeurs arrondies à visée pédagogique).}
\end{popfigure}

\begin{corrige}[Exercice : Répartition par branche]
\begin{enumerate}
\item La branche la plus importante est le \textbf{bâtiment et les travaux publics} (30\,\% des artisans), portée par la construction rapide des villes.
\item Part cumulée BTP + habillement : $30 + 20 = 50\,\%$. La moitié des artisans camerounais exerce donc dans ces deux branches.
\item L'artisanat d'art ne représente qu'environ 8\,\% des effectifs, mais il reste stratégique : il attire les touristes (devises), exporte des produits à forte valeur ajoutée, préserve l'identité culturelle nationale et fait vivre des pôles entiers comme Foumban.
\end{enumerate}
\end{corrige}

\begin{experience}[TP : Croquis de la répartition de l'artisanat camerounais]
\textbf{Objectif :} Réaliser un croquis légendé de la répartition des pôles artisanaux au Cameroun.
\textbf{Matériel :} Carte de base du Cameroun (contour, fleuves, villes principales), crayons de couleur, règle.
\textbf{Consignes :}
\begin{enumerate}
\item Tracer le contour du Cameroun et positionner les 4 pôles majeurs (Foumban, Bamenda, Maroua, Douala/Yaoundé).
\item Choisir un code couleur par branche dominante (ex: orange pour art, violet pour vannerie, vert pour cuir, bleu pour urbain).
\item Figurer les pôles par des surfaces ou points proportionnels à leur importance.
\item Tracer des flèches d'approvisionnement (ex: bois vers Foumban, peaux vers Maroua).
\item Compléter la légende, le titre, l'échelle et l'orientation.
\end{enumerate}
\textbf{Attendu :} Un croquis propre, lisible, expliquant la localisation par les ressources et la demande.
\end{experience}

\begin{aretenir}[L'essentiel de la section]
\begin{itemize}
\item L'artisanat camerounais comprend six grandes branches : artisanat d'art, BTP, mécanique et réparation, habillement, agroalimentaire, services.
\item Quatre pôles dominent : \textbf{Foumban} (sculpture, cité des arts), \textbf{Bamenda et le Nord-Ouest} (vannerie, tissage), \textbf{Maroua et l'Extrême-Nord} (cuir, poterie), \textbf{Douala et Yaoundé} (mécanique, menuiserie, couture).
\item La localisation s'explique par les \textbf{ressources locales} (bois, argile, cuir, raphia) et les \textbf{traditions culturelles} (ateliers royaux, chefferies), complétées par la \textbf{demande urbaine} dans les métropoles.
\item La lecture d'une carte de répartition suit cinq étapes : figurés, concentrations, description, explication, conclusion.
\end{itemize}
\end{aretenir}

Cette répartition spatiale, héritée de l'histoire et des ressources, n'est pas qu'une curiosité géographique : elle explique le poids économique et social de l'artisanat dans chaque région, poids que nous mesurerons dans la section suivante consacrée à l'importance de l'artisanat au Cameroun.

\coursec{L'importance de l'artisanat au Cameroun}

Nous venons de décrire les branches de l'artisanat camerounais (artisanat d'art, artisanat de production, artisanat de service) et leur répartition spatiale (villages d'artisans de Foumban et de Bafoussam, quartiers spécialisés de Douala et de Yaoundé, artisanat rural dispersé). Mais à quoi sert concrètement ce secteur ? Pourquoi l'État, les collectivités et les partenaires internationaux s'intéressent-ils à lui ? C'est toute la question de l'\cle{importance} de l'artisanat, que nous allons étudier sous quatre angles : \cle{social}, \cle{économique}, \cle{culturel et touristique}, \cle{fiscal et monétaire}.

\courssub{L'importance sociale : le premier employeur du Cameroun}

\begin{definition}[Importance sociale de l'artisanat]
L'importance sociale de l'artisanat désigne l'ensemble des rôles que joue ce secteur dans la vie de la population : création d'\cle{emplois}, insertion professionnelle des jeunes et des non-diplômés, fixation des populations rurales et réduction de la pauvreté.
\end{definition}

\textbf{Intuition.} Dans un quartier de Douala ou de Yaoundé, comptez les ateliers : menuisiers, couturiers, coiffeurs, mécaniciens, soudeurs, cordonniers, tailleurs... Chacun emploie un patron, des ouvriers et des apprentis. Multipliez ce constat par des milliers de quartiers et de villages : vous comprenez que l'artisanat est le plus grand « recruteur » du pays, bien devant l'administration ou les grandes entreprises.

\textbf{Un rôle d'employeur majeur.} L'artisanat est le cœur du \cle{secteur informel}, c'est-à-dire l'ensemble des activités économiques non déclarées ou peu structurées. Or ce secteur informel est le \cle{premier employeur} du Cameroun : il occupe la grande majorité de la population active, loin devant l'agriculture moderne, l'industrie ou la fonction publique.

\begin{exemplebox}[Un exemple chiffré : 9 actifs sur 10]
Selon les estimations de l'Institut National de la Statistique (INS) et de l'Organisation Internationale du Travail (OIT) issues des enquêtes sur l'emploi et le secteur informel (EESI 2010, EESI 2021), le secteur informel — dominé par l'artisanat et le petit commerce — occupe environ \textbf{90 \%} de la population active occupée, soit près de \textbf{9 actifs sur 10}. Les emplois salariés du secteur moderne (administration, grandes entreprises) ne concernent qu'une minorité (environ 10 \%). L'artisanat est donc, en volume, le premier pourvoyeur d'emplois du pays.
\end{exemplebox}

\textbf{L'insertion des jeunes et des non-diplômés.} L'artisanat offre une voie d'\cle{insertion professionnelle} aux jeunes sortis tôt du système scolaire : l'\cle{apprentissage} sur le tas, chez un maître-artisan, permet d'acquérir un métier en deux à cinq ans, sans diplôme exigé à l'entrée. C'est une réponse concrète au \cle{chômage} et au sous-emploi des jeunes, dont le taux dépasse 13 \% chez les 15--24 ans (INS, 2021).

\textbf{La lutte contre l'exode rural.} En offrant des revenus dans les villages et les petites villes (poterie, vannerie, tissage, menuiserie rurale), l'artisanat \cle{fixe les populations} sur place et freine l'\cle{exode rural} vers Douala et Yaoundé, déjà surchargées.

\begin{popfigure}
\begin{center}
\begin{tikzpicture}
\begin{axis}[
  ybar, bar width=0.7cm,
  width=0.95\linewidth, height=6cm,
  ymin=0, ymax=100,
  ylabel={Part de l'emploi (\%)},
  symbolic x coords={Informel,AgriModerne,Industrie,AdminFormel},
  xtick=data,
  xticklabels={Informel (artisanat, petit commerce),Agriculture moderne,Industrie,Administration et services formels},
  x tick label style={rotate=20, anchor=east, font=\small},
  nodes near coords, nodes near coords style={font=\small},
  every axis plot/.append style={fill=popblue, draw=popdark},
]
\addplot coordinates {(Informel,90) (AgriModerne,4) (Industrie,3) (AdminFormel,3)};
\end{axis}
\end{tikzpicture}
\end{center}
\legende{Répartition indicative de la population active occupée au Cameroun par grand secteur (estimations INS/OIT, en \%). Le secteur informel, dont l'artisanat est la composante principale, écrase les autres secteurs. Source : ordres de grandeur issus des enquêtes camerounaises sur l'emploi (EESI) et rapports OIT.}
\end{popfigure}

\courssub{L'importance économique : produire, sous-traiter, faire circuler les revenus}

\begin{definition}[Importance économique de l'artisanat]
L'importance économique de l'artisanat se mesure par sa contribution au \cle{PIB} (Produit Intérieur Brut, c'est-à-dire la richesse produite par le pays en un an), par sa production de biens de consommation accessibles et par ses liens avec les autres secteurs (industrie, BTP, commerce).
\end{definition}

\textbf{Une contribution réelle au PIB.} Même si une partie de l'activité artisanale échappe aux statistiques officielles (car informelle), le secteur informel dans son ensemble représente une part importante du PIB camerounais — les estimations (INS, Comptes Nationaux, FMI) lui attribuent couramment de l'ordre de \textbf{25 \% à 35 \% du PIB}. L'artisanat y contribue par la production de meubles, d'habillement, d'agroalimentaire de proximité, de réparations et de services.

\textbf{Des biens de consommation accessibles.} L'artisanat produit des biens \cle{adaptés aux revenus} des ménages modestes : meubles en bois local moins chers que les meubles importés, pagnes et vêtements cousus sur mesure, ustensiles de cuisine, chaussures réparées plutôt que rachetées. Il joue ainsi un rôle de \cle{substitution aux importations} : ce qui est fabriqué localement n'a pas besoin d'être acheté à l'étranger en devises.

\textbf{La sous-traitance pour l'industrie et le BTP.} Les artisans travaillent aussi \emph{pour} les grandes entreprises : c'est la \cle{sous-traitance}. Menuisiers et ferronniers interviennent sur les chantiers du \cle{BTP} (bâtiments et travaux publics) ; les mécaniciens et soudeurs entretiennent les flottes de véhicules des sociétés ; les couturiers confectionnent les uniformes. L'artisanat est donc un maillon indispensable de la chaîne de production nationale.

\begin{popfigure}
\begin{center}
\begin{tikzpicture}[font=\small,
  box/.style={draw=popdark, fill=popblueL, rounded corners=2pt, minimum width=3.1cm, minimum height=0.95cm, align=center},
  centre/.style={draw=popdark, fill=poporangeL, rounded corners=2pt, minimum width=3.4cm, minimum height=1.1cm, align=center, font=\small\bfseries},
  fleche/.style={-{Stealth[length=3mm]}, thick, draw=popdark}]
\node[centre, align=center] (A) at (0,0){ARTISANAT\\camerounais};
\node[box, align=center] (E) at (-4.6,2.2){Emplois\\(informel, jeunes)};
\node[box, align=center] (R) at (4.6,2.2){Revenus des ménages\\(lutte contre la pauvreté)};
\node[box, align=center] (T) at (-4.6,-2.2){Tourisme\\(marchés, villages d'artisans)};
\node[box, align=center] (P) at (0,-3.1){Patrimoine culturel\\(savoir-faire préservés)};
\node[box, align=center] (D) at (4.6,-2.2){Devises\\(exportations d'art)};
\draw[fleche] (A) -- (E);
\draw[fleche] (A) -- (R);
\draw[fleche] (A) -- (T);
\draw[fleche] (A) -- (P);
\draw[fleche] (A) -- (D);
\end{tikzpicture}
\end{center}
\legende{Schéma fléché des retombées de l'artisanat camerounais : 1. emplois ; 2. revenus ; 3. tourisme ; 4. patrimoine ; 5. devises. Ces cinq retombées se renforcent mutuellement.}
\end{popfigure}

\courssub{L'importance culturelle et touristique}

\textbf{La préservation du patrimoine.} Les artisans sont les dépositaires de \cle{savoir-faire} transmis de génération en génération : sculpture sur bois et sur bronze de l'Ouest (Foumban, Bafoussam), vannerie et poterie (Baham, Batcham), tissage du ndop chez les Bamiléké et les Bamoun, travail du cuir et bijoux du Grand-Nord (Maroua, Garoua). Sans artisans, ce \cle{patrimoine culturel} disparaîtrait. L'artisanat est donc un conservatoire vivant des identités culturelles camerounaises.

\textbf{Une attraction touristique.} Les \cle{marchés artisanaux} et les \cle{villages d'artisans} (Foumban et son Palais des rois Bamoun, Bafoussam, Bamenda, Maroua) attirent les visiteurs nationaux et étrangers. Le touriste qui achète un masque, une statuette ou un pagne injecte directement de l'argent dans l'économie locale : hébergement, restauration, transport et achats d'objets d'art.

\begin{savaistu}[Des marchés qui rayonnent sur toute l'Afrique centrale]
Les marchés artisanaux de Yaoundé (notamment le marché de Tsinga) et de Douala (marché des fleurs, marché Congo) attirent des acheteurs venus de toute l'Afrique centrale : Gabon, Congo, Guinée équatoriale, Tchad, Centrafrique. Revendeurs et touristes y achètent sculptures, pagnes, bijoux et vannerie, qu'ils revendent ensuite dans leurs pays. L'artisanat camerounais est ainsi un produit d'exportation régional informel mais structurant.
\end{savaistu}

\begin{attention}[Ne pas réduire l'artisanat au tourisme]
Erreur fréquente dans les copies : présenter l'artisanat uniquement comme une attraction pour les touristes. C'est faux et réducteur. Le rôle \textbf{majeur} de l'artisanat est d'abord \cle{l'emploi} (premier employeur du pays, ~90 \% des actifs) et la \cle{production de proximité} (biens et services accessibles aux ménages). La dimension touristique et culturelle est réelle mais secondaire en volume : commencez toujours votre analyse par l'emploi et la production.
\end{attention}

\courssub{L'importance fiscale et monétaire}

\textbf{Des devises grâce aux exportations.} L'\cle{artisanat d'art} (sculptures, masques, bronzes, pagnes, bijoux) est exporté vers l'Afrique centrale, l'Europe et l'Amérique du Nord. Ces ventes rapportent des \cle{devises} (euros, dollars), monnaies indispensables au pays pour importer ce qu'il ne produit pas. Même modeste à l'échelle nationale (quelques milliards de FCFA/an), cette entrée de devises est vitale pour les artisans exportateurs et leurs filières.

\textbf{Des recettes pour les collectivités.} Même informel, l'artisanat contribue aux finances publiques : \cle{patentes} et licences pour les ateliers déclarés, taxes de marché et tickets de place perçus par les \cle{communes}, impôt libératoire pour les petites unités formalisées. Ces recettes alimentent les budgets des collectivités territoriales décentralisées (voirie, marchés, écoles, santé).

\courssub{Exploiter des données chiffrées : démarche d'analyse}

Au Probatoire et au Baccalauréat technique, on vous donnera souvent un tableau ou un graphique sur l'emploi artisanal ou la contribution au PIB, avec la consigne : « Analysez ce document ». Voici la méthode attendue.

\begin{methode}[Analyser un tableau statistique en 4 étapes]
\begin{enumerate}
\item \textbf{Lire le document} : titre, unité (\%, milliers de personnes, milliards de FCFA), source et date.
\item \textbf{Dégager la tendance générale} en une phrase : hausse, baisse, domination d'un secteur (ex. : « le secteur informel domine largement l'emploi »).
\item \textbf{Citer deux chiffres précis} du document pour appuyer la tendance (ex. : « environ 90 \% des actifs contre moins de 5 \% dans l'industrie »).
\item \textbf{Conclure} en reliant le chiffre à la notion du cours (ex. : « l'artisanat est donc le premier employeur du Cameroun »).
\end{enumerate}
\end{methode}

\begin{exoresolu}[Analyse d'un document statistique]
\textbf{Énoncé.} Document : « Répartition indicative de la population active occupée au Cameroun : secteur informel (artisanat, petit commerce, services de rue) : environ 90 \% ; agriculture moderne : environ 4 \% ; industrie : environ 3 \% ; administration et services formels : environ 3 \%. » Question : dégagez l'importance sociale de l'artisanat à partir de ce document.
\tcblower
\textbf{Solution détaillée.}
\begin{enumerate}
\item \emph{Lecture} : le document donne la répartition de la population active occupée, en pourcentage, par grand secteur.
\item \emph{Tendance générale} : l'emploi camerounais est massivement informel ; le secteur moderne (industrie, administration) est marginal en volume.
\item \emph{Deux chiffres d'appui} : le secteur informel occupe environ 90 \% des actifs, soit 9 actifs sur 10, contre seulement 3 \% environ pour l'industrie.
\item \emph{Conclusion} : l'artisanat, qui constitue le cœur du secteur informel, est le premier employeur du Cameroun. Il assure l'insertion des jeunes et des non-diplômés, lutte contre le chômage et, en offrant des revenus en milieu rural, freine l'exode vers les villes. Son importance sociale est donc capitale.
\end{enumerate}
\end{exoresolu}

\begin{exercice}[À vous de jouer]
Un rapport indique : « Le secteur informel contribue pour environ 30 \% au PIB du Cameroun et les exportations d'artisanat d'art rapportent plusieurs milliards de FCFA de devises par an. » 1) Dégagez la tendance générale. 2) Citez deux chiffres du texte. 3) Concluez sur l'importance économique et monétaire de l'artisanat.
\end{exercice}

\begin{corrige}[Exercice]
1) \emph{Tendance} : l'artisanat et le secteur informel pèsent lourd dans l'économie nationale, à la fois dans la richesse produite et dans les recettes extérieures. 2) \emph{Chiffres} : environ 30 \% du PIB pour le secteur informel ; plusieurs milliards de FCFA de devises pour l'artisanat d'art exporté. 3) \emph{Conclusion} : l'artisanat contribue significativement au PIB par la production de biens et services de proximité, et il rapporte des devises grâce aux exportations d'objets d'art. Son importance est donc à la fois économique (production, sous-traitance) et monétaire (devises, recettes fiscales des collectivités).
\end{corrige}

\begin{aretenir}[L'essentiel]
\begin{itemize}
\item \textbf{Social} : l'artisanat est le premier employeur du Cameroun (secteur informel : environ 90 \% des actifs, soit 9 sur 10) ; il insère les jeunes et les non-diplômés et freine l'exode rural.
\item \textbf{Économique} : il contribue au PIB (environ 25--35 \% avec l'informel), produit des biens accessibles et sous-traite pour l'industrie et le BTP.
\item \textbf{Culturel et touristique} : il préserve le patrimoine (sculpture, ndop, vannerie, cuir) et attire les touristes (marchés de Tsinga à Yaoundé, villages d'artisans de Foumban).
\item \textbf{Fiscal et monétaire} : exportations d'artisanat d'art = devises ; patentes et taxes de marché = recettes des communes.
\item \textbf{Méthode} : analyser un document = tendance générale + deux chiffres précis + conclusion reliée au cours.
\end{itemize}
\end{aretenir}

\coursec{Les problèmes de l'artisanat camerounais}

Nous avons vu dans les sections précédentes que l'artisanat joue un rôle considérable au Cameroun : il fait vivre des centaines de milliers de familles, forme la jeunesse, fournit des biens et services de proximité et conserve des savoir-faire précieux. Pourtant, ce secteur si utile reste fragile. Malgré son importance, l'artisan camerounais travaille souvent dans des conditions difficiles qui l'empêchent de se développer. Comprendre ces difficultés, c'est préparer la recherche de solutions, qui fera l'objet du travail dirigé suivant (TD2). Cette section recense et classe les problèmes de l'artisanat camerounais, puis propose une méthode pour les hiérarchiser.

\courssub{Problèmes financiers : le manque d'argent pour investir}

L'intuition est simple : pour acheter une machine, un stock de matière première ou agrandir son atelier, il faut de l'argent. Or l'artisan camerounais part presque toujours de très peu.

\begin{definition}[Problèmes financiers de l'artisanat]
Les \cle{problèmes financiers} désignent l'ensemble des difficultés d'accès aux capitaux (fonds propres, crédits) qui empêchent l'artisan d'investir dans son outillage, ses matières premières et son local.
\end{definition}

Trois faits dominants caractérisent cette situation :

\begin{itemize}
\item \textbf{La difficulté d'accès au crédit bancaire.} Les banques exigent des garanties (titre foncier, salaire régulier, comptabilité tenue) que l'artisan informel ne possède pas. Elles jugent son activité trop risquée. Résultat : la grande majorité des artisans n'a jamais obtenu de prêt bancaire.
\item \textbf{La faiblesse des capitaux propres.} L'artisan démarre avec son épargne personnelle, souvent minime. Ses faibles bénéfices sont absorbés par les dépenses quotidiennes de la famille : il ne peut donc pas accumuler d'épargne pour investir.
\item \textbf{Le recours aux tontines.} Pour contourner le refus des banques, les artisans se tournent vers des formes traditionnelles de financement.
\end{itemize}

\begin{definition}[La tontine]
Une \cle{tontine} est une association d'épargne et de crédit rotatif entre personnes de confiance : chaque membre cotise régulièrement une somme fixe, et le total collecté est remis à tour de rôle à l'un des membres. Très utilisée par les artisans camerounais, elle permet de financer un achat important sans passer par la banque.
\end{definition}

\begin{exemplebox}[Un exemple chiffré : financer sans banque]
Les enquêtes de terrain menées auprès des artisans camerounais montrent que la majorité d'entre eux déclarent n'avoir \textbf{jamais obtenu de prêt bancaire}. Le financement repose sur trois sources principales : l'épargne personnelle, les tontines et l'aide familiale. Par exemple, un menuisier de Bafoussam qui souhaite acheter une scie circulaire à \num{150000} FCFA cotise \num{10000} FCFA par semaine dans une tontine de 15 membres : lorsque son tour arrive, il reçoit \num{150000} FCFA d'un coup et peut acheter sa machine.
\end{exemplebox}

\courssub{Problèmes techniques : un outillage rudimentaire}

Sans argent, pas de machines modernes. Le deuxième groupe de problèmes est donc la conséquence directe du premier.

\begin{itemize}
\item \textbf{Outillage rudimentaire.} Beaucoup d'artisans travaillent encore avec des outils manuels (rabot, marteau, forge au charbon) ou des machines d'occasion vétustes, souvent importées après des années de service en Europe.
\item \textbf{Faible productivité.} Avec un outillage manuel, la production est lente : un menuisier équipé d'outils électriques produit en une journée ce qu'un artisan non équipé réalise en une semaine. La \cle{productivité} (quantité produite par travailleur et par unité de temps) reste donc faible.
\item \textbf{Qualité irrégulière des produits.} L'absence de machines de précision et de normes de fabrication entraîne des finitions inégales : deux chaises du même atelier peuvent présenter des dimensions différentes. Cette irrégularité nuit à la réputation des produits artisanaux face aux produits industriels standardisés.
\end{itemize}

\begin{attention}[Ne pas confondre cause et conséquence]
Une erreur fréquente consiste à citer « outillage rudimentaire » comme un problème isolé. En réalité, il est souvent la \textbf{conséquence} du problème financier : l'artisan n'a pas choisi ses vieux outils, il n'a pas les moyens d'en acheter de nouveaux. Dans une rédaction, relie toujours les problèmes entre eux.
\end{attention}

\courssub{Problèmes de formation : un savoir-faire menacé}

\begin{itemize}
\item \textbf{Apprentissage informel non certifié.} La formation se fait presque toujours « sur le tas », chez un patron d'atelier. L'apprenti acquiert des gestes, mais ne reçoit \textbf{aucun diplôme} reconnu : il ne peut pas prouver officiellement sa qualification, ce qui limite sa mobilité et sa crédibilité.
\item \textbf{Perte de savoir-faire.} Certains métiers traditionnels (poterie, vannerie, tissage, forge) disparaissent avec les anciens maîtres artisans, car les techniques ne sont ni écrites ni enseignées dans des écoles.
\item \textbf{Désintérêt d'une partie de la jeunesse.} Beaucoup de jeunes rêvent de fonctions publiques ou de métiers « modernes » et considèrent l'artisanat comme un métier de pauvres, sale et sans avenir. Les ateliers peinent ainsi à trouver des relèves motivées.
\end{itemize}

\courssub{Problèmes commerciaux : vendre est un combat}

Produire ne suffit pas : il faut vendre. Or le marché est difficile.

\begin{itemize}
\item \textbf{Concurrence des produits industriels importés.} Les chaussures, meubles, tissus et ustensiles importés d'Asie ou d'Europe, souvent bon marché et bien présentés, concurrencent directement les produits artisanaux locaux.
\item \textbf{Faible accès aux marchés.} L'artisan vend surtout devant son atelier ou sur les marchés locaux. Il n'a ni boutique en ville, ni accès aux grandes surfaces, ni présence sur Internet : sa clientèle reste limitée à son quartier.
\item \textbf{Le poids des intermédiaires.} Pour atteindre les marchés lointains, l'artisan passe par des revendeurs qui achètent à bas prix et revendent cher : une part importante du bénéfice échappe au producteur.
\item \textbf{La contrebande.} Des produits importés frauduleusement, sans payer les droits de douane, arrivent sur le marché à des prix imbattables, aggravant la concurrence déloyale.
\end{itemize}

\courssub{Problèmes institutionnels : un cadre peu favorable}

Enfin, l'environnement administratif et matériel pèse sur l'artisan.

\begin{itemize}
\item \textbf{L'informalité.} La plupart des ateliers ne sont pas immatriculés au registre de commerce : ils échappent aux statistiques, ne peuvent pas obtenir de marchés publics et restent invisibles pour les programmes d'appui de l'État.
\item \textbf{Une fiscalité perçue comme lourde.} Les artisans qui se formalisent dénoncent la multiplicité des taxes (impôt libératoire, taxes communales, redevances diverses) qu'ils jugent disproportionnées par rapport à leurs revenus.
\item \textbf{Faible organisation professionnelle.} Les groupements et coopératives d'artisans existent mais restent peu nombreux et peu influents : isolés, les artisans ne peuvent pas défendre collectivement leurs intérêts ni négocier avec les pouvoirs publics.
\item \textbf{Accès difficile à l'énergie et aux infrastructures.} Les coupures fréquentes d'électricité paralysent les machines, le coût du raccordement est élevé, et les mauvaises routes encombrent le transport des produits vers les marchés.
\end{itemize}

\begin{attention}[Erreur fréquente : citer sans classer]
Au Probatoire et au Baccalauréat technique, lister dix problèmes en vrac ne suffit pas. Il faut les \textbf{classer} et distinguer :
\begin{itemize}
\item les \cle{problèmes internes}, qui relèvent de l'artisan lui-même et de son atelier (gestion, outillage, formation, qualité) ;
\item les \cle{problèmes externes}, qui relèvent de l'environnement (concurrence importée, fiscalité, contrebande, infrastructures, accès au crédit).
\end{itemize}
Cette distinction est essentielle : les problèmes internes peuvent être résolus par l'artisan et ses organisations, les problèmes externes exigent l'action de l'État.
\end{attention}

\courssub{Vue d'ensemble : le diagramme d'Ishikawa}

Pour visualiser l'ensemble des causes des difficultés de l'artisanat, on utilise un outil appelé \cle{diagramme en arêtes de poisson} (ou diagramme d'Ishikawa) : chaque « arête » représente une famille de causes qui convergent vers le même effet, la faiblesse du secteur artisanal.

\begin{popfigure}
\begin{tikzpicture}[scale=0.92, transform shape]
% Colonne vertébrale (jusqu'à la boîte effet)
\draw[line width=1.6pt, popdark, -{Stealth[length=4mm]}] (-7,0) -- (7.5,0);
% Tête (effet)
\node[draw=popink, fill=poppinkL, rounded corners, align=center, font=\bfseries\small, inner sep=5pt] at (9.5,0) {Difficultés et\\faiblesse de\\l'artisanat};
% Arêtes supérieures
\draw[line width=1.1pt, popblue] (-4.6,2.6) -- (-3.2,0);
\node[draw=popblue, fill=popblueL, rounded corners, font=\small\bfseries, inner sep=3pt] at (-4.9,2.9) {Financement};
\node[font=\scriptsize, align=left, text=popdark] at (-4.9,2.1) {pas de crédit bancaire\\capitaux propres faibles};
\draw[line width=1.1pt, popblue] (-0.6,2.6) -- (0.6,0);
\node[draw=popblue, fill=popblueL, rounded corners, font=\small\bfseries, inner sep=3pt] at (-0.9,2.9) {Technique};
\node[font=\scriptsize, align=left, text=popdark] at (-0.9,2.1) {outillage rudimentaire\\faible productivité};
\draw[line width=1.1pt, popblue] (3.4,2.6) -- (4.6,0);
\node[draw=popblue, fill=popblueL, rounded corners, font=\small\bfseries, inner sep=3pt] at (3.1,2.9) {Formation};
\node[font=\scriptsize, align=left, text=popdark] at (3.1,2.1) {apprentissage informel\\perte de savoir-faire};
% Arêtes inférieures
\draw[line width=1.1pt, poporange] (-2.6,-2.6) -- (-1.2,0);
\node[draw=poporange, fill=poporangeL, rounded corners, font=\small\bfseries, inner sep=3pt] at (-2.9,-2.9) {Marché};
\node[font=\scriptsize, align=left, text=popdark] at (-2.9,-3.6) {concurrence importée\\intermédiaires, contrebande};
\draw[line width=1.1pt, poporange] (1.4,-2.6) -- (2.6,0);
\node[draw=poporange, fill=poporangeL, rounded corners, font=\small\bfseries, inner sep=3pt] at (1.1,-2.9) {Institutions};
\node[font=\scriptsize, align=left, text=popdark] at (1.1,-3.6) {informalité, fiscalité\\énergie, infrastructures};
\end{tikzpicture}
\legende{Diagramme d'Ishikawa (arêtes de poisson) des difficultés de l'artisanat camerounais. Chaque arête regroupe une famille de causes : financement, technique, formation, marché, institutions. Toutes convergent vers l'effet final : la faiblesse du secteur artisanal.}
\end{popfigure}

\begin{popfigure}
\begin{tikzpicture}[scale=0.95, transform shape]
% Sol
\fill[popcream] (0,0) rectangle (10,0.25);
\draw[popdark] (0,0.25) -- (10,0.25);
% Hangar
\draw[fill=popblueL, draw=popdark] (1,0.25) rectangle (7,3.2);
\draw[fill=popgoldL, draw=popdark] (0.7,3.2) -- (4,4.3) -- (7.3,3.2) -- cycle;
% Porte ouverte
\draw[fill=popdark!60] (3.2,0.25) rectangle (4.8,2.4);
% Véhicule en réparation (schéma)
\draw[fill=popmuted!40, draw=popdark, rounded corners] (5.4,0.6) rectangle (8.9,1.4);
\draw[fill=popmuted!40, draw=popdark, rounded corners] (6.1,1.4) rectangle (8.1,2.0);
\draw[fill=popdark] (6.0,0.55) circle (0.32);
\draw[fill=popdark] (8.3,0.55) circle (0.32);
% Mécanicien sous le véhicule
\draw[fill=popgreen, draw=popdark] (6.6,0.28) rectangle (7.8,0.5);
% Outils épars
\draw[popink, line width=1pt] (1.4,0.4) -- (2.0,0.4);
\draw[popink, line width=1pt] (2.3,0.35) -- (2.3,0.6);
\draw[popink, line width=1pt] (2.2,0.48) -- (2.4,0.48);
% Bidons
\draw[fill=poporangeL, draw=popdark] (0.6,0.25) rectangle (1.0,0.9);
\draw[fill=poporangeL, draw=popdark] (9.2,0.25) rectangle (9.6,0.85);
% Flèches de légende
\draw[-{Stealth}, popink, thick] (8.9,3.6) -- (8.2,2.1);
\node[font=\scriptsize, align=left, text=popink, anchor=west] at (8.6,3.8) {1. véhicule réparé\\à même le sol};
\draw[-{Stealth}, popink, thick] (0.5,2.6) -- (1.6,0.5);
\node[font=\scriptsize, align=left, text=popink, anchor=east] at (0.5,2.8) {2. outils rudimentaires\\épars sur le sol};
\draw[-{Stealth}, popink, thick] (4.0,5.0) -- (4.0,4.35);
\node[font=\scriptsize, align=center, text=popink] at (4.0,5.3) {3. toit de tôle, atelier ouvert\\sur la rue, sans équipement de sécurité};
\end{tikzpicture}
\legende{Croquis d'après photographie : un atelier de mécanique informelle (le « quartier général »). On observe : 1. un véhicule réparé sans pont élévateur ; 2. un outillage sommaire posé au sol ; 3. un local précaire ouvert sur la rue. Ces conditions illustrent les problèmes techniques (outillage rudimentaire) et institutionnels (informalité, absence d'infrastructures) de l'artisanat camerounais.}
\end{popfigure}

\courssub{Hiérarchiser les problèmes : une activité guidée}

Tous ces problèmes n'ont pas le même poids. Pour agir efficacement, il faut savoir les \cle{hiérarchiser}, c'est-à-dire les ranger du plus important au moins important.

\begin{methode}[Hiérarchiser des problèmes à partir d'une enquête]
\begin{enumerate}
\item \textbf{Recueillir les données} : interroger des artisans (questionnaire, interview) ou exploiter des documents (rapports, articles, statistiques).
\item \textbf{Classer les problèmes cités} dans les cinq familles vues ci-dessus (financement, technique, formation, marché, institutions).
\item \textbf{Croiser fréquence et gravité} : pour chaque problème, noter combien d'artisans le citent (fréquence) et à quel point il bloque l'activité (gravité, par exemple de 1 à 3).
\item \textbf{Établir le classement} : le problème le plus fréquent ET le plus grave arrive en tête.
\item \textbf{Justifier en une phrase} : « Le problème X arrive en premier car il est cité par ... artisans sur ... et il empêche ... ».
\end{enumerate}
\end{methode}

\begin{exoresolu}[Hiérarchiser les problèmes d'un quartier]
Énoncé. Une enquête menée auprès de 20 artisans d'un quartier de Douala donne les résultats suivants : le manque de crédit est cité par 16 artisans (gravité 3), les coupures d'électricité par 12 artisans (gravité 2), la concurrence des produits importés par 9 artisans (gravité 2), l'outillage rudimentaire par 7 artisans (gravité 2), la faible organisation professionnelle par 3 artisans (gravité 1). Classe ces problèmes du plus important au moins important et justifie le premier rang.
\tcblower
Solution détaillée. On croise la fréquence (nombre de citations) et la gravité (sur 3) :
\begin{center}
\begin{tabularx}{\linewidth}{|X|c|c|c|}\hline
\rowcolor{popblueL} \textbf{Problème} & \textbf{Fréquence} & \textbf{Gravité} & \textbf{Rang} \\ \hline
Manque de crédit & 16/20 & 3 & 1 \\ \hline
Coupures d'électricité & 12/20 & 2 & 2 \\ \hline
Concurrence des importations & 9/20 & 2 & 3 \\ \hline
Outillage rudimentaire & 7/20 & 2 & 4 \\ \hline
Faible organisation professionnelle & 3/20 & 1 & 5 \\ \hline
\end{tabularx}
\end{center}
Justification du premier rang : le manque de crédit arrive en tête car il est cité par 16 artisans sur 20 (soit 80 \% des enquêtés) et sa gravité est maximale (3/3) : sans financement, l'artisan ne peut ni moderniser son outillage ni acheter ses matières premières, ce qui aggrave plusieurs autres problèmes.
\end{exoresolu}

\begin{aretenir}[L'essentiel de la section]
Les problèmes de l'artisanat camerounais se classent en cinq familles : \textbf{financiers} (pas de crédit bancaire, capitaux propres faibles, recours aux tontines), \textbf{techniques} (outillage rudimentaire, faible productivité, qualité irrégulière), \textbf{de formation} (apprentissage informel non certifié, perte de savoir-faire, désintérêt des jeunes), \textbf{commerciaux} (concurrence des importations, intermédiaires, contrebande) et \textbf{institutionnels} (informalité, fiscalité lourde, faible organisation, énergie et infrastructures déficientes). On distingue les problèmes \textbf{internes} (atelier) des problèmes \textbf{externes} (environnement). Pour les hiérarchiser, on croise leur fréquence et leur gravité, puis on justifie le classement.
\end{aretenir}

\begin{exercice}[À toi de jouer]
Dans ton quartier ou ton village, identifie trois ateliers artisanaux (menuiserie, couture, mécanique, coiffure...). Pour chacun, relève deux problèmes observés, classe-les dans les cinq familles de la leçon, puis indique s'il s'agit d'un problème interne ou externe. Propose enfin un classement hiérarchique justifié en une phrase.
\end{exercice}

\begin{corrige}[Exercice]
Exemple de réponse attendue (à adapter selon le terrain). Atelier de couture : machine à coudre manuelle vétuste (problème technique, interne) et concurrence des friperies importées (problème commercial, externe). Atelier de mécanique : absence de crédit pour acheter un pont élévateur (problème financier, externe) et coupures d'électricité (problème institutionnel, externe). Salon de coiffure : apprenties non certifiées (problème de formation, interne) et taxes multiples (problème institutionnel, externe). Classement : 1. absence de crédit, car elle bloque l'achat de tout équipement et aggrave les autres problèmes ; 2. coupures d'électricité, car elles arrêtent le travail plusieurs heures par jour ; 3. concurrence des friperies, car elle réduit la clientèle du tailleur.
\end{corrige}

\coursec{TD2 — Les solutions aux problèmes de l'artisanat}

Dans la section précédente, nous avons dressé le diagnostic : l'artisanat camerounais souffre de difficultés de financement, d'un outillage vieillissant, d'une formation insuffisante, de l'informalité, de la concurrence des produits importés et de la contrefaçon. Un bon diagnostic n'a de valeur que s'il débouche sur des remèdes. C'est l'objet de ce travail dirigé : à chaque problème identifié, associer une solution concrète, réaliste et portée par un acteur précis. Nous terminerons par une conclusion argumentée appliquée au cas de l'atelier de sculpture de Foumban présenté en accroche de la leçon.

\courssub{Les solutions financières : donner des moyens aux artisans}

\textbf{L'intuition.} Un potier de Maroua qui gagne 2\,000 FCFA par jour ne peut pas économiser de quoi acheter un four moderne à 800\,000 FCFA. Sans épargne suffisante et sans accès au crédit bancaire classique (les banques exigent des garanties que l'artisan informel ne possède pas), il reste bloqué. La solution financière consiste donc à rapprocher l'argent de l'artisan.

\textbf{Les mécanismes.} Quatre pistes complémentaires existent :

\begin{itemize}
\item La \cle{microfinance} : les établissements de microfinance (EMF) comme les réseaux \cle{MC\textsuperscript{2}}, les \cle{COOPEC} (coopératives d'épargne et de crédit) ou le réseau \cle{CamCCUL} accordent de petits prêts (microcrédits) adaptés aux revenus modestes, souvent sans garantie immobilière, parfois sur la base de la caution solidaire du groupe.
\item Les \cle{subventions} et \cle{fonds d'appui} publics ou des partenaires au développement, qui financent l'équipement ou la formation sans exiger de remboursement intégral.
\item Les \cle{groupements d'artisans} (GIC, coopératives) qui mutualisent les moyens : à dix, on achète ensemble une machine, on négocie les matières premières en gros, on partage un local.
\item La \cle{tontine} traditionnelle, forme d'épargne rotative bien ancrée dans la culture camerounaise, qui reste un premier levier de financement.
\end{itemize}

\begin{exemplebox}[La microfinance en chiffres]
Le secteur de la microfinance au Cameroun (MC\textsuperscript{2}, COOPEC, CamCCUL, etc.) compte \textbf{plus de deux millions de membres} et finance chaque année de nombreux petits artisans : menuisiers, couturiers, coiffeurs, potiers. Un microcrédit de 150\,000 à 500\,000 FCFA suffit souvent à acheter une machine à coudre industrielle, un tour de potier ou un stock de bois. Le remboursement s'échelonne sur 6 à 24 mois, avec des échéances adaptées aux revenus irréguliers de l'artisan.
\end{exemplebox}

\begin{attention}[Piège à éviter]
Un prêt n'est pas une solution magique : il faut le rembourser. Un microcrédit mal calculé (échéances trop lourdes, investissement non rentable) peut enfoncer l'artisan dans l'endettement. La solution financière doit toujours s'accompagner d'une formation à la gestion.
\end{attention}

\courssub{Les solutions techniques : moderniser et former}

\textbf{L'intuition.} Un sculpteur qui travaille avec des outils émoussés produit lentement et moins finement qu'un concurrent équipé de machines électriques. La technique conditionne la qualité et la quantité.

\textbf{Les mécanismes.}

\begin{itemize}
\item La \cle{modernisation de l'outillage} : acquisition de machines simples et robustes (scies électriques, tours de potier, machines à coudre industrielles), adaptées au contexte (coupures d'électricité fréquentes, entretien local possible).
\item Les \cle{centres de formation et de perfectionnement} : les lycées techniques, les centres de formation professionnelle et les programmes de perfectionnement permettent aux artisans d'apprendre de nouvelles techniques, le dessin technique, la gestion et même l'usage d'Internet.
\item Les \cle{normes de qualité} : respecter des normes (dimensions, finitions, hygiène pour l'agroalimentaire artisanal) permet d'accéder à des marchés plus exigeants, y compris à l'exportation. L'\cle{ANOR} (Agence des normes et de la qualité) joue un rôle clé dans ce domaine.
\end{itemize}

\begin{experience}[Observer pour comprendre : deux ateliers de menuiserie]
\textbf{Protocole.} On observe deux menuisiers de Bafoussam pendant une semaine : l'un travaille entièrement à la main, l'autre dispose d'une scie circulaire et d'une ponceuse électriques.\par
\textbf{Observations.} Le premier fabrique une table en trois jours ; le second en fabrique une par jour, avec des finitions plus régulières, et vend au même prix.\par
\textbf{Interprétation.} À prix de vente égal, la productivité triple grâce à l'outillage moderne : le revenu mensuel peut être multiplié par deux ou trois. La modernisation technique est donc un levier direct d'augmentation des revenus, à condition que l'investissement initial soit financé (d'où le lien avec les solutions financières).
\end{experience}

\courssub{Les solutions institutionnelles : l'État aux côtés des artisans}

\begin{definition}[Le MINPMEESA]
Le \cle{MINPMEESA} (ministère des Petites et Moyennes Entreprises, de l'Économie sociale et de l'Artisanat) est l'organe d'État chargé de la promotion de l'artisanat au Cameroun. Il élabore la politique artisanale, délivre la carte professionnelle d'artisan, appuie la formation et organise la promotion des produits artisanaux.
\end{definition}

\textbf{Les mécanismes institutionnels.}

\begin{itemize}
\item La \cle{formalisation accompagnée} : l'État encourage les artisans à déclarer leur activité en simplifiant les démarches et en réduisant les coûts, car la formalisation donne accès au crédit, aux marchés publics et à la protection sociale.
\item Les \cle{guichets uniques} (Centres de formalités des entreprises -- CFE) : un seul bureau regroupe toutes les démarches de création d'entreprise, ce qui réduit les délais et les déplacements.
\item Le \cle{statut de l'artisan} et la carte professionnelle d'artisan : ils reconnaissent officiellement le métier et ouvrent des droits.
\item Les \cle{foires et salons artisanaux}, soutenus par l'État, qui offrent une vitrine aux productions locales.
\end{itemize}

\begin{savaistu}[Le SIARC, vitrine de l'artisanat camerounais]
Le Salon international de l'artisanat du Cameroun (SIARC), organisé régulièrement à Yaoundé, valorise les productions des artisans des dix régions : sculptures de Foumban, poteries du Nord, vannerie, tissage, bijoux, cuir. C'est une occasion unique pour les artisans de rencontrer des acheteurs, des exportateurs et des partenaires financiers.
\end{savaistu}

\begin{popfigure}
\begin{center}
\begin{tikzpicture}[
  box/.style={draw=popblue, fill=popblueL, rounded corners=2pt, text width=4.6cm, align=center, font=\small, minimum height=0.95cm},
  arr/.style={-{Stealth[length=2.5mm]}, thick, popdark}
]
\node[box, align=center] (a) at (0,0){\textbf{1. Décision} \\ l'artisan choisit de formaliser son atelier};
\node[box, align=center] (b) at (0,-1.7){\textbf{2. Guichet unique (CFE)} \\ dépôt du dossier, choix du statut};
\node[box, align=center] (c) at (0,-3.4){\textbf{3. Immatriculation} \\ registre de commerce, numéro contribuable};
\node[box, align=center] (d) at (0,-5.1){\textbf{4. Carte professionnelle d'artisan} \\ délivrée via le MINPMEESA};
\node[box, fill=popgreenL, draw=popgreen, align=center] (e) at (0,-6.8){\textbf{5. Atelier formalisé} \\ accès au crédit, aux marchés, aux salons};
\draw[arr] (a) -- (b);
\draw[arr] (b) -- (c);
\draw[arr] (c) -- (d);
\draw[arr] (d) -- (e);
\end{tikzpicture}
\end{center}
\legende{Organigramme : le parcours simplifié de formalisation d'un atelier artisanal au Cameroun (schéma indicatif).}
\end{popfigure}

\courssub{Les solutions commerciales : vendre mieux et plus loin}

\textbf{L'intuition.} Produire un beau pagne tissé ne suffit pas : encore faut-il trouver des clients. Beaucoup d'artisans vendent uniquement au bord de la route ou aux touristes de passage. Les solutions commerciales élargissent les débouchés.

\begin{itemize}
\item La promotion du \cle{made in Cameroun} : campagnes de valorisation des produits locaux, labels, journées de consommation locale, qui poussent les consommateurs camerounais à préférer les produits artisanaux nationaux aux importations.
\item Le \cle{e-commerce} : grâce au téléphone mobile, l'artisan peut photographier ses produits et les vendre sur les réseaux sociaux ou des plateformes de vente en ligne, touchant des clients à Douala, Yaoundé et même à l'étranger.
\item Les \cle{coopératives de vente} : les artisans se regroupent pour ouvrir des boutiques communes, participer collectivement aux foires et négocier de meilleurs prix.
\item La \cle{protection contre la contrefaçon} : marques déposées, labels d'origine et actions de l'État contre les imitations importées qui copient les motifs traditionnels.
\end{itemize}

\courssub{Le rôle des artisans eux-mêmes : les solutions viennent aussi d'en bas}

L'État et les partenaires ne peuvent pas tout. Les artisans sont les premiers acteurs de leur développement :

\begin{itemize}
\item S'\cle{organiser en GIC et coopératives} : l'union fait la force pour acheter, produire et vendre.
\item \cle{Améliorer la gestion} : tenir un cahier de comptes simple (recettes, dépenses, bénéfices), séparer l'argent de l'atelier de l'argent de la famille, fixer des prix qui couvrent les coûts.
\item \cle{Transmettre le savoir-faire} : former des apprentis garantit la survie des métiers (sculpture, forge, tissage) et évite la disparition des techniques traditionnelles.
\end{itemize}

\courssub{Tableau synoptique : problème, solution, acteur}

\begin{popfigure}
\begin{center}
\begin{tikzpicture}[
  pb/.style={draw=poppink, fill=poppinkL, rounded corners=2pt, text width=3.4cm, align=center, font=\footnotesize, minimum height=1.15cm},
  so/.style={draw=popblue, fill=popblueL, rounded corners=2pt, text width=4.2cm, align=center, font=\footnotesize, minimum height=1.15cm},
  ac/.style={draw=popgreen, fill=popgreenL, rounded corners=2pt, text width=3.4cm, align=center, font=\footnotesize, minimum height=1.15cm},
  arr/.style={-{Stealth[length=2.2mm]}, thick, popdark}
]
\node[font=\small\bfseries, poppink] at (-4.6,0.7) {PROBLÈME};
\node[font=\small\bfseries, popblue] at (0,0.7) {SOLUTION};
\node[font=\small\bfseries, popgreen] at (4.6,0.7) {ACTEUR};
\node[pb] (p1) at (-4.6,-0.8) {Manque de financement};
\node[so] (s1) at (0,-0.8) {Microcrédit, COOPEC, tontines, fonds d'appui};
\node[ac] (a1) at (4.6,-0.8) {EMF (MC\textsuperscript{2}, CamCCUL), État};
\node[pb] (p2) at (-4.6,-2.4) {Outillage vétuste};
\node[so] (s2) at (0,-2.4) {Modernisation, mutualisation du matériel};
\node[ac] (a2) at (4.6,-2.4) {GIC, coopératives, artisans};
\node[pb] (p3) at (-4.6,-4.0) {Formation insuffisante};
\node[so] (s3) at (0,-4.0) {Centres de formation, perfectionnement, apprentissage};
\node[ac] (a3) at (4.6,-4.0) {MINPMEESA, lycées techniques, maîtres artisans};
\node[pb] (p4) at (-4.6,-5.6) {Informalité};
\node[so] (s4) at (0,-5.6) {Guichets uniques, statut de l'artisan, carte professionnelle};
\node[ac] (a4) at (4.6,-5.6) {État (CFE, MINPMEESA)};
\node[pb] (p5) at (-4.6,-7.2) {Débouchés limités, contrefaçon};
\node[so] (s5) at (0,-7.2) {Made in Cameroun, e-commerce, salons, labels};
\node[ac] (a5) at (4.6,-7.2) {Artisans, coopératives, État};
\foreach \i in {1,...,5} { \draw[arr] (p\i) -- (s\i); \draw[arr] (s\i) -- (a\i); }
\end{tikzpicture}
\end{center}
\legende{Tableau synoptique : à chaque problème de l'artisanat camerounais correspond une solution portée par un acteur précis.}
\end{popfigure}

\begin{methode}[Justifier une solution en géographie]
Pour proposer et justifier une solution à un problème de l'artisanat, suis trois étapes :
\begin{enumerate}
\item \textbf{Rappeler le problème} précisément (ex. : « l'atelier n'a pas accès au crédit bancaire »).
\item \textbf{Relier explicitement la solution au problème} qu'elle résout (ex. : « le microcrédit d'une COOPEC fournit les fonds que la banque refuse »).
\item \textbf{Nommer l'acteur} qui la met en œuvre (ex. : « l'établissement de microfinance, appuyé par le MINPMEESA ») et vérifier qu'elle est \emph{réaliste} dans le contexte local.
\end{enumerate}
\end{methode}

\begin{attention}[Une solution doit être adaptée au contexte local]
Une solution réaliste doit être adaptée au contexte local (coût, accessibilité, culture) et non copiée d'un modèle étranger. Une machine solaire coûteuse importée d'Europe peut être inutilisable dans un village sans technicien de maintenance ; une tontine ou une COOPEC de quartier, elle, fonctionne parce qu'elle repose sur la confiance locale. Au Probatoire (et plus tard au Baccalauréat technique), une solution hors contexte est pénalisée.
\end{attention}

\courssub{Exercice résolu et conclusion argumentée : le cas de l'atelier de Foumban}

\begin{exoresolu}[Sauver l'atelier de sculpture de Foumban]
\textbf{Énoncé.} L'atelier de sculpture sur bois de M. Kamdem, à Foumban, emploie quatre apprentis. Il fait face à trois problèmes : (1) il ne peut pas acheter une scie électrique à 400\,000 FCFA ; (2) ses sculptures sont copiées par des produits importés bon marché ; (3) il n'est pas déclaré et ne peut pas participer aux salons officiels. Propose une solution justifiée pour chaque problème, puis rédige une courte conclusion indiquant la solution prioritaire.
\tcblower
\textbf{Solution détaillée.}
\begin{itemize}
\item \emph{Problème 1 (financement)} : M. Kamdem peut adhérer à la COOPEC de sa localité et solliciter un microcrédit de 400\,000 FCFA, remboursable sur 18 mois. Acteur : l'établissement de microfinance. La scie électrique augmentant sa productivité, le remboursement est couvert par les ventes supplémentaires.
\item \emph{Problème 2 (contrefaçon)} : il peut rejoindre la coopérative des sculpteurs de Foumban pour créer un label « authentique sculpture de Foumban » et vendre via les réseaux sociaux à des clients de Yaoundé, Douala et de la diaspora, prêts à payer le vrai prix de l'artisanat. Acteurs : les artisans regroupés, appuyés par l'État pour la protection des labels.
\item \emph{Problème 3 (informalité)} : il se rend au guichet unique (CFE) pour immatriculer son atelier, puis demande la carte professionnelle d'artisan auprès du MINPMEESA, ce qui lui ouvre l'accès au SIARC. Acteur : l'État.
\end{itemize}
\emph{Conclusion.} La solution prioritaire est la formalisation : c'est la clé qui débloque les autres (crédit plus facile, participation aux salons, protection). Mais elle ne suffit pas seule : c'est la combinaison des trois solutions, adaptées au contexte de Foumban, capitale artisanale de l'Ouest, qui assurera la survie et la croissance de l'atelier.
\end{exoresolu}

\begin{exercice}[À toi de jouer]
Une coopérative de potières de Maroua veut exporter ses poteries vers le Tchad voisin. Elle manque de fonds pour un four moderne (1\,200\,000 FCFA) et ses membres ne connaissent pas les normes d'emballage. Rédige un paragraphe de 8 à 10 lignes proposant deux solutions justifiées (problème → solution → acteur) et conclus en indiquant laquelle mettre en œuvre en premier.
\end{exercice}

\begin{corrige}[Exercice — potières de Maroua]
Premier problème : le manque de financement pour le four. Solution : la coopérative peut solliciter un microcrédit groupé auprès d'un établissement de microfinance du Grand-Nord (réseau MC\textsuperscript{2} ou COOPEC locale), la caution solidaire du groupe remplaçant les garanties bancaires ; elle peut aussi candidater à un fonds d'appui du MINPMEESA. Acteurs : l'EMF et l'État. Deuxième problème : la méconnaissance des normes d'emballage. Solution : suivre une session de perfectionnement dans un centre de formation, avec l'appui du MINPMEESA et de l'ANOR pour les normes. Acteurs : centres de formation, État. Conclusion : la formation est prioritaire, car un four financé sans maîtrise des normes ne permettrait pas d'exporter ; ensuite seulement, le microcrédit rendra l'investissement rentable. Ces deux solutions sont réalistes car elles s'appuient sur des institutions présentes dans la région.
\end{corrige}

\begin{aretenir}[L'essentiel du TD2]
Les solutions aux problèmes de l'artisanat camerounais sont de quatre ordres : \textbf{financières} (microfinance, COOPEC, fonds d'appui, groupements), \textbf{techniques} (modernisation, formation, normes), \textbf{institutionnelles} (MINPMEESA, guichets uniques, statut de l'artisan, salons comme le SIARC) et \textbf{commerciales} (made in Cameroun, e-commerce, coopératives de vente, lutte contre la contrefaçon). Les artisans eux-mêmes, organisés en GIC et coopératives, sont des acteurs essentiels. Toute solution doit être reliée au problème qu'elle résout, portée par un acteur nommé et adaptée au contexte local.
\end{aretenir}

\coursec{TP — Cartographie et représentation graphique de l'artisanat}

Dans la section précédente, nous avons recherché des \cle{solutions} aux problèmes de l'artisanat camerounais : formation, financement, organisation en coopératives, amélioration des débouchés. Mais pour proposer des solutions efficaces, encore faut-il \textbf{savoir où sont les artisans, combien ils sont et dans quelles branches ils travaillent}. C'est précisément le rôle des outils du géographe : la \cle{carte}, le \cle{diagramme} et le \cle{croquis}. Ce travail pratique (TP) vous apprend à construire ces trois documents à partir de données réelles, puis à les commenter — une compétence directement exigible au Probatoire et au Baccalauréat technique.

\courssub{TP1 : Construire une carte des pôles artisanaux du Cameroun}

\textbf{Intuition.} Une carte est un \emph{langage} : au lieu d'écrire « Foumban est un grand centre de sculpture et de bronze, Bali de la poterie, Bafoussam de la vannerie », on le \emph{montre}. Mais un langage a des règles : sans elles, la carte devient illisible ou mensongère.

\begin{definition}[Fond de carte muet et figuré]
Un \cle{fond de carte muet} est un contour géographique vierge (ici, le contour du Cameroun) que l'élève doit compléter. Un \cle{figuré} est un symbole (point, cercle, carré, couleur, flèche) qui représente une donnée sur la carte. Un figuré est dit \cle{proportionnel} lorsque sa taille varie avec la valeur représentée (ex. : un cercle deux fois plus grand pour deux fois plus d'artisans).
\end{definition}

\begin{methode}[Les quatre règles d'or d'une carte réussie]
Une carte de géographie comporte \textbf{obligatoirement} :
\begin{enumerate}
\item un \cle{titre} précis (quoi ? où ? quand ?) ;
\item une \cle{légende} complète, placée hors du contour, qui explique chaque figuré ;
\item une \cle{orientation} (flèche du nord) et, si possible, une échelle indicative ;
\item des \cle{figurés proportionnés} aux données et limités en nombre (4 figurés maximum pour rester lisible).
\end{enumerate}
\end{methode}

\textbf{Document de travail.} Le tableau ci-dessous fournit les données à reporter sur le fond de carte muet.

\begin{center}
\begin{tabularx}{\linewidth}{|l|X|c|}\hline
\rowcolor{popblueL} \textbf{Pôle artisanal} & \textbf{Branche dominante} & \textbf{Importance (nombre d'ateliers, estimation)} \\ \hline
Foumban (Ouest) & Sculpture, bronze, cuir & 300 \\ \hline
Bafoussam (Ouest) & Vannerie, tissage & 200 \\ \hline
Bali (Nord-Ouest) & Poterie & 100 \\ \hline
Maroua (Extrême-Nord) & Tannage, cuir, couture & 150 \\ \hline
Yaoundé (Centre) & Ebénisterie, mécanique artisanale & 250 \\ \hline
Douala (Littoral) & Mécanique, soudure, menuiserie & 250 \\ \hline
\end{tabularx}
\end{center}

\begin{exemplebox}[Choix des figurés : l'exemple chiffré]
À partir d'un tableau de 6 pôles et plusieurs branches, on choisit \textbf{au maximum 4 figurés} pour garder une carte lisible. Ici, on peut regrouper :
\begin{itemize}
\item une \textbf{couleur de point} par grande famille de branches (bronze/sculpture = orange, poterie = violet, cuir/tannage = vert, mécanique/ébénisterie/menuiserie = bleu, vannerie/tissage = pourpre) ;
\item une \textbf{taille de point proportionnelle} au nombre d'ateliers (300 ateliers = grand cercle, 100 = petit cercle).
\end{itemize}
On obtient ainsi deux informations (famille de branche + importance) avec un seul type de figuré.
\end{exemplebox}

\begin{popfigure}
\begin{center}
\begin{tikzpicture}[scale=0.9]
% Contour simplifié du Cameroun (fond muet)
\draw[thick,popdark,fill=popcream] (1.2,0) -- (3.4,0) -- (4.4,0.8) -- (4.6,2.2) -- (4.2,3.4) -- (4.6,4.6) -- (4.0,5.6) -- (3.0,6.4) -- (2.2,6.2) -- (1.6,5.2) -- (1.8,4.2) -- (1.2,3.4) -- (0.6,2.4) -- (0.4,1.2) -- cycle;
% Flèche du nord
\draw[->,thick,popdark] (5.3,5.4) -- (5.3,6.2) node[above]{\textbf{N}};
% Repères (villes) en points vides à compléter : 6 pôles
\foreach \p in {(2.0,1.0), (2.6,2.2), (1.4,3.6), (2.4,4.0), (3.2,5.4), (1.0,5.8)}
  \draw[popmuted,fill=white] \p circle (0.09);
% Titre
\node[align=center] at (2.5,7.0){\textbf{Fond de carte muet : les pôles artisanaux du Cameroun}};
\node[align=left,anchor=north west] at (5.0,4.4){\footnotesize \textbf{À faire :}};
\node[align=left,anchor=north west] at (5.0,4.1){\footnotesize 1. Placer les 6 pôles du tableau ;};
\node[align=left,anchor=north west] at (5.0,3.8){\footnotesize 2. Choisir un figuré coloré par famille de branche ;};
\node[align=left,anchor=north west] at (5.0,3.5){\footnotesize 3. Proportionner la taille à l'importance ;};
\node[align=left,anchor=north west] at (5.0,3.2){\footnotesize 4. Rédiger titre et légende.};
\end{tikzpicture}
\end{center}
\legende{Fond de carte muet à compléter. Contour très simplifié du Cameroun ; les 6 points gris indiquent les emplacements approximatifs des pôles à identifier (Douala et Yaoundé au sud, Bafoussam et Foumban à l'ouest, Bali au nord-ouest, Maroua au nord). Échelle indicative : 1 cm $\approx$ 150 km.}
\end{popfigure}

\begin{attention}[Erreur fréquente : la carte fausse]
L'erreur la plus grave est d'\textbf{oublier la légende} ou d'utiliser des \textbf{figurés non proportionnels} : si Foumban (300 ateliers) et Bali (100 ateliers) sont représentés par des cercles identiques, la carte affirme quelque chose de faux. Au Bac, une carte sans titre, sans légende ou sans nord est \textbf{systématiquement pénalisée}, même si les localisations sont exactes.
\end{attention}

\courssub{TP2 : Construire un diagramme de répartition des artisans}

\textbf{Intuition.} La carte dit \emph{où} ; le diagramme dit \emph{combien}. Deux diagrammes sont au programme : le \cle{diagramme circulaire} (ou « camembert »), adapté aux \textbf{parts d'un total} (en \%), et le \cle{diagramme en bâtons}, adapté à la \textbf{comparaison de valeurs} entre catégories.

\begin{propriete}[Calcul des angles d'un diagramme circulaire]
Chaque secteur a un angle proportionnel à sa part dans le total :
\[ \text{angle du secteur} = \frac{\text{effectif de la catégorie}}{\text{effectif total}} \times 360°. \]
La somme des angles doit faire exactement $360°$ (arrondir au degré le plus proche et vérifier la somme).
\end{propriete}

\begin{exoresolu}[Diagramme circulaire : répartition des artisans par branche]
\textbf{Énoncé.} Une coopérative de la région de l'Ouest compte 600 artisans répartis ainsi : vannerie 180, sculpture 150, poterie 90, couture 120, bijouterie 60. Construire le diagramme circulaire.
\tcblower
\textbf{Solution détaillée.} On calcule chaque angle :
\begin{align*}
\text{Vannerie : } & \frac{180}{600}\times 360 = 108° \\
\text{Sculpture : } & \frac{150}{600}\times 360 = 90° \\
\text{Poterie : } & \frac{90}{600}\times 360 = 54° \\
\text{Couture : } & \frac{120}{600}\times 360 = 72° \\
\text{Bijouterie : } & \frac{60}{600}\times 360 = 36°
\end{align*}
Vérification : $108+90+54+72+36 = 360°$. Le diagramme est correct. On trace ensuite chaque secteur au rapporteur, on colore, on légende et on donne un titre : « Répartition des 600 artisans de la coopérative par branche ».
\end{exoresolu}

\begin{popfigure}
\begin{center}
\begin{tikzpicture}
\begin{axis}[
  ybar, width=0.85\linewidth, height=6cm,
  bar width=14pt,
  ymin=0, ymax=220,
  ylabel={Nombre d'artisans},
  symbolic x coords={Vannerie,Sculpture,Poterie,Couture,Bijouterie},
  xtick=data, x tick label style={font=\footnotesize},
  nodes near coords, nodes near coords style={font=\footnotesize},
  ylabel style={font=\small},
]
\addplot[fill=poporange,draw=popdark] coordinates {(Vannerie,180) (Sculpture,150) (Poterie,90) (Couture,120) (Bijouterie,60)};
\end{axis}
\end{tikzpicture}
\end{center}
\legende{Diagramme en bâtons : répartition des 600 artisans de la coopérative par branche (données de l'exercice résolu). La vannerie domine (180 artisans, soit 30 \% de l'effectif).}
\end{popfigure}

\begin{attention}[Piège : choisir le mauvais diagramme]
N'utilisez \textbf{jamais} un diagramme circulaire pour comparer des valeurs qui ne forment pas un total (ex. : production de plusieurs années). À l'inverse, un diagramme en bâtons ne montre pas les \emph{parts} d'un ensemble. Question à se poser : « mes données forment-elles un tout (100 \%) ? » Si oui : circulaire. Si non : bâtons.
\end{attention}

\courssub{TP3 : Réaliser un croquis d'un espace artisanal}

\textbf{Intuition.} Le croquis est une \emph{photo dessinée et simplifiée} d'un espace réel : le village artisanal de Foumban, le marché artisanal de Yaoundé, un quartier d'ateliers de soudure à Douala. Contrairement à la carte, il montre l'\textbf{organisation interne} de l'espace : où sont les ateliers, où se vendent les produits, d'où viennent les matières premières.

\begin{methode}[Réussir un croquis en 4 étapes]
\begin{enumerate}
\item \textbf{Observer} l'espace (visite, photo, plan) et repérer les éléments essentiels.
\item \textbf{Sélectionner} : ne représenter que \textbf{3 à 5 éléments} (ateliers, zone de vente, stockage des matières premières, voie d'accès, point d'eau).
\item \textbf{Dessiner} des formes simples (rectangles, traits, points) à peu près à leur place réelle.
\item \textbf{Légender} avec des \textbf{flèches sobres} et numérotées, renvoyant à une légende placée en marge, plus un titre et le nord.
\end{enumerate}
\end{methode}

\begin{popfigure}
\begin{center}
\begin{tikzpicture}[scale=1.0]
% Cadre du village
\draw[thick,popdark,fill=popcream] (0,0) rectangle (8,5);
% Route
\draw[very thick,popmuted] (0,1.0) -- (8,1.4);
\node[popmuted,font=\footnotesize] at (6.9,1.05){route};
% Ateliers (zone ouest)
\fill[poporangeL] (0.4,2.0) rectangle (2.6,4.4);
\node[font=\footnotesize,align=center] at (1.5,3.2){Ateliers\\(sculpture,\\bronze)};
% Zone de vente (centre, près de la route)
\fill[popgreenL] (3.2,1.7) rectangle (5.2,3.0);
\node[font=\footnotesize,align=center] at (4.2,2.35){Boutiques\\et vente};
% Stockage matières premières (nord-est)
\fill[poppurpleL] (5.8,3.2) rectangle (7.6,4.6);
\node[font=\footnotesize,align=center] at (6.7,3.9){Stockage\\bois, cuir,\\argile};
% Habitations
\fill[popblueL] (3.2,3.5) rectangle (5.2,4.6);
\node[font=\footnotesize] at (4.2,4.05){Habitations};
% Nord
\draw[->,thick,popdark] (8.5,4.2) -- (8.5,4.9) node[above]{\footnotesize \textbf{N}};
% Légende fléchée
\node[font=\footnotesize,anchor=west] at (8.8,3.6){1. Ateliers (production)};
\draw[->,poporange,thick] (8.75,3.55) -- (2.65,3.3);
\node[font=\footnotesize,anchor=west] at (8.8,3.1){2. Zone de vente};
\draw[->,popgreen,thick] (8.75,3.05) -- (5.25,2.4);
\node[font=\footnotesize,anchor=west] at (8.8,2.6){3. Matières premières};
\draw[->,poppurple,thick] (8.75,2.55) -- (7.65,3.9);
\node[font=\footnotesize,anchor=west] at (8.8,2.1){4. Habitations};
\draw[->,popblue,thick] (8.75,2.05) -- (5.25,4.0);
\node[font=\footnotesize,anchor=west] at (8.8,1.6){5. Route (accès clients)};
\draw[->,popmuted,thick] (8.75,1.55) -- (7.0,1.35);
% Titre
\node at (4,5.35){\textbf{Croquis modèle : organisation d'un village artisanal}};
\end{tikzpicture}
\end{center}
\legende{Modèle de croquis légendé. La logique spatiale est visible : les ateliers (1) sont à l'écart, la vente (2) est près de la route (5) pour attirer la clientèle, le stockage des matières premières (3) est proche des ateliers. Seuls 5 éléments sont représentés : c'est la règle de l'essentiel.}
\end{popfigure}

\begin{attention}[Piège du croquis « fourre-tout »]
Un croquis surchargé (toits, arbres, personnages, motos...) cache l'essentiel et perd sa valeur géographique. Le correcteur attend une \textbf{sélection raisonnée} : chaque élément dessiné doit servir l'analyse de l'organisation de l'espace. Mieux vaut 4 éléments bien placés et bien légendés que 15 détails décoratifs.
\end{attention}

\courssub{Analyse croisée des productions des groupes}

Chaque groupe présente sa carte, son diagramme et son croquis au tableau. La classe compare les productions selon une grille commune :

\begin{center}
\begin{tabularx}{\linewidth}{|l|X|X|}\hline
\rowcolor{popblueL} \textbf{Critère} & \textbf{Question de contrôle} & \textbf{Points} \\ \hline
Titre & Précise-t-il le quoi, le où, le quand ? & /2 \\ \hline
Légende & Tous les figurés sont-ils expliqués ? & /3 \\ \hline
Orientation & Le nord est-il indiqué ? & /1 \\ \hline
Proportionnalité & Les tailles reflètent-elles les données ? & /2 \\ \hline
Lisibilité & 4 figurés au maximum, pas de surcharge ? & /2 \\ \hline
\end{tabularx}
\end{center}

La \textbf{correction collective} consiste à repérer les écarts entre les cartes : deux groupes partant des mêmes données peuvent produire des cartes très différentes selon le choix des figurés. On discute alors : quel choix rend le mieux visible la concentration de l'artisanat dans l'Ouest et les grandes villes ? C'est ainsi que se construit le raisonnement géographique.

\courssub{Complément utile au Bac : le commentaire de carte en trois temps}

\begin{methode}[Rédiger un commentaire de carte (méthode exigible)]
\begin{enumerate}
\item \textbf{Le constat} : décrire ce que montre la carte, du général au particulier (« l'artisanat camerounais est inégalement réparti : il se concentre dans l'Ouest, le Centre et le Littoral... »). Citer des faits et des chiffres, sans encore expliquer.
\item \textbf{L'explication} : rendre compte de la répartition par les facteurs étudiés (traditions culturelles, matières premières disponibles, présence de clients et de touristes, accès aux transports).
\item \textbf{Le bilan} : dégager l'essentiel et ouvrir (« cette concentration justifie les solutions proposées au TD2 : former et financer les artisans des régions délaissées »).
\end{enumerate}
\end{methode}

\begin{exercice}[Pour s'entraîner]
À partir de votre carte complétée des pôles artisanaux, rédigez un commentaire de 15 à 20 lignes en suivant le plan en trois temps (constat, explication, bilan). Vous appuierez votre constat sur au moins trois localisations précises et deux chiffres du tableau de données.
\end{exercice}

\begin{corrige}[Exercice — éléments de correction attendus]
\textbf{Constat} : la carte montre une répartition inégale ; les pôles majeurs sont Foumban (300 ateliers), Yaoundé et Douala (250 chacun) ; le Nord est moins doté (Maroua, 150 ; Bali, 100). \textbf{Explication} : l'Ouest doit sa place aux traditions royales (bronze de Foumban), les grandes villes à la demande urbaine et aux débouchés, le Nord au cuir lié à l'élevage. \textbf{Bilan} : l'artisanat est partout présent mais concentré ; d'où la nécessité de soutenir les pôles secondaires (formation, crédit, coopératives vues au TD2).
\end{corrige}

\begin{aretenir}[L'essentiel du TP]
\begin{itemize}
\item Une carte exige : \cle{titre}, \cle{légende}, \cle{nord}, figurés \cle{proportionnés} (4 maximum).
\item Parts d'un total $\rightarrow$ diagramme \cle{circulaire} (angle $= \frac{\text{effectif}}{\text{total}}\times 360°$) ; comparaison $\rightarrow$ diagramme en \cle{bâtons}.
\item Un croquis ne montre que \cle{l'essentiel} (3 à 5 éléments), légendé par des flèches sobres.
\item Le commentaire de carte se rédige en trois temps : \cle{constat}, \cle{explication}, \cle{bilan}.
\end{itemize}
\end{aretenir}

\begin{savaistu}[Le village artisanal de Foumban, vitrine du Cameroun]
Le quartier des artisans de Foumban, adossé au Palais des rois Bamoun, accueille sculpteurs, fondeurs de bronze et tisserands depuis des générations. Il attire chaque année des milliers de visiteurs et exporte vers toute l'Afrique centrale : c'est l'exemple parfait d'un espace artisanal où production, vente et patrimoine culturel se combinent — et un excellent sujet de croquis au Bac technique.
\end{savaistu}

\newpage

\coursec{Activité d'intégration}

\courssub{Objectifs de l'activité}

À l'issue de cette activité d'intégration, l'élève doit être capable de :
\begin{itemize}
\item mobiliser les notions de la leçon (définition, branches, importance et problèmes de l'artisanat) pour analyser une situation concrète de la vie courante ;
\item exploiter un dossier documentaire composé d'une photographie, d'un tableau de comptes simplifiés, d'une carte et de données touristiques ;
\item diagnostiquer et hiérarchiser les problèmes d'un atelier artisanal camerounais ;
\item proposer des solutions réalistes en identifiant les acteurs compétents (artisan, GIC, microfinance, MINPMEESA, commune) ;
\item rédiger et justifier une conclusion argumentée, puis la défendre devant la classe.
\end{itemize}

\courssub{Situation}

\textbf{Relancer l'atelier de sculpture de Foumban.}\par

Foumban, chef-lieu du département du Noun (région de l'Ouest), est la capitale historique de l'artisanat d'art camerounais. La ville abrite le palais royal des sultans Bamoun, un musée et le célèbre quartier des artisans, où travaillent sculpteurs sur bois, bronziers, potiers et tisserands. L'artisanat d'art y est la première activité non agricole et attire chaque année des touristes nationaux et étrangers.

Maître Ibrahim Njoya (nom d'emprunt), 52 ans, dirige un atelier de sculpture sur bois hérité de son père. L'atelier emploie 6 apprentis et produit des statues, des masques, des tabourets royaux et des objets de décoration. Depuis trois ans, l'atelier est en difficulté et Maître Ibrahim envisage de fermer. Le maire de Foumban, soucieux de préserver ce patrimoine, demande à ta classe de Première technique d'étudier le dossier et de proposer une stratégie de relance.

\textbf{Document 1 : photographie de l'atelier.} On observe un local vétuste, des outils manuels traditionnels (ciseaux à bois, gouges, herminettes) mais aucune machine électrique, des statues entassées sans emballage, et une devanture sans enseigne ni vitrine.

\textbf{Document 2 : comptes simplifiés de l'atelier (année écoulée).}

\begin{center}
\begin{tabularx}{\linewidth}{|l|X|}\hline
\rowcolor{popblueL} \textbf{Poste} & \textbf{Montant (FCFA)} \\ \hline
Recettes (ventes) & \num{2400000} \\ \hline
Dépenses : bois et matériaux & \num{900000} \\ \hline
Dépenses : rémunération des 6 apprentis & \num{1080000} \\ \hline
Dépenses : loyer, électricité, divers & \num{600000} \\ \hline
\textbf{Total des dépenses} & \textbf{\num{2580000}} \\ \hline
\end{tabularx}
\end{center}

Le prix moyen d'une statue est de \num{15000} FCFA ; l'atelier en a vendu 160 dans l'année, dont 110 à des touristes de passage et 50 à des revendeurs de Douala qui les revendent trois fois plus cher.

\textbf{Document 3 : carte des pôles artisanaux de l'Ouest.} Foumban (sculpture, bronze), Bafoussam (broderie, vannerie), Dschang (vannerie, poterie), Bandjoun et Baham (arts graphiques et sculpture). Foumban est reliée à Bafoussam (environ 70 km) puis à Douala par la route nationale n°5.

\textbf{Document 4 : tourisme dans la région.} Le palais royal et le musée de Foumban accueillent environ \num{20000} visiteurs par an ; la grande fête traditionnelle du Nguon attire à elle seule plusieurs milliers de visiteurs. Moins de 1 \% des visiteurs achètent actuellement un objet dans l'atelier de Maître Ibrahim.

\courssub{Consignes}

Par groupes de 4 à 6 élèves, tu répondras aux questions suivantes :

\begin{enumerate}
\item À partir du cours et du dossier, rappelle la définition de l'artisanat et précise à quelle branche de l'artisanat appartient l'atelier de Maître Ibrahim.
\item Calcule le résultat (bénéfice ou déficit) de l'atelier pour l'année écoulée. Que constates-tu ?
\item À partir des documents 1, 2 et 4, identifie les problèmes de l'atelier en les classant en trois catégories : problèmes financiers, problèmes techniques et problèmes commerciaux.
\item Hiérarchise ces problèmes : quel est, selon ton groupe, le problème le plus urgent ? Justifie ton choix par un chiffre du dossier.
\item Propose \textbf{trois solutions réalistes} pour relancer l'atelier. Pour chacune, précise l'acteur qui doit la mettre en œuvre (l'artisan lui-même, un GIC, une microfinance, le MINPMEESA, la commune de Foumban) et le problème qu'elle résout.
\item Rédige une conclusion argumentée d'une dizaine de lignes : quelle solution doit être mise en œuvre en priorité et pourquoi ? Appuie-toi sur au moins deux données chiffrées du dossier.
\item Présente la conclusion de ton groupe à la classe (3 minutes). La classe vote ensuite la stratégie jugée la plus pertinente.
\end{enumerate}

\courssub{Ressources à mobiliser}

\begin{aretenir}[Notions utiles de la leçon]
\begin{itemize}
\item L'\cle{artisanat} est une activité de production ou de services exercée par un artisan qui travaille essentiellement avec ses mains et des outils simples, souvent aidé d'apprentis, dans un petit atelier.
\item Branches : artisanat d'art (sculpture, bronze, poterie, vannerie), artisanat de production (menuiserie, couture, mécanique), artisanat de services (coiffure, réparation).
\item Importance : source d'emplois et de revenus, valorisation des matières premières locales, attraction touristique, préservation de la culture.
\item Problèmes : manque de capitaux et difficulté d'accès au crédit, outillage vétuste, absence de formation, difficultés d'écoulement des produits, concurrence des produits importés, absence d'organisation.
\item Solutions : création de GIC (Groupements d'Initiative Commune) et de coopératives, recours à la microfinance, appui de l'État (MINPMEESA : Ministère des Petites et Moyennes Entreprises, de l'Économie Sociale et de l'Artisanat), formation, foires et promotion touristique.
\end{itemize}
\end{aretenir}

\begin{methode}[Calculs utiles]
\begin{enumerate}
\item Résultat = recettes $-$ total des dépenses.
\item Part en pourcentage = (partie $\div$ total) $\times 100$.
\item Recette supplémentaire possible = nombre de clients visés $\times$ prix moyen.
\end{enumerate}
\end{methode}

\courssub{Démarche et corrigé commenté}

\begin{exoresolu}[Relancer l'atelier de sculpture de Foumban]
Énoncé : à partir du dossier, diagnostiquer les problèmes de l'atelier, proposer trois solutions avec les acteurs compétents et rédiger une conclusion justifiée.
\tcblower
\textbf{Étape 1 : branche de l'atelier.} La sculpture sur bois (statues, masques, tabourets royaux) relève de l'\emph{artisanat d'art}, branche qui produit des objets à valeur culturelle et touristique. Foumban est d'ailleurs le principal pôle national de cette branche.

\textbf{Étape 2 : résultat de l'atelier.}
\begin{align*}
\text{Résultat} &= \text{recettes} - \text{dépenses}\\
&= 2\,400\,000 - 2\,580\,000\\
&= -180\,000 \text{ FCFA}
\end{align*}
L'atelier est en \textbf{déficit de \num{180000} FCFA} : il dépense plus qu'il ne gagne. La fermeture menace donc réellement.

\textbf{Étape 3 : diagnostic des problèmes.}
\begin{itemize}
\item \emph{Financiers} : déficit de \num{180000} FCFA ; pas de capitaux pour investir ; pas d'accès au crédit bancaire classique.
\item \emph{Techniques} : outillage uniquement manuel (aucune machine), local vétuste, pas d'emballage des produits.
\item \emph{Commerciaux} : dépendance aux touristes de passage et aux revendeurs de Douala qui captent la marge (revente au triple) ; ni enseigne, ni vitrine ; moins de 1 \% des \num{20000} visiteurs annuels achètent.
\end{itemize}

\textbf{Étape 4 : hiérarchisation.} Le problème le plus urgent est \emph{commercial} : avec \num{20000} visiteurs par an, l'atelier ne vend que 110 statues aux touristes, soit 0,55 \% des visiteurs. Le marché existe mais n'est pas capté.

\textbf{Étape 5 : trois solutions.}
\begin{enumerate}
\item \emph{Créer un point de vente-vitrine près du palais royal et participer à la fête du Nguon} (acteurs : artisan + commune de Foumban). Si seulement 5 \% des \num{20000} visiteurs achetaient une statue à \num{15000} FCFA, la recette serait de $1\,000 \times 15\,000 = 15\,000\,000$ FCFA, très supérieure aux recettes actuelles.
\item \emph{Obtenir un crédit de microfinance via un GIC d'artisans} pour acheter des machines simples (ponceuse, scie électrique) et rénover l'atelier (acteurs : GIC + microfinance). Cela augmente la production et la qualité.
\item \emph{Solliciter le MINPMEESA} pour la formation des apprentis, l'enregistrement officiel de l'atelier et la participation aux foires artisanales nationales (acteur : MINPMEESA), afin de trouver des débouchés au-delà des revendeurs de Douala.
\end{enumerate}

\textbf{Étape 6 : conclusion argumentée (modèle).} « L'atelier de Maître Ibrahim est en déficit de \num{180000} FCFA, mais son problème principal n'est pas le manque de clients : c'est l'absence de débouchés organisés. En effet, Foumban reçoit environ \num{20000} visiteurs par an, dont plusieurs milliers lors de la fête du Nguon, et moins de 1 \% d'entre eux achètent à l'atelier, faute de vitrine et de signalisation. Nous proposons donc en priorité la création d'un point de vente-vitrine près du palais royal, avec l'appui de la commune : capter seulement 5 \% des visiteurs rapporterait environ 15 millions de FCFA, soit plus de six fois les recettes actuelles. Cette solution, peu coûteuse et rapide, dégagera ensuite les ressources nécessaires pour moderniser l'outillage grâce à un crédit de microfinance obtenu par le GIC des artisans. Ainsi, l'artisanat d'art de Foumban continuera de créer des emplois, de faire vivre la culture bamoun et d'attirer les touristes. »

\textbf{Étape 7 : présentation et vote.} Chaque groupe présente sa stratégie ; la classe compare les arguments (données chiffrées, réalisme, acteurs identifiés) et vote la stratégie la plus pertinente.
\end{exoresolu}

\courssub{Bilan de l'activité}

Cette activité a permis de mettre en évidence que :
\begin{itemize}
\item l'artisanat d'art, bien que porteur d'emplois (6 apprentis ici) et de valeur culturelle, reste fragile lorsqu'il manque de capitaux, d'outillage moderne et surtout de débouchés commerciaux ;
\item un diagnostic rigoureux s'appuie sur des \cle{données chiffrées} (déficit de \num{180000} FCFA, taux de captation des visiteurs inférieur à 1 \%) et non sur des impressions ;
\item la relance de l'artisanat exige la \cle{concertation de plusieurs acteurs} : l'artisan lui-même, regroupé en GIC, les établissements de microfinance, l'État à travers le MINPMEESA, et les collectivités locales comme la commune de Foumban ;
\item rédiger une conclusion justifiée suppose de hiérarchiser les problèmes, de choisir une solution prioritaire et de la défendre par des arguments chiffrés — compétence directement mobilisable dans la vie courante et à l'examen.
\end{itemize}

\newpage

\coursec{Méthodes et exercices résolus}

\coursec{Les activités artisanales}

\courssub{Définition et caractéristiques de l'artisanat}

\begin{methode}[Analyser une définition et caractériser l'artisanat]
\textbf{Objectif :} Savoir définir l'artisanat et en dégager les caractéristiques principales à partir d'un document (texte, image, témoignage).
\begin{enumerate}
    \item \textbf{Identifier l'activité :} Repérer le secteur (transformation de matières premières locales), le statut de l'acteur (artisan) et l'échelle (unité de production individuelle ou familiale).
    \item \textbf{Relever les critères distinctifs :} 
    \begin{itemize}
        \item Main-d'œuvre \cle{prépondérante manuelle} (outils simples, machines légères).
        \item Matières premières \cle{locales} (bois, argile, cuir, fibres, métaux de récupération).
        \item \cle{Savoir-faire traditionnel} transmis oralement ou par imitation (maître/apprenti).
        \item Production \cle{unique ou en petite série}, à forte valeur identitaire/culturelle.
        \item Structure \cle{informelle} prédominante (peu capitalisée, gestion familiale).
    \end{itemize}
    \item \textbf{Différencier :} Artisanat d'art (esthétique), artisanat de service (réparation, coiffure), artisanat de production (objets utilitaires).
    \item \textbf{Rédiger la définition synthétique :} « L'artisanat est une activité de transformation de matières premières locales, utilisant un savoir-faire traditionnel et une main-d'œuvre manuelle prépondérante, pour produire des biens uniques ou en petite série, souvent dans le secteur informel. »
\end{enumerate}
\end{methode}

\begin{exoresolu}[Caractérisation de l'artisanat à partir d'un témoignage]
\textbf{Document :} Extrait d'un entretien avec M. Njoya, sculpteur sur bois à Foumban. 
« J'ai appris le métier auprès de mon père dès l'âge de 12 ans. Je travaille dans mon atelier attenant à la case familiale. J'utilise du bois d'iroko et de l'ébène que j'achète aux exploitants forestiers locaux. Mes outils sont des ciseaux à bois, des gouges et une vieille perceuse à colonne. Je fabrique des masques, des statuettes et des chaises royales sur commande. Chaque pièce est unique. Je forme actuellement trois apprentis. Je vends aux touristes de passage et aux grossistes de Douala, mais je n'ai pas de compte en banque, tout se fait en espèces. »

\textbf{Question :} En vous appuyant sur le document, citez et expliquez quatre caractéristiques de l'artisanat illustrées par le cas de M. Njoya.

\tcblower
\textbf{Solution détaillée :}

\textbf{1. Transmission du savoir-faire par apprentissage traditionnel (maître/apprenti).}
\textit{Justification :} Le texte précise : « J'ai appris le métier auprès de mon père dès l'âge de 12 ans » et « Je forme actuellement trois apprentis ». Cela illustre la transmission orale et gestuelle, hors système scolaire formel, caractéristique majeure de l'artisanat.

\textbf{2. Utilisation de matières premières locales.}
\textit{Justification :} M. Njoya utilise « du bois d'iroko et de l'ébène que j'achète aux exploitants forestiers locaux ». L'approvisionnement local (forêt équatoriale proche) réduit les coûts et ancre l'activité dans le terroir.

\textbf{3. Main-d'œuvre manuelle prépondérante et outillage léger.}
\textit{Justification :} Les outils cités sont des « ciseaux à bois, des gouges et une vieille perceuse à colonne ». L'absence de machines à commande numérique ou de chaîne de montage automatisée confirme le caractère manuel. La perceuse à colonne est une machine légère, seule concession à la modernité.

\textbf{4. Production unique, à forte valeur identitaire, structure informelle.}
\textit{Justification :} « Chaque pièce est unique » (pas de standardisation industrielle). Les objets (masques, statuettes, chaises royales) portent une charge culturelle forte (art Bamoun). La gestion « en espèces », sans « compte en banque », et l'atelier « attenant à la case familiale » confirment le caractère \cle{informel} et \cle{familial} de l'unité de production.

\textbf{Conclusion :} Le cas de M. Njoya est archétypal de l'\cle{artisanat d'art} camerounais : ancrage culturel, transmission intergénérationnelle, ressources locales, informel.
\end{exoresolu}

\courssub{Les branches et la répartition spatiale de l'artisanat camerounais}

\begin{methode}[Réaliser un croquis de répartition spatiale de l'artisanat au Cameroun]
\textbf{Objectif :} Construire un croquis légendé montrant la spécialisation régionale de l'artisanat (TP Cartographie).
\begin{enumerate}
    \item \textbf{Choisir le fond de carte :} Contour simplifié du Cameroun (forme triangle rectangle), grands axes (RN1, RN3, RN6, Transcamerounais), fleuves (Bénoué, Sanaga, Wouri), 6-7 villes repères (Yaoundé, Douala, Garoua, Maroua, Bamenda, Foumban, Bafoussam).
    \item \textbf{Structurer la légende (3 à 4 items max) :} Grouper par \cle{spécialisation dominante} par région naturelle/culturelle.
    \begin{itemize}
        \item \textbf{Zone Soudano-sahélienne (Nord) :} Cuir (Maroua, Garoua), Potterie (Mora, Rhumsiki), Tissage (Mandara), Métallurgie traditionnelle (Guider).
        \item \textbf{Zone des Hauts Plateaux (Ouest/Nord-Ouest) :} Sculpture sur bois (Foumban, Bafoussam, Baham), Potterie (Baham, Bandjoun), Perles/Bijouterie (Bafoussam), Tissage Ndop (Bamenda).
        \item \textbf{Zone Forestière (Sud, Est, Centre) :} Sculpture (Ebolowa, Sangmélima), Vannerie/Objets en raphia, Pirogues (Littoral/Sud-Ouest).
        \item \textbf{Grandes métropoles (Yaoundé, Douala) :} Artisanat de service (mécanique, couture, coiffure, ferronnerie), Artisanat d'art de revente (marchés artisanaux : Mfoundi, Akwa).
    \end{itemize}
    \item \textbf{Choisir la sémiologie :} 
    \begin{itemize}
        \item \cle{Signes surfaciques} (hachures, pointillés, couleurs pastel) pour les zones de production brute.
        \item \cle{Signes ponctuels} (symboles distincts : ciseau, pot, peau, métier à tisser) pour les pôles spécialisés (villes).
        \item \cle{Signes linéaires} (flèches épaisses) pour les flux de commercialisation vers Douala/Yaoundé/Export.
    \end{itemize}
    \item \textbf{Finaliser :} Titre précis, légende organisée, échelle indicative, flèche Nord, source.
\end{enumerate}
\end{methode}

\begin{exoresolu}[Croquis : Spécialisation artisanale régionale au Cameroun]
\textbf{Consigne :} À l'aide d'un fond de carte du Cameroun fourni, réalisez un croquis intitulé « Les grandes régions de spécialisation artisanale au Cameroun ». La légende comportera trois items pour les zones de production et un item pour les flux commerciaux.

\tcblower
\textbf{Solution détaillée :}

\textbf{1. Analyse de la demande :} 3 items zones + 1 item flux = 4 items au total. Titre imposé.

\textbf{2. Proposition de légende structurée :}
\begin{center}
\begin{tabularx}{\linewidth}{|l|X|}\hline
\textbf{Symbole} & \textbf{Signification} \\ \hline
\cellcolor{poporangeL}\rule{1cm}{4pt} & \textbf{Zone soudano-sahélienne :} Cuir (Maroua, Garoua), Potterie, Tissage, Métallurgie \\ \hline
\cellcolor{popgreenL}\rule{1cm}{4pt} & \textbf{Zone des Hauts Plateaux (Ouest/Nord-Ouest) :} Sculpture bois (Foumban, Baham), Potterie, Perles, Tissage Ndop \\ \hline
\cellcolor{popblueL}\rule{1cm}{4pt} & \textbf{Zone forestière (Sud/Centre/Est) :} Sculpture, Vannerie/Raphia, Pirogues, Ébénisterie \\ \hline
$\blacktriangleright$ \textcolor{popdark}{Epaisseur $\propto$ volume} & \textbf{Flux de commercialisation} vers Douala, Yaoundé et ports d'exportation \\ \hline
\end{tabularx}
\end{center}

\textbf{3. Code TikZ du croquis attendu (schéma simplifié pour l'élève) :}
\begin{popfigure}
\begin{tikzpicture}[scale=0.85, every node/.style={font=\footnotesize}]
    % Contour Cameroun simplifié (polygone approximatif)
    \draw[popdark, thick] (0,0) -- (8,0) -- (9,5) -- (7,9) -- (3,10) -- (0,6) -- cycle;
    % Villes repères
    \node[circle, fill=popdark, inner sep=1.5pt, label=below:\textbf{Maroua}] (Mar) at (1.5, 1.2) {};
    \node[circle, fill=popdark, inner sep=1.5pt, label=below:\textbf{Garoua}] (Gar) at (3.5, 2.5) {};
    \node[circle, fill=popdark, inner sep=1.5pt, label=left:\textbf{Foumban}] (Fou) at (3.5, 6.5) {};
    \node[circle, fill=popdark, inner sep=1.5pt, label=left:\textbf{Bafoussam}] (Baf) at (2.8, 7.2) {};
    \node[circle, fill=popdark, inner sep=1.5pt, label=above:\textbf{Bamenda}] (Bam) at (4.5, 8.5) {};
    \node[circle, fill=popdark, inner sep=1.5pt, label=right:\textbf{Yaoundé}] (Yao) at (5.5, 6.0) {};
    \node[circle, fill=popdark, inner sep=1.5pt, label=right:\textbf{Douala}] (Dou) at (7.5, 3.5) {};
    % Zones surfaciques (hachures via patterns)
    \pattern[pattern=north east lines, pattern color=poporange!60] (0,0) -- (8,0) -- (5,3) -- (1,3) -- cycle; % Nord
    \pattern[pattern=north west lines, pattern color=popgreen!60] (2,5) -- (5,5) -- (5,9) -- (2,9.5) -- cycle; % Ouest
    \pattern[pattern=dots, pattern color=popblue!60] (5,3) -- (9,5) -- (7,9) -- (5,5) -- cycle; % Sud/Est
    % Flux (flèches vers Douala/Yaoundé)
    \draw[popdark, very thick, -{Stealth[length=3mm]}] (Mar) -- (Dou);
    \draw[popdark, very thick, -{Stealth[length=3mm]}] (Gar) -- (Dou);
    \draw[popdark, thick, -{Stealth[length=2mm]}] (Fou) -- (Dou);
    \draw[popdark, thick, -{Stealth[length=2mm]}] (Bam) -- (Yao);
    \draw[popdark, thick, -{Stealth[length=2mm]}] (4,8) -- (Yao);
    % Flèche Nord
    \node[align=center] at (8.5, 9.5){\textbf{N}\\\uparrow};
    % Échelle
    \node at (0.5, -0.5) {Échelle approx. 1/10 000 000};
\end{tikzpicture}
\legende{Croquis schématique : Les grandes régions de spécialisation artisanale au Cameroun. Les hachures indiquent les zones de production dominante ; les flèches épaisses les flux principaux vers les métropoles et l'export.}
\end{popfigure}

\textbf{4. Justification des choix :}
\begin{itemize}
    \item Le Nord (orange) : Zone sahélienne/soudanienne, élevage $\rightarrow$ cuir, argile $\rightarrow$ poterie, coton $\rightarrow$ tissage, fer $\rightarrow$ métallurgie.
    \item L'Ouest (vert) : Forêt galerie/savane boisée, bois durs $\rightarrow$ sculpture, chefferies $\rightarrow$ art de cour (perles, Ndop), argile $\rightarrow$ poterie.
    \item Le Sud/Est (bleu) : Forêt dense $\rightarrow$ bois précieux (ébène), raphia, voies navigables $\rightarrow$ pirogues.
    \item Flux : Convergence vers Douala (port export) et Yaoundé (admin/tourisme).
\end{itemize}
\end{exoresolu}

\courssub{L'importance de l'artisanat au Cameroun}

\begin{methode}[Démontrer l'importance socio-économique de l'artisanat (ESS)]
\textbf{Objectif :} Structurer une argumentation (rédaction type Probatoire/Baccalauréat) autour des 4 piliers : Emploi, Revenu/PIB, Culture/Identité, Tourisme/Export.
\begin{enumerate}
    \item \textbf{Introduction :} Définition rapide + Accroche chiffrée (ex: « 2ème pourvoyeur d'emplois après l'agriculture », « $\approx$ 20\% population active »). Annonce du plan (3 axes).
    \item \textbf{Axe 1 : Pourvoyeur d'emplois et de revenus (Économie réelle).}
    \begin{itemize}
        \item Absorption main-d'œuvre non qualifiée / jeunes / femmes (secteur informel).
        \item Revenus quotidiens pour ménages vulnérables (seuil de pauvreté).
        \item Formation par apprentissage (insertion pro sans diplôme).
    \end{itemize}
    \item \textbf{Axe 2 : Valorisation du patrimoine et levier touristique.}
    \begin{itemize}
        \item Objets = supports culturels (rites, chefferies, musées).
        \item Attractivité touristique (marchés artisanaux, routes des chefferies, FESTAC).
        \item Devises : Export objets d'art (Europe, USA, Asie).
    \end{itemize}
    \item \textbf{Axe 3 : Dynamique locale et maillage territorial.}
    \begin{itemize}
        \item Fixation populations en milieu rural (lutte exode).
        \item Valorisation ressources locales (bois, argile, cuir) $\rightarrow$ valeur ajoutée locale.
        \item Essor PME/Coopératives (ex: Groupements d'artisans de Foumban, Baham).
    \end{itemize}
    \item \textbf{Conclusion :} Synthèse + Ouverture (nécessité structuration/formalisation pour passer à l'échelle).
\end{enumerate}
\cle{Formule type Examen :} « L'artisanat est un \textbf{amortisseur social} majeur et un \textbf{vecteur d'identité nationale}, mais sa contribution au PIB formel reste sous-estimée par l'informel. »
\end{methode}

\begin{exoresolu}[L'artisanat, pilier de l'économie locale : rédaction guidée]
\textbf{Sujet :} « L'artisanat camerounais est un secteur clé pour le développement local. » En vous appuyant sur des exemples précis, montrez la validité de cette affirmation en structurant votre réponse en trois parties.

\tcblower
\textbf{Solution détaillée (Plan rédigé type Examen - 15/20 visé) :}

\textbf{Introduction}
[Accroche] Au Cameroun, le secteur informel emploie près de 90\% de la population active (INS, EESI 2021). L'artisanat en constitue l'épine dorsale productive.
[Définition] Rappel : activité de transformation manuelle de ressources locales par savoir-faire traditionnel.
[Problématique] Comment un secteur faiblement capitalisé devient-il un levier de développement territorial ?
[Annonce] Nous analyserons tour à tour son rôle d'amortisseur social (I), de gardien de l'identité culturelle et levier touristique (II), puis de dynamiseur des territoires ruraux (III).

\textbf{I. Un amortisseur social majeur : emploi et revenus pour les exclus du formel}
L'artisanat absorbe une main-d'œuvre jeune, peu scolarisée et féminine, rejetée par le marché salarial formel.
\begin{itemize}
    \item \textbf{Chiffres :} Environ 2,5 à 3 millions d'artisans (estimation MINPMEESA), soit $\approx$ 15-20\% pop. active.
    \item \textbf{Exemple :} À Douala, le quartier \textit{New Bell} ou \textit{Makepe} concentre des milliers de soudeurs, menuisiers, mécaniciens « de rue » qui forment des apprentis (coût formation quasi nul pour l'État).
    \item \textbf{Impact :} Revenu journalier permettant la scolarisation des enfants, l'alimentation, la santé. C'est une \cle{économie de survie} qui prévient l'instabilité sociale.
\end{itemize}

\textbf{II. Patrimoine culturel et levier touristique : la « marque Cameroun »}
L'artisanat d'art transforme la culture en valeur marchande.
\begin{itemize}
    \item \textbf{Identité :} Le \cle{Ndop} (Bamenda), le \cle{Toghu} (Nord-Ouest), les \cle{masques Bamoun} (Foumban), les \cle{pipes en terre} (Baham) sont des marqueurs ethniques.
    \item \textbf{Tourisme :} Le \cle{Marché des Artisans de Foumban} (réhabilité 2020) attire 50 000 visiteurs/an. La \cle{Route des Chefferies} (Ouest) génère des retombées directes (vente objets, hébergement, restauration).
    \item \textbf{Export :} Les statues « colon » (Bamileke), masques Fang (Sud) s'exportent. Devises non négligeables mais mal tracées (informel).
\end{itemize}

\textbf{III. Ancrage territorial et valorisation des ressources locales}
L'artisanat fixe la population là où la ressource est.
\begin{itemize}
    \item \textbf{Ruralité :} Potiers de \cle{Baham} (argile locale), Tanneurs de \cle{Maroua} (peaux bovins locale), Vanniers du \cle{Sanaga} (raphia).
    \item \textbf{Valeur ajoutée :} Transformation sur place du brut (bois rond $\rightarrow$ statue ; peau brute $\rightarrow$ sac). Évite l'export de matière première brute.
    \item \textbf{Structuration naissante :} Création de \cle{coopératives} (ex: COOPAF à Foumban) et \cle{centres de formation} (CFP de Mbalmayo, Institut de Formation Artisanale de Foumban) soutenus par le MINPMEESA et partenaires (UE, GIZ).
\end{itemize}

\textbf{Conclusion}
[Bilan] L'artisanat est bien un pilier : social (emploi), culturel (identité), spatial (maillage rural), économique (revenus, devises).
[Ouverture] Cependant, sa contribution au PIB formel reste faible ($< 5\%$ estimé). Le défi est la \cle{formalisation} (statut d'auto-entrepreneur, accès crédit, protection propriété intellectuelle) pour passer de l'économie de survie à l'économie de croissance.
\end{exoresolu}

\courssub{Les problèmes de l'artisanat camerounais}

\begin{methode}[Analyser les problèmes de l'artisanat (Grille d'analyse SWOT simplifiée)]
\textbf{Objectif :} Classer les contraintes en 4 catégories pour une réponse exhaustive (TD2 préparation).
\begin{enumerate}
    \item \textbf{Approvisionnement (Amont) :} Pénurie/coût matières premières (bois légal rare, argile surexploitée), absence de circuits d'approvisionnement organisés, concurrence industrie (cuir, textile).
    \item \textbf{Production (Processus) :} Outillage vétuste, pénibilité, faible productivité, manque formation technique moderne (gestion, design, normes), transmission orale fragile (perte gestes).
    \item \textbf{Commercialisation (Aval) :} Circuits courts inexistants, dépendance intermédiaires (revendeurs, tour-opérateurs), marchés locaux saturés, faible accès export (normes qualité, emballage, logistique), contrefaçon (imports asiatiques bon marché).
    \item \textbf{Environnement institutionnel/Financier :} Statut informel (pas de crédit bancaire, pas assurance, pas protection sociale), fiscalité harcelante (contrôles multiples), absence de propriété intellectuelle (brevets dessins/modèles), locaux précaires (expulsions).
\end{enumerate}
\cle{Astuce Examen :} Citer l'exemple précis : « L'artisan de Foumban achète l'iroko 3x plus cher qu'il y a 10 ans car les scieries exportent les grumes » (Amont). « Il vend son masque 50 000 FCFA au revendeur qui le revend 500 000 FCFA à Paris » (Aval).
\end{methode}

\begin{exoresolu}[Étude de cas : Les problèmes des potiers de Baham (Ouest)]
\textbf{Document :} Synthèse d'une enquête de terrain (2023) chez le groupement « Femmes Potières de Baham ».
\begin{itemize}
    \item \textbf{Ressource :} L'argile de qualité se trouve à 8 km (site de \textit{Nkouong}). Transport à dos d'homme ou moto-taxi (coût : 15 000 FCFA/voyage). Le site est menacé par l'urbanisation (lotissements).
    \item \textbf{Production :} Cuisson à ciel ouvert (feu de bois). Taux de casse $\approx$ 40\% (chocs thermiques). Pas de four amélioré. Pénibilité : port de charges lourdes, fumées.
    \item \textbf{Produit :} Marmites, canari, pots à eau. Design inchangé depuis 30 ans. Concurrence des marmites en aluminium (importées, 2 000 FCFA vs poterie 3 500 FCFA) et plastique.
    \item \textbf{Vente :} Marché hebdomadaire de Baham (samedi). Clients : ménagères locales + quelques revendeurs Bafoussam. Prix imposés par acheteurs. Pas de label, pas d'emballage.
    \item \textbf{Statut :} Groupement informel (reconnu mairie seulement). Pas de compte bancaire groupé. Crédit refusé par microfinance (pas de garanties).
\end{itemize}

\textbf{Question :} En utilisant la grille d'analyse (Amont, Processus, Aval, Institutionnel), dressez le diagnostic complet des contraintes du groupement.

\tcblower
\textbf{Solution détaillée :}

\begin{center}
\begin{tabularx}{\linewidth}{|l|X|}\hline
\textbf{Dimension} & \textbf{Contraintes identifiées (Preuves du doc)} \\ \hline
\textbf{AMONT (Approvisionnement)} & \begin{itemize}\item \textbf{Pénurie spatiale :} Gisement unique à 8 km (éloignement). \item \textbf{Coût logistique élevé :} 15 000 FCFA/voyage (moto/portage) $\rightarrow$ renchérissement coût de revient. \item \textbf{Menace foncière :} Urbanisation (lotissements) grignote le gisement $\rightarrow$ risque rupture définitive.\end{itemize} \\ \hline
\textbf{PROCESSUS (Production)} & \begin{itemize}\item \textbf{Technologie archaïque :} Cuisson à ciel ouvert (bois) $\rightarrow$ \textbf{Taux casse 40\%} (perte sèche massive). \item \textbf{Pénibilité/Santé :} Port charges, fumées toxiques (maladies respiratoires). \item \textbf{Innovation nulle :} Design figé (marmites traditionnelles) $\rightarrow$ inadaptation besoins modernes (cuisinières gaz/électrique).\end{itemize} \\ \hline
\textbf{AVAL (Commercialisation)} & \begin{itemize}\item \textbf{Marché captif/étroit :} Marché hebdo local uniquement. \item \textbf{Puissance acheteurs :} Prix imposés par revendeurs (asymétrie info). \item \textbf{Concurrence déloyale :} Aluminium/Plastique moins chers, plus légers, adaptés cuisson moderne. \item \textbf{Absence marketing :} Pas d'emballage, pas de label « Baham », pas de vente en ligne/boutique fixe.\end{itemize} \\ \hline
\textbf{INSTITUTIONNEL / FINANCIER} & \begin{itemize}\item \textbf{Informel :} Groupement non coopérative (statut associatif simple) $\rightarrow$ pas de personnalité morale forte. \item \textbf{Exclusion bancaire :} Pas de compte, pas de garanties $\rightarrow$ crédit refusé. \item \textbf{Pas de protection :} Pas de marque collective, pas de brevet design $\rightarrow$ copie possible.\end{itemize} \\ \hline
\end{tabularx}
\end{center}

\textbf{Synthèse diagnostique :} Le groupement est piégé dans un \cle{cercle vicieux} : Faible productivité (processus) $\rightarrow$ Coût élevé + Qualité faible $\rightarrow$ Prix non compétitif (Aval) $\rightarrow$ Revenu faible $\rightarrow$ Pas d'investissement (Four, camion, formation) $\rightarrow$ Maintien faible productivité. L'insécurité foncière (Amont) menace l'existence même de la ressource.
\end{exoresolu}

\courssub{TD2 — Les solutions aux problèmes de l'artisanat}

\begin{methode}[Proposer des solutions structurées (TD2 : Réponses aux problèmes)]
\textbf{Objectif :} Faire correspondre \textbf{chaque problème} identifié à une \textbf{solution concrète, acteur responsable, et indicateur de réussite}. Structure type « Fiche action ».
\begin{enumerate}
    \item \textbf{Problème Amont (Ressources) :} 
    \begin{itemize}
        \item \textit{Solution :} Création de \cle{plantations artisanales} (bois à croissance rapide : Gmelina, Acacia) / \cle{Zonage foncier} protégeant gisements argile / Groupements d'achat mutualisés.
        \item \textit{Acteurs :} MINFOF, MINDAF, Communes, Groupements.
    \end{itemize}
    \item \textbf{Problème Processus (Production) :} 
    \begin{itemize}
        \item \textit{Solution :} Diffusion \cle{fours améliorés} (type \textit{Anagama} ou fours à gaz/électrique pour poterie), \cle{formation design/innovation} (centres CFP, designers), \cle{ergonomie} postes de travail.
        \item \textit{Acteurs :} MINPMEESA, MINEFOP, ONG (GIZ, SNV), Écoles de design (ESSTIC, IBA).
    \end{itemize}
    \item \textbf{Problème Aval (Marché) :} 
    \begin{itemize}
        \item \textit{Solution :} \cle{Marque collective} / Label « Made in Cameroon » / \cle{Boutiques partagées} en ville / \cle{E-commerce} (Jumia, plateforme dédiée) / \cle{Salons professionnels} (PROMOTE, SIAO Ouaga) / \cle{Normes qualité} (ANOR).
        \item \textit{Acteurs :} MINCOMMERCE, ANOR, API, Chambres de Métiers.
    \end{itemize}
    \item \textbf{Problème Institutionnel/Financier :} 
    \begin{itemize}
        \item \textit{Solution :} \cle{Formalisation} (Statut Auto-entrepreneur 2023, Carte d'artisan), \cle{Fonds de garantie} / \cle{Crédit adapté} (FNE, PIAASI, Microfinance conventionnée), \cle{Protection PI} (OAPI - dessins/modèles), \cle{Assurance maladie} (CNPS secteur informel).
        \item \textit{Acteurs :} MINPMEESA, MINFI, CNPS, OAPI, Banques.
    \end{itemize}
\end{enumerate}
\cle{Réflexe Examen :} Toujours lier \textbf{Problème $\rightarrow$ Solution $\rightarrow$ Acteur $\rightarrow$ Effet attendu}.
\end{methode}

\begin{exoresolu}[Fiche Action : Solutions pour les Potiers de Baham]
\textbf{Consigne :} À partir du diagnostic de l'exercice précédent, remplissez le tableau de planification d'actions prioritaires (3 actions max) pour le groupement « Femmes Potières de Baham ».

\tcblower
\textbf{Solution détaillée :}

\begin{center}
\begin{tabularx}{\linewidth}{|l|X|X|X|X|}\hline
\textbf{N°} & \textbf{Action Prioritaire (Solution)} & \textbf{Problème Cible} & \textbf{Acteurs Principaux} & \textbf{Indicateur de Résultat (à 2 ans)} \\ \hline
1 & \textbf{Construction d'un four amélioré collectif (type four à gaz ou Anagama isolé) + Formation cuisson.} & Taux casse 40\% (Processus) ; Pénibilité/Bois (Environnement) & MINPMEESA (financement), CFP Bafoussam (formation), Groupement (gestion) & Taux casse $< 10\%$ ; Conso bois $-50\%$ ; Production $+30\%$. \\ \hline
2 & \textbf{Sécurisation foncière du gisement d'argile (Nkouong) + Achat tricycle motorisé transport.} & Éloignement 8 km + Coût transport + Menace lotissement (Amont) & Commune Baham (délibération protection), MINDAF (titre foncier), Projet d'appui (don tricycle) & Gisement classé « Zone artisanale » ; Coût transport $-70\%$ ; Stock argile sécurisé 10 ans. \\ \hline
3 & \textbf{Création Marque Collective « Poterie de Baham » + Boutique showroom Bafoussam + Formation design (cocottes, tajines, déco).} & Design obsolète + Concurrence Alu/Plastique + Prix imposés (Aval) & OAPI (dépôt marque), ANOR (normes), Designer (création), MINCOMMERCE (boutique) & Marque déposée ; 15 nouveaux modèles ; CA $+50\%$ ; Vente directe 60\% vs revendeurs. \\ \hline
\end{tabularx}
\end{center}

\textbf{Justification pédagogique :}
\begin{itemize}
    \item \textbf{Action 1} : Impact immédiat sur le revenu (moins de casse = plus de stock à vendre) et santé. Coût modéré, tech maîtrisée.
    \item \textbf{Action 2} : Condition \textit{sine qua non} de la survie de l'activité (ressource). Implique la puissance publique (Commune/MINDAF).
    \item \textbf{Action 3} : Rupture stratégique. Sortie de la concurrence par le prix (perdue d'avance) vers la concurrence par la \cle{différenciation} (design, origine, qualité). Nécessite appui technique pointu (Design, Juridique OAPI).
\end{itemize}
\end{exoresolu}

\courssub{TP — Cartographie et représentation graphique de l'artisanat}

\begin{methode}[Traiter des données statistiques et réaliser une représentation graphique (TP Cartographie/Graphiques)]
\textbf{Objectif :} Passer d'un tableau brut à un graphique choisi (barres, secteurs, courbe) + commentaire géographe.
\begin{enumerate}
    \item \textbf{Lire le tableau :} Identifier variables (années, secteurs, effectifs, CA), unités (milliers, FCFA, \%), source.
    \item \textbf{Calculer si besoin :} Parts relatives (\%), Taux de variation annuel moyen (TVAM), Rang.
    \item \textbf{Choisir le graphique adapté :}
    \begin{itemize}
        \item \textbf{Diagramme en barres :} Comparer effectifs/CA par secteur \textbf{à une date} (ex: 2022).
        \item \textbf{Diagramme circulaire (camembert) :} Montrer la \textbf{structure} (\% de chaque secteur dans le total) \textbf{à une date}.
        \item \textbf{Courbe d'évolution :} Montrer la \textbf{dynamique temporelle} d'un agrégat (ex: Effectif total 2010-2022).
    \end{itemize}
    \item \textbf{Construire rigoureusement :} Titre complet (Quoi ? Où ? Quand ?), Axes légendés (grandeur + unité), Échelle adaptée, Source, Légende des séries.
    \item \textbf{Commenter (3 temps) :} 
    \begin{enumerate}
        \item \textbf{Lecture globale :} « Le graphique montre la répartition de l'artisanat par branche en 2022... »
        \item \textbf{Analyse chiffrée :} « Le secteur dominant est la menuiserie (34,3\%, 115 000 artisans), suivi de la couture (26,9\%). La métallurgie ne pèse que 17,3\%. »
        \item \textbf{Interprétation géographique :} « La dominance du bois s'explique par la ressource forestière (Sud/Est/Ouest) et l'habitat. La couture répond à un besoin vestimentaire universel. La faiblesse relative de la métallurgie tient au coût des intrants (fer importé). »
    \end{enumerate}
\end{enumerate}
\end{methode}

\begin{exoresolu}[Statistiques et graphique : Évolution de l'emploi artisanal par secteur (2015-2022)]
\textbf{Données :} Effectifs estimés (en milliers d'artisans) - Source : MINPMEESA / INS (Projections).
\begin{center}
\begin{tabular}{|c|c|c|c|c|c|}\hline
\textbf{Secteur} & \textbf{2015} & \textbf{2018} & \textbf{2020} & \textbf{2022} & \textbf{Part 2022 (\%)} \\ \hline
Menuiserie / Ébénisterie & 85 & 95 & 105 & 115 & 34,3 \% \\ \hline
Couture / Habillement & 60 & 70 & 80 & 90 & 26,9 \% \\ \hline
Mécanique / Soudure & 35 & 42 & 50 & 58 & 17,3 \% \\ \hline
Alimentation (Boulangerie, boissons) & 25 & 30 & 35 & 40 & 11,9 \% \\ \hline
Art d'art / Sculpture / Potterie & 18 & 20 & 22 & 24 & 7,2 \% \\ \hline
Autres (Coiffure, Esthétique, etc.) & 5 & 7 & 8 & 8 & 2,4 \% \\ \hline
\textbf{TOTAL} & \textbf{228} & \textbf{264} & \textbf{300} & \textbf{335} & \textbf{100 \%} \\ \hline
\end{tabular}
\end{center}

\textbf{Consignes :}
\begin{enumerate}
    \item Calculer le Taux de Variation Annuel Moyen (TVAM) de l'effectif total entre 2015 et 2022.
    \item Construire un \textbf{diagramme en barres groupées} comparant les effectifs 2015 et 2022 par secteur.
    \item Commenter l'évolution de la structure sectorielle.
\end{enumerate}

\tcblower
\textbf{Solution détaillée :}

\textbf{1. Calcul du TVAM (Effectif Total)}
Formule : $TVAM = \left( \frac{V_{fin}}{V_{deb}} \right)^{\frac{1}{n}} - 1$ avec $n = 2022 - 2015 = 7$ ans.
$V_{deb} = 228$ ; $V_{fin} = 335$.
$\frac{335}{228} \approx 1,4693$.
$1,4693^{\frac{1}{7}} \approx 1,0564$ (calculatrice : $1,4693^{(1/7)}$).
$TVAM \approx 0,0564 = \textbf{5,6 \% par an}$.
\textit{Interprétation :} L'emploi artisanal croît rapidement (+5,6\%/an), bien plus que la croissance démographique ($\approx$ 2,6\%), confirmant son rôle d'amortisseur.

\textbf{2. Diagramme en barres groupées (Code pgfplots / Description attendue)}
\begin{popfigure}
\begin{tikzpicture}
\begin{axis}[
    ybar,
    width=\linewidth,
    height=10cm,
    enlargelimits=0.15,
    ylabel={Effectifs (milliers d'artisans)},
    symbolic x coords={Menuiserie, Couture, Mécanique, Alimentation, Art d'art, Autres},
    xtick=data,
    xticklabel style={rotate=45, anchor=east, font=\tiny},
    legend style={at={(0.5,-0.25)}, anchor=north, legend columns=-1, font=\footnotesize},
    ymin=0, ymax=130,
    bar width=6pt,
    nodes near coords, nodes near coords align={vertical}, every node near coord/.style={font=\tiny, rotate=90, anchor=west, yshift=2pt},
    popbar1/.style={fill=popblue},
    popbar2/.style={fill=poporange}
]
\addplot[popbar1] coordinates {(Menuiserie,85) (Couture,60) (Mécanique,35) (Alimentation,25) (Art d'art,18) (Autres,5)};
\addplot[popbar2] coordinates {(Menuiserie,115) (Couture,90) (Mécanique,58) (Alimentation,40) (Art d'art,24) (Autres,8)};
\legend{2015, 2022}
\end{axis}
\end{tikzpicture}
\legende{Évolution des effectifs artisans par secteur (2015 vs 2022). Source : MINPMEESA/INS. La menuiserie et la couture concentrent 61,2\% des effectifs en 2022.}
\end{popfigure}

\textbf{3. Commentaire structuré}
\textbf{Lecture globale :} Le graphique met en évidence une croissance généralisée des effectifs dans tous les secteurs entre 2015 et 2022, avec une hiérarchie stable dominée par la menuiserie et la couture.

\textbf{Analyse chiffrée :}
\begin{itemize}
    \item \textbf{Croissance absolue :} Le total gagne 107 000 artisans (+47\% en 7 ans).
    \item \textbf{Leaders :} Menuiserie (+30 000, +35\%), Couture (+30 000, +50\%), Mécanique (+23 000, +66\% - plus forte croissance relative).
    \item \textbf{Structure 2022 :} Menuiserie (34,3\%) + Couture (26,9\%) = \textbf{61,2\%} des effectifs. L'art d'art stagne (7,2\%).
\end{itemize}

\textbf{Interprétation géographique et économique :}
\begin{itemize}
    \item \textbf{Menuiserie :} Tirée par le boom immobilier urbain (portes, fenêtres, meubles) et la ressource bois (malgré raréfaction). Secteur « refuge » pour jeunes non scolarisés (apprentissage court).
    \item \textbf{Couture :} Féminisation forte, demande locale soutenue (tenues traditionnelles, uniformes scolaires), faible barrière à l'entrée (machine à coudre accessible).
    \item \textbf{Mécanique/Soudure :} Forte croissance (+66\%) liée au parc automobile vieillissant (réparation > neuf) et construction métallique (portails, toitures).
    \item \textbf{Art d'art :} Croissance faible (+33\% en 7 ans vs +47\% total). Freins : marché de niche (tourisme/export), transmission difficile, concurrence imports, manque de financement pour innover.
\end{itemize}
\textbf{Conclusion :} L'artisanat camerounais est un \cle{secteur de masse} (services/construction) plus que d'art. Les politiques publiques doivent cibler la modernisation des secteurs porteurs (menuiserie, mécanique) tout en préservant le secteur identitaire (art d'art) par des mesures spécifiques (label, tourisme).
\end{exoresolu}

\begin{attention}[Les 5 erreurs qui coûtent des points à l'examen]
\begin{enumerate}
    \item \textbf{Confondre Artisanat et Industrie / Commerce :} Ne pas citer « utilisation de machines lourdes », « production de masse standardisée », « achat-revente sans transformation » comme caractéristiques de l'artisanat. \textbf{Mot-clé :} \cle{Main-d'œuvre manuelle prépondérante} + \cle{Matières premières locales} + \cle{Savoir-faire traditionnel}.
    \item \textbf{Lister les problèmes sans les hiérarchiser ni les localiser :} Dire « il y a des problèmes de bois » au lieu de « \textbf{Pénurie de bois légal} dans la zone forestière (Sud/Est) due à l'exportation des grumes par les industriels, renchérissant le coût pour l'artisan de Foumban ». \textbf{Exigez :} Localisation + Mécanisme causal + Acteur concerné.
    \item \textbf{Proposer des solutions « magique » ou hors de portée des acteurs :} « L'État doit tout donner » / « Il faut interdire les imports ». \textbf{Bonne rédaction :} « L'État (MINPMEESA) facilite l'accès au crédit via le FNE ; La Chambre de Métiers organise la formation design ; Les artisans se structurent en coopératives pour mutualiser l'achat de matière première. » \cle{Responsabilité partagée}.
    \item \textbf{Oublier le vocabulaire technique de la cartographie/graphique :} « J'ai fait un dessin » au lieu de « J'ai réalisé un \textbf{croquis} / un \textbf{diagramme en barres groupées} ». Oubli du \textbf{Titre complet (Quoi/Où/Quand)}, de la \textbf{Légende hiérachisée}, de l'\textbf{Échelle}, de la \textbf{Flèche Nord}, de la \textbf{Source}. Zéro point pour la carte/graphique sans ces éléments.
    \item \textbf{Ne pas calculer / Ne pas utiliser les chiffres du sujet :} Dans l'exercice statistique, ne pas calculer le TVAM ou les pourcentages, ou écrire « ça a beaucoup augmenté » au lieu de « +47\% en 7 ans, soit un TVAM de 5,6\%/an ». \textbf{L'examen sanctionne l'absence de quantification.} Toujours : Calcul $\rightarrow$ Unité $\rightarrow$ Phrase d'interprétation.
\end{enumerate}
\end{attention}

\newpage

\coursec{Exercices}

\courssub{Je vérifie mes connaissances}

\begin{exercice}[QCM : notions de base sur l'artisanat]
Pour chaque question, une seule réponse est exacte. Recopie la lettre correspondante.
\begin{enumerate}
\item L'artisanat est une activité de production :
\begin{enumerate}
\item[a)] entièrement mécanisée et de masse ;
\item[b)] manuelle ou peu mécanisée, exercée par un artisan dans une petite unité ;
\item[c)] réservée aux grandes entreprises publiques.
\end{enumerate}
\item Le ministère en charge de la promotion de l'artisanat au Cameroun est :
\begin{enumerate}
\item[a)] le MINEPAT ;
\item[b)] le MINPMEESA ;
\item[c)] le MINADER.
\end{enumerate}
\item La sculpture sur bois, la poterie et le tissage appartiennent à la branche :
\begin{enumerate}
\item[a)] de l'artisanat d'art ;
\item[b)] de l'artisanat de service ;
\item[c)] de l'artisanat de production agroalimentaire.
\end{enumerate}
\item La microfinance aide les artisans surtout en leur donnant accès :
\begin{enumerate}
\item[a)] aux matières premières importées ;
\item[b)] au crédit et à l'épargne ;
\item[c)] aux débouchés extérieurs.
\end{enumerate}
\end{enumerate}
\end{exercice}

\begin{exercice}[Vrai ou faux, avec justification]
Indique si chaque affirmation est vraie ou fausse, puis justifie ta réponse en une ou deux phrases.
\begin{enumerate}
\item L'artisanat camerounais emploie une part importante de la population active, surtout dans le secteur informel.
\item L'artisanat se distingue de l'industrie par l'importance du capital et le nombre élevé de salariés.
\item Les activités artisanales sont réparties de manière parfaitement égale entre les dix régions du Cameroun.
\item Le manque de financement est l'un des principaux problèmes de l'artisanat camerounais.
\end{enumerate}
\end{exercice}

\begin{exercice}[Définitions]
Définis en deux ou trois phrases chacun des termes suivants :
\begin{enumerate}
\item artisanat ;
\item artisan ;
\item artisanat d'art ;
\item coopérative artisanale.
\end{enumerate}
\end{exercice}

\begin{exercice}[Calculs directs sur des données artisanales]
Une enquête menée dans une ville camerounaise recense \num{2400} artisans répartis ainsi : artisanat d'art : 720 ; artisanat de production (menuiserie, métallerie, agroalimentaire) : \num{1080} ; artisanat de service (coiffure, réparation, couture) : 600.
\begin{enumerate}
\item Calcule la part en pourcentage de chaque branche dans l'effectif total.
\item Représente ces parts par un diagramme circulaire (indique l'angle de chaque secteur).
\item L'année suivante, l'effectif total passe à \num{2760} artisans. Calcule le taux de variation en pourcentage.
\end{enumerate}
\end{exercice}

\courssub{Je m'entraîne}

\begin{exercice}[Branches et métiers]
\begin{enumerate}
\item Cite les trois grandes branches de l'artisanat camerounais et donne deux métiers pour chacune.
\item Classe les métiers suivants dans la branche qui convient : potier, mécanicien, sculpteur, coiffeur, vannerie, couturière, forgeron, réparateur de téléphones.
\item Explique en quelques lignes pourquoi l'artisanat de service se développe surtout dans les villes.
\end{enumerate}
\end{exercice}

\begin{exercice}[Importance économique et sociale de l'artisanat]
À partir de tes connaissances du cours :
\begin{enumerate}
\item Présente trois contributions de l'artisanat à l'économie camerounaise (emploi, revenus, utilisation des matières premières locales).
\item Montre en quatre lignes comment l'artisanat participe à la préservation de la culture camerounaise et au tourisme.
\item Rédige un court paragraphe justifiant l'intérêt de l'artisanat pour l'insertion professionnelle des jeunes diplômés de l'enseignement technique.
\end{enumerate}
\end{exercice}

\begin{exercice}[Les problèmes de l'artisanat]
Le tableau ci-dessous présente les difficultés déclarées par 100 artisans interrogés à Douala et Yaoundé (plusieurs réponses possibles).
\begin{center}
\begin{tabularx}{\linewidth}{|l|X|}\hline
\rowcolor{popblueL} Difficulté déclarée & Nombre d'artisans concernés \\ \hline
Difficulté d'accès au crédit bancaire & 68 \\ \hline
Concurrence des produits importés & 55 \\ \hline
Manque d'équipements modernes & 49 \\ \hline
Difficulté d'accès aux matières premières & 37 \\ \hline
Absence de formation professionnelle & 31 \\ \hline
\end{tabularx}
\end{center}
\begin{enumerate}
\item Classe ces difficultés de la plus fréquente à la moins fréquente.
\item Calcule le pourcentage d'artisans concernés par le problème de financement.
\item Regroupe ces difficultés en deux catégories : problèmes financiers et problèmes techniques. Justifie ton classement.
\item Propose une solution adaptée à la difficulté la plus fréquente.
\end{enumerate}
\end{exercice}

\begin{exercice}[Cartographie des pôles artisanaux]
Sur un fond de carte muet du Cameroun (reproduit ou fourni par l'enseignant) :
\begin{enumerate}
\item Place les cinq pôles artisanaux suivants avec un figuré ponctuel adapté : Foumban (sculpture et artisanat d'art), Maroua (poterie et cuir), Bamenda (vannerie et tissage), Douala (menuiserie et métallerie), Yaoundé (artisanat de service et d'art).
\item Ajoute un titre, une légende et l'orientation (flèche du nord).
\item Rédige un commentaire de cinq à huit lignes décrivant la répartition spatiale de ces pôles et proposant une explication.
\end{enumerate}
\end{exercice}

\courssub{Je me prépare au Bac}

\begin{exercice}[Type Probatoire : statistiques de l'emploi artisanal par région]
Le tableau ci-dessous donne le nombre estimé d'artisans (en milliers) dans six régions du Cameroun (données fictives d'entraînement).
\begin{center}
\begin{tabularx}{\linewidth}{|l|X|X|X|X|X|X|}\hline
\rowcolor{popblueL} Région & Ouest & Extrême-Nord & Littoral & Centre & Nord-Ouest & Est \\ \hline
Artisans (milliers) & 95 & 80 & 70 & 65 & 45 & 15 \\ \hline
\end{tabularx}
\end{center}
\begin{enumerate}
\item Calcule le nombre total d'artisans des six régions, puis la part en pourcentage de la région Ouest et de la région Est.
\item Construis un diagramme en bâtons représentant l'emploi artisanal par région (échelle conseillée : 1 cm pour 10 milliers d'artisans).
\item Dégage deux enseignements sur la répartition spatiale de l'artisanat au Cameroun.
\item Propose une explication au contraste observé entre l'Ouest et l'Est.
\end{enumerate}
\end{exercice}

\begin{exercice}[Type Baccalauréat : étude de cas d'une coopérative de potières de Maroua]
Une coopérative de 40 potières de Maroua (région de l'Extrême-Nord) produit chaque mois \num{1200} poteries vendues en moyenne \num{1500} FCFA l'unité sur les marchés locaux. Elle souhaite exporter vers le Tchad et le Nigeria. Elle dispose d'un local loué, de fours traditionnels et d'une épargne de \num{900000} FCFA. Un microcrédit de \num{3000000} FCFA lui serait accessible au taux annuel de 8 \%.
\begin{enumerate}
\item Calcule le chiffre d'affaires mensuel actuel de la coopérative.
\item Identifie quatre obstacles que la coopérative doit lever pour réussir son projet d'exportation (financement, équipement, qualité et normalisation, organisation commerciale).
\item Calcule le montant des intérêts annuels du microcrédit envisagé.
\item Rédige une note d'une quinzaine de lignes présentant trois solutions hiérarchisées pour permettre l'exportation, en précisant pour chacune les acteurs concernés (coopérative, microfinance, MINPMEESA, communes, partenaires privés).
\end{enumerate}
\end{exercice}

\begin{exercice}[Question de synthèse type Baccalauréat technique]
Sujet : \og L'artisanat camerounais est riche mais fragile. \fg

On te fournit deux documents :
\begin{itemize}
\item \textbf{Document 1} : carte schématique des pôles artisanaux du Cameroun (Foumban, Maroua, Bamenda, Douala, Yaoundé, Bafoussam), montrant une concentration dans l'ouest et les grandes villes.
\item \textbf{Document 2} : tableau indiquant que 68 \% des artisans n'ont pas accès au crédit bancaire, 55 \% subissent la concurrence des produits importés et 49 \% manquent d'équipements modernes.
\end{itemize}

Travail à faire :
\begin{enumerate}
\item À l'aide du document 1 et de tes connaissances, montre la richesse et la diversité de l'artisanat camerounais (branches, pôles, emplois, culture).
\item À l'aide du document 2, identifie trois problèmes majeurs qui expliquent la fragilité de l'artisanat camerounais.
\item Propose trois solutions justifiées, en précisant les acteurs qui doivent les mettre en œuvre (État, MINPMEESA, microfinances, coopératives, artisans).
\item Rédige une conclusion de quatre à six lignes répondant à la question : l'artisanat peut-il devenir un moteur du développement local au Cameroun ?
\end{enumerate}
\end{exercice}

\newpage

\coursec{Corrigés des exercices}

\begin{corrige}[Exercice 1 — QCM : notions de base sur l'artisanat]
\begin{enumerate}
\item \textbf{Réponse b}. L'artisanat est une activité de production \cle{manuelle ou peu mécanisée}, exercée par un artisan dans une petite unité de production (atelier, échoppe), avec peu de capital et peu de salariés. La production entièrement mécanisée et de masse (a) caractérise l'\emph{industrie}, et non l'artisanat.
\item \textbf{Réponse b}. Le \cle{MINPMEESA} (Ministère des Petites et Moyennes Entreprises, de l'Économie Sociale et de l'Artisanat) est le ministère en charge de la promotion de l'artisanat au Cameroun. Le MINEPAT s'occupe de la planification et le MINADER de l'agriculture.
\item \textbf{Réponse a}. La sculpture sur bois, la poterie et le tissage relèvent de l'\cle{artisanat d'art}, qui produit des objets à valeur esthétique et culturelle. L'artisanat de service rend des services (coiffure, réparation) et l'artisanat de production fabrique des biens utilitaires (meubles, produits agroalimentaires).
\item \textbf{Réponse b}. La \cle{microfinance} donne aux artisans l'accès au \emph{crédit} (petits prêts) et à l'\emph{épargne}, dont ils sont exclus par les banques classiques qui exigent des garanties qu'ils ne possèdent pas.
\end{enumerate}
\end{corrige}

\begin{corrige}[Exercice 2 — Vrai ou faux, avec justification]
\begin{enumerate}
\item \textbf{Vrai.} L'artisanat camerounais occupe une part importante de la population active : il constitue l'un des premiers employeurs du pays, mais la grande majorité de ces emplois relève du \cle{secteur informel} (activités non déclarées, sans contrat formel).
\item \textbf{Faux.} C'est l'inverse : l'artisanat se distingue de l'industrie par la \emph{faiblesse} du capital investi, le \emph{faible} nombre de salariés (souvent moins de 10) et le caractère manuel du travail. L'importance du capital et des effectifs salariés caractérise l'industrie.
\item \textbf{Faux.} La répartition est très \emph{inégale} : les activités artisanales se concentrent dans l'ouest du pays (Ouest, Nord-Ouest, Littoral, Centre) et dans les grandes villes (Douala, Yaoundé), tandis que des régions comme l'Est sont faiblement dotées.
\item \textbf{Vrai.} Le manque de financement est l'un des principaux obstacles : les artisans n'ont pas accès au crédit bancaire classique (absence de garanties), ce qui les empêche d'acheter des équipements modernes et des matières premières en quantité.
\end{enumerate}
\end{corrige}

\begin{corrige}[Exercice 3 — Définitions]
\begin{enumerate}
\item \textbf{Artisanat} : ensemble des activités de production de biens et de services exercées de manière manuelle ou peu mécanisée, dans de petites unités, par des artisans qui mettent en œuvre un savoir-faire souvent traditionnel. Il se distingue de l'industrie par le faible capital, le faible effectif et l'importance du geste manuel.
\item \textbf{Artisan} : travailleur qualifié qui fabrique, transforme ou répare des biens, ou rend des services, grâce à un savoir-faire technique manuel (sculpteur, potier, menuisier, coiffeur). Il travaille le plus souvent à son compte, dans un petit atelier, parfois avec quelques aides ou apprentis.
\item \textbf{Artisanat d'art} : branche de l'artisanat qui produit des objets à forte valeur esthétique et culturelle, destinés à la décoration, au culte ou au tourisme : sculpture sur bois, poterie décorative, vannerie, tissage, bijoux, bronzes de Foumban.
\item \textbf{Coopérative artisanale} : groupement d'artisans qui mettent en commun leurs moyens (local, équipements, achats de matières premières, commercialisation) pour réduire les coûts, accéder plus facilement au crédit et défendre leurs intérêts. Les décisions y sont prises collectivement et les bénéfices partagés entre les membres.
\end{enumerate}
\end{corrige}

\begin{corrige}[Exercice 4 — Calculs directs sur des données artisanales]
\begin{enumerate}
\item Effectif total : $\num{720} + \num{1080} + \num{600} = \num{2400}$ artisans (vérification conforme à l'énoncé).
\begin{itemize}
\item Artisanat d'art : $\dfrac{\num{720}}{\num{2400}} \times 100 = 30\,\%$.
\item Artisanat de production : $\dfrac{\num{1080}}{\num{2400}} \times 100 = 45\,\%$.
\item Artisanat de service : $\dfrac{\num{600}}{\num{2400}} \times 100 = 25\,\%$.
\end{itemize}
Vérification : $30 + 45 + 25 = 100\,\%$. Le calcul est juste.
\item L'angle de chaque secteur s'obtient en multipliant la part par $360$ degrés :
\begin{align*}
\text{Art} &: 0{,}30 \times 360 = 108\text{ degrés}\\
\text{Production} &: 0{,}45 \times 360 = 162\text{ degrés}\\
\text{Service} &: 0{,}25 \times 360 = 90\text{ degrés}
\end{align*}
Vérification : $108 + 162 + 90 = 360$ degrés. Le diagramme circulaire comporte donc trois secteurs de 108, 162 et 90 degrés, légendés avec les pourcentages.
\item Taux de variation :
\[ \text{Taux} = \frac{\num{2760} - \num{2400}}{\num{2400}} \times 100 = \frac{360}{\num{2400}} \times 100 = 15\,\%. \]
L'effectif des artisans a augmenté de 15\,\% en un an.
\end{enumerate}
\textbf{Points de vigilance} : ne pas oublier de vérifier que la somme des pourcentages vaut 100\,\% et celle des angles 360 degrés ; le taux de variation se calcule toujours par rapport à la valeur de \emph{départ}.
\end{corrige}

\begin{corrige}[Exercice 5 — Branches et métiers]
\begin{enumerate}
\item Les trois grandes branches de l'artisanat camerounais :
\begin{itemize}
\item \textbf{Artisanat d'art} : sculpture sur bois, poterie (aussi vannerie, tissage, bronzerie) ;
\item \textbf{Artisanat de production} : menuiserie, métallerie (aussi transformation agroalimentaire, briqueterie) ;
\item \textbf{Artisanat de service} : coiffure, réparation (aussi couture, mécanique).
\end{itemize}
\item Classement des métiers :
\begin{center}
\begin{tabularx}{\linewidth}{|l|X|}\hline
\rowcolor{popblueL} Branche & Métiers \\ \hline
Artisanat d'art & potier, sculpteur, vannerie, forgeron (objets d'art et utilitaires traditionnels) \\ \hline
Artisanat de production & forgeron (outils), couturière (production de vêtements) \\ \hline
Artisanat de service & mécanicien, coiffeur, réparateur de téléphones, couturière (retouches) \\ \hline
\end{tabularx}
\end{center}
Remarque : certains métiers sont frontaliers. Le forgeron relève de l'artisanat d'art ou de production selon ce qu'il fabrique ; la couturière produit (confection) ou rend un service (retouches). L'essentiel est de \emph{justifier} le classement.
\item L'artisanat de service se développe surtout en ville parce que les villes concentrent une \cle{demande} forte et diversifiée : population nombreuse, véhicules et appareils à réparer, besoins de coiffure et d'habillement. Le citadin dispose en outre d'un revenu monétaire qui lui permet de payer ces services, et la densité urbaine assure à l'artisan une clientèle régulière, contrairement à la campagne où la demande est dispersée.
\end{enumerate}
\end{corrige}

\begin{corrige}[Exercice 6 — Importance économique et sociale de l'artisanat]
\begin{enumerate}
\item Trois contributions économiques :
\begin{itemize}
\item \textbf{Emploi} : l'artisanat est l'un des premiers pourvoyeurs d'emplois du pays, notamment pour les jeunes et les personnes sans diplôme, absorbant une large part de la main-d'œuvre du secteur informel ;
\item \textbf{Revenus} : il procure des revenus monétaires à des millions de ménages, en ville comme à la campagne, et contribue ainsi à la lutte contre la pauvreté ;
\item \textbf{Valorisation des matières premières locales} : bois, argile, cuir, fibres végétales (raffia, raphia), produits agricoles sont transformés sur place, ce qui crée de la valeur ajoutée locale et réduit les importations.
\end{itemize}
\item L'artisanat conserve et transmet le patrimoine culturel camerounais : les motifs sculptés de Foumban, les poteries de l'Extrême-Nord, les tissus du Nord-Ouest expriment l'identité des peuples. Ces objets attirent les touristes (villages artisanaux, marchés d'art, musées), ce qui génère des devises et fait connaître la culture camerounaise à l'étranger.
\item Pour les jeunes diplômés de l'enseignement technique, l'artisanat offre une voie concrète d'\cle{insertion professionnelle} : il demande peu de capital de départ, valorise directement les compétences acquises (menuiserie, mécanique, électricité, couture) et permet de créer rapidement sa propre petite entreprise. Dans un contexte de chômage des jeunes et de rareté des emplois salariés, l'auto-emploi artisanal constitue une solution réaliste, d'autant que l'État et la microfinance accompagnent de plus en plus les porteurs de projets.
\end{enumerate}
\end{corrige}

\begin{corrige}[Exercice 7 — Les problèmes de l'artisanat]
\begin{enumerate}
\item Classement de la plus fréquente à la moins fréquente :
\begin{enumerate}
\item Difficulté d'accès au crédit bancaire (68) ;
\item Concurrence des produits importés (55) ;
\item Manque d'équipements modernes (49) ;
\item Difficulté d'accès aux matières premières (37) ;
\item Absence de formation professionnelle (31).
\end{enumerate}
\item Pourcentage concerné par le problème de financement : $\dfrac{68}{100} \times 100 = 68\,\%$. Soit 68\,\% des artisans interrogés (plus des deux tiers).
\item Regroupement en deux catégories :
\begin{itemize}
\item \textbf{Problèmes financiers} : difficulté d'accès au crédit bancaire (68) et concurrence des produits importés (55). Le crédit est par nature un problème de financement ; la concurrence des importations en est une conséquence indirecte : faute de capitaux, l'artisan produit cher et en petite quantité, et ne peut pas rivaliser en prix avec les produits industriels importés.
\item \textbf{Problèmes techniques} : manque d'équipements modernes (49), difficulté d'accès aux matières premières (37) et absence de formation professionnelle (31). Ils concernent directement les moyens et les savoir-faire de production.
\end{itemize}
(On accepte aussi un classement plaçant la concurrence des importations dans les problèmes commerciaux, à condition de le justifier.)
\item Solution à la difficulté la plus fréquente (accès au crédit) : développer le recours à la \cle{microfinance} (MC2, CamCCUL, coopératives d'épargne et de crédit) qui accorde de petits prêts sans garanties bancaires lourdes, et encourager les artisans à se regrouper en coopératives ou en \emph{tontines} pour mutualiser l'épargne et servir de caution collective. L'État, à travers le MINPMEESA, peut en outre créer ou renforcer des fonds de garantie et des subventions d'équipement.
\end{enumerate}
\end{corrige}

\begin{corrige}[Exercice 8 — Cartographie des pôles artisanaux]
\begin{enumerate}
\item Figurés ponctuels conseillés (un symbole par spécialité dominante, ou un point unique avec étiquette) :
\begin{itemize}
\item \textbf{Foumban} (Ouest) : sculpture, bronzes, artisanat d'art ;
\item \textbf{Maroua} (Extrême-Nord) : poterie, cuir (tannerie) ;
\item \textbf{Bamenda} (Nord-Ouest) : vannerie, tissage ;
\item \textbf{Douala} (Littoral) : menuiserie, métallerie ;
\item \textbf{Yaoundé} (Centre) : artisanat de service et d'art.
\end{itemize}
Sur la carte, Foumban se place dans l'ouest montagneux, Maroua à l'extrême nord (plaine du Diamaré), Bamenda au nord-ouest (hautes terres), Douala sur la côte et Yaoundé au centre-sud.
\item Éléments obligatoires : un \textbf{titre} (ex. \og Les principaux pôles artisanaux du Cameroun \fg), une \textbf{légende} expliquant les figurés, et la \textbf{flèche du nord}. Une échelle indicative est un plus.
\item \textbf{Commentaire attendu (corrigé type)} : la répartition des pôles artisanaux est très inégale. Quatre des cinq pôles se situent dans la moitié ouest du pays : Foumban et Bamenda dans les hautes terres de l'Ouest et du Nord-Ouest, Douala et Yaoundé dans les deux plus grandes villes. Seul Maroua représente le grand nord. Cette concentration s'explique d'abord par la densité de population de l'ouest, qui offre main-d'œuvre et clientèle ; ensuite par la longue tradition artisanale des royaumes de l'Ouest (Bamoun, Bamiléké), où la sculpture et le tissage sont liés à la chefferie ; enfin par le poids des métropoles Douala et Yaoundé, pôles de consommation et d'équipement. L'Est et le Sud, faiblement peuplés et enclavés, restent en marge de cette dynamique.
\end{enumerate}
\textbf{Points de vigilance} : une carte sans titre, sans légende ou sans orientation est sanctionnée ; le commentaire doit décrire (constats localisés) \emph{puis} expliquer, et non se contenter d'énumérer les villes.
\end{corrige}

\begin{corrige}[Exercice 9 — Type Probatoire : statistiques de l'emploi artisanal par région]
\begin{enumerate}
\item Total : $95 + 80 + 70 + 65 + 45 + 15 = 370$ milliers d'artisans.
\begin{itemize}
\item Part de l'Ouest : $\dfrac{95}{370} \times 100 \approx 25{,}7\,\%$ (environ 26\,\%).
\item Part de l'Est : $\dfrac{15}{370} \times 100 \approx 4{,}1\,\%$ (environ 4\,\%).
\end{itemize}
Vérification : $25{,}7 + 4{,}1 \approx 29{,}8\,\%$ pour les deux régions, ce qui est cohérent avec $(95+15)/370 = 110/370 \approx 29{,}7\,\%$ (écart dû aux arrondis).
\item Diagramme en bâtons : axe horizontal = régions, axe vertical = effectifs en milliers, échelle 1 cm pour 10 milliers. Hauteurs des bâtons : Ouest 9,5 cm ; Extrême-Nord 8 cm ; Littoral 7 cm ; Centre 6,5 cm ; Nord-Ouest 4,5 cm ; Est 1,5 cm. Chaque bâton est séparé du suivant (variable qualitative), le graphique porte un titre et les axes sont légendés.
\item Deux enseignements :
\begin{itemize}
\item L'artisanat est \textbf{concentré dans l'ouest du pays} : l'Ouest, l'Extrême-Nord, le Littoral et le Centre totalisent $95+80+70+65 = 310$ milliers, soit $\dfrac{310}{370} \times 100 \approx 83{,}8\,\%$ de l'effectif des six régions ;
\item Il existe de \textbf{forts contrastes régionaux} : la région Ouest compte à elle seule plus de six fois plus d'artisans que l'Est ($95/15 \approx 6{,}3$).
\end{itemize}
\item Explication du contraste Ouest/Est : l'Ouest est une région de hautes terres très densément peuplée, dotée d'une tradition artisanale séculaire (royaumes bamiléké et bamoun : sculpture, tissage, vannerie) et de marchés dynamiques ; l'Est est au contraire faiblement peuplé, enclavé (mauvaises routes, éloignement des grands marchés), couvert de forêt dense, et sa population vit davantage de l'agriculture vivrière et de la cueillette que de l'artisanat marchand.
\end{enumerate}
\textbf{Barème indicatif (sur 20)} : calculs exacts et vérifiés (6 pts) ; diagramme correct avec titre, axes légendés et échelle respectée (6 pts) ; deux enseignements chiffrés (4 pts) ; explication géographique argumentée (4 pts).
\textbf{Points de vigilance} : les bâtons d'un diagramme en bâtons sont \emph{séparés} (ne pas les coller comme dans un histogramme) ; tout enseignement dégagé doit être appuyé par un chiffre calculé.
\end{corrige}

\begin{corrige}[Exercice 10 — Type Baccalauréat : étude de cas d'une coopérative de potières de Maroua]
\begin{enumerate}
\item Chiffre d'affaires mensuel actuel :
\[ \num{1200} \times \num{1500} = \num{1800000} \text{ FCFA par mois.} \]
\item Quatre obstacles à lever :
\begin{itemize}
\item \textbf{Financement} : l'épargne de \num{900000} FCFA est insuffisante pour moderniser l'outil et financer l'exportation (transport, formalités douanières) ;
\item \textbf{Équipement} : les fours traditionnels donnent une cuisson irrégulière et une production limitée ; il faut des fours améliorés pour produire en quantité et en qualité constantes ;
\item \textbf{Qualité et normalisation} : exporter exige des produits homogènes, résistants au transport, répondant à des normes (dimensions, solidité, emballage) ;
\item \textbf{Organisation commerciale} : la coopérative ne connaît pas les marchés tchadien et nigérian ; il lui faut des circuits de distribution, des contacts commerciaux et une logistique transfrontalière.
\end{itemize}
\item Intérêts annuels du microcrédit :
\[ I = \num{3000000} \times \frac{8}{100} = \num{240000} \text{ FCFA par an.} \]
Le montant total à rembourser en un an serait donc de $\num{3000000} + \num{240000} = \num{3240000}$ FCFA.
\item \textbf{Note rédigée (corrigé type)} : pour réussir son exportation, la coopérative doit hiérarchiser trois solutions. \emph{Premièrement}, sécuriser le financement : contracter le microcrédit de \num{3000000} FCFA auprès d'une microfinance (MC2, CamCCUL) en complément de l'épargne de \num{900000} FCFA, et solliciter les fonds d'appui du MINPMEESA ; la coopérative devra présenter un plan d'affaires crédible. \emph{Deuxièmement}, moderniser l'équipement et la qualité : acheter des fours améliorés et un matériel d'emballage, et faire former les 40 potières aux normes de production (stages organisés avec le MINPMEESA et les centres de formation), afin de garantir des poteries homogènes et transportables. \emph{Troisièmement}, organiser la commercialisation : nouer des partenariats avec des commerçants de Moundou et de Maiduguri, s'appuyer sur la commune de Maroua et les partenaires privés (transporteurs) pour la logistique, et participer aux foires artisanales régionales pour se faire connaître. La coopérative pilote l'ensemble, la microfinance apporte le crédit, l'État (MINPMEESA) la formation et la garantie, les communes et partenaires privés l'appui logistique et commercial.
\end{enumerate}
\textbf{Barème indicatif (sur 20)} : calcul du chiffre d'affaires (3 pts) ; quatre obstacles correctement identifiés (6 pts) ; calcul des intérêts (3 pts) ; note structurée avec trois solutions hiérarchisées et acteurs cités (8 pts).
\textbf{Points de vigilance} : distinguer chiffre d'affaires et bénéfice (les coûts ne sont pas fournis, on ne calcule donc pas de bénéfice) ; dans la note, chaque solution doit nommer explicitement les acteurs ; le taux de 8\,\% est \emph{annuel}, ne pas l'appliquer au mois.
\end{corrige}

\begin{corrige}[Exercice 11 — Question de synthèse type Baccalauréat technique]
\begin{enumerate}
\item \textbf{La richesse de l'artisanat camerounais.} Le document 1 montre une diversité de pôles spécialisés : Foumban (sculpture, bronzes), Maroua (poterie, cuir), Bamenda (vannerie, tissage), Douala (menuiserie, métallerie), Yaoundé (art et services), Bafoussam (artisanat d'art). Ces pôles illustrent les trois branches de l'artisanat : artisanat d'art, artisanat de production et artisanat de service. Cet artisanat emploie une large part de la population active, surtout dans le secteur informel, valorise les matières premières locales (bois, argile, cuir, fibres) et constitue un patrimoine culturel vivant, hérité des royaumes de l'Ouest et des traditions du grand nord, qui attire en outre les touristes.
\item \textbf{La fragilité de l'artisanat.} Le document 2 révèle trois problèmes majeurs : d'abord le \emph{manque de financement} (68\,\% des artisans sans accès au crédit bancaire), qui bloque tout investissement ; ensuite la \emph{concurrence des produits importés} (55\,\%), souvent moins chers et mieux présentés ; enfin l'\emph{archaïsme de l'outil de production} (49\,\% sans équipements modernes), qui limite la quantité et la qualité. À cela s'ajoute, d'après le document 1, une répartition spatiale déséquilibrée qui marginalise des régions entières.
\item \textbf{Trois solutions justifiées} :
\begin{itemize}
\item \emph{Faciliter l'accès au financement} : développer la microfinance et les fonds de garantie ; acteurs : microfinances (MC2, CamCCUL), État et MINPMEESA, artisans regroupés en coopératives et tontines ;
\item \emph{Moderniser et former} : subventionner l'équipement des ateliers et organiser la formation professionnelle et l'apprentissage ; acteurs : État (MINPMEESA), communes, maîtres artisans ;
\item \emph{Protéger et promouvoir la production locale} : contrôler les importations concurrentes, labelliser les produits camerounais, organiser foires et marchés d'art ; acteurs : État (douanes, MINPMEESA), coopératives, partenaires privés et touristiques.
\end{itemize}
\item \textbf{Conclusion (corrigé type)} : l'artisanat camerounais possède des atouts réels : savoir-faire reconnus, matières premières abondantes, main-d'œuvre disponible et forte demande locale et touristique. Mais sa fragilité financière et technique limite aujourd'hui son essor. Si les pouvoirs publics, la microfinance et les artisans eux-mêmes mettent en œuvre les solutions proposées — crédit, formation, équipement, organisation en coopératives —, l'artisanat peut devenir un véritable moteur du développement local : créateur d'emplois et de revenus, il fixerait les populations, valoriserait les ressources de chaque région et rayonnerait culturellement à l'étranger.
\end{enumerate}
\textbf{Barème indicatif (sur 20)} : richesse et diversité démontrées à partir du document 1 et des connaissances (5 pts) ; trois problèmes tirés du document 2 et chiffrés (5 pts) ; trois solutions justifiées avec acteurs précisés (6 pts) ; conclusion argumentée répondant à la question posée (4 pts).
\textbf{Points de vigilance} : mobiliser \emph{explicitement} les deux documents (citer les chiffres 68\,\%, 55\,\%, 49\,\% et les pôles de la carte) ; ne pas confondre les acteurs (le MINPMEESA encadre et promeut, il ne fabrique pas) ; la conclusion doit répondre à la question posée par un jugement nuancé, pas par une simple répétition de l'introduction.
\end{corrige}

\newpage

\coursec{Fiche bilan}

\begin{popfigure}
\begin{tikzpicture}[mindmap, grow cyclic, every node/.style={concept, font=\small},
concept color=popblue!60,
level 1/.append style={level distance=3.4cm, sibling angle=60},
level 2/.append style={level distance=2.1cm, sibling angle=45, font=\scriptsize}]
\node[concept, concept color=popink!80, text=white, font=\small\bfseries]{L'artisanat au Cameroun}
child[concept color=poporange!70]{ node[concept]{Définition et caractéristiques}
  child{ node[concept]{Activité manuelle, petit atelier} }
  child{ node[concept]{Savoir-faire traditionnel} }
}
child[concept color=popgreen!70]{ node[concept]{Branches}
  child{ node[concept]{Bois, vannerie, poterie} }
  child{ node[concept]{Textile, bijouterie, mécanique} }
}
child[concept color=poppurple!70]{ node[concept]{Répartition spatiale}
  child{ node[concept]{Foumban : arts du bois} }
  child{ node[concept]{Grand Nord : cuir, vannerie} }
}
child[concept color=popgold!80]{ node[concept]{Importance}
  child{ node[concept]{Emplois, revenus} }
  child{ node[concept]{Culture, tourisme} }
}
child[concept color=poppink!70]{ node[concept]{Problèmes}
  child{ node[concept]{Financement, équipement} }
  child{ node[concept]{Concurrence, formation} }
}
child[concept color=popblue!50]{ node[concept]{Solutions (TD2)}
  child{ node[concept]{Crédit, coopératives} }
  child{ node[concept]{Formation, promotion} }
};
\end{tikzpicture}
\legende{Carte mentale de la leçon : les six idées forces sur l'artisanat camerounais.}
\end{popfigure}

\begin{bilan}
\textbf{Définitions clés}
\begin{itemize}
\item \cle{Artisanat} : activité de production ou de service exercée à petite échelle, dans un atelier, à l'aide d'un \cle{savoir-faire} essentiellement manuel, souvent transmis de génération en génération.
\item \cle{Artisan} : travailleur qualifié qui conçoit et réalise lui-même des objets (sculpteur, potier, forgeron, couturier, mécanicien).
\item \cle{Coopérative artisanale} : groupement d'artisans qui mutualisent moyens de production, achats de matières premières et ventes.
\end{itemize}

\textbf{Branches et répartition spatiale (à connaître par cœur)}
\begin{center}
\begin{tabularx}{\linewidth}{|l|X|X|}\hline
\rowcolor{popblueL} \textbf{Branche} & \textbf{Produits} & \textbf{Zones principales} \\ \hline
Sculpture sur bois & Statues, masques, meubles & Foumban (Ouest), régions du Centre et du Sud \\ \hline
Vannerie & Paniers, nattes, chapeaux & Grand Nord (Extrême-Nord, Nord, Adamaoua) \\ \hline
Poterie & Jarres, marmites, objets décoratifs & Extrême-Nord, Ouest \\ \hline
Cuir et peaux & Sacs, sandales, poufs & Maroua, Ngaoundéré \\ \hline
Textile et habillement & Ndop, tissus, couture & Foumban, Douala, Yaoundé \\ \hline
Petits métiers techniques & Mécanique, soudure, bijouterie & Grandes villes (Douala, Yaoundé) \\ \hline
\end{tabularx}
\end{center}

\textbf{Importance de l'artisanat}
\begin{itemize}
\item Source d'\cle{emplois} et de revenus, surtout dans le secteur informel (villes et campagnes).
\item Valorisation des \cle{matières premières locales} (bois, argile, cuir, fibres).
\item Préservation du \cle{patrimoine culturel} et attraction touristique (marchés artisanaux de Foumban).
\item Contribution à l'économie locale et à la réduction des importations.
\end{itemize}

\textbf{Problèmes et solutions (TD2)}
\begin{itemize}
\item \emph{Problèmes} : difficulté d'accès au \cle{financement} ; outillage rudimentaire ; concurrence des produits importés ; formation insuffisante ; faible organisation commerciale.
\item \emph{Solutions} : microcrédit et coopératives ; formation professionnelle et apprentissage ; équipement moderne des ateliers ; promotion (foires, labels, tourisme) ; appui de l'État et des collectivités.
\end{itemize}

\textbf{Méthodes express}
\begin{itemize}
\item \emph{Définir} : activité + échelle (petit atelier) + savoir-faire manuel.
\item \emph{Localiser} : associer toujours une branche à une région précise (ex. vannerie $\rightarrow$ Grand Nord).
\item \emph{Problème $\rightarrow$ solution} : présenter les solutions en miroir des problèmes, avec un exemple camerounais.
\item \emph{Cartographier} : croquis = titre, légende, flèche du nord, symboles simples par branche.
\end{itemize}
\end{bilan}

\begin{aretenir}[Checklist avant le Bac]
\begin{itemize}
\item $\square$ Je sais définir l'artisanat et citer ses caractéristiques (petite échelle, travail manuel, savoir-faire).
\item $\square$ Je connais au moins 5 branches artisanales camerounaises.
\item $\square$ J'associe chaque branche à sa zone de production (Foumban, Grand Nord, grandes villes).
\item $\square$ Je peux citer 3 produits artisanaux typiques (masque, natte, jarre).
\item $\square$ J'explique l'importance économique de l'artisanat (emplois, revenus, secteur informel).
\item $\square$ J'explique son importance culturelle et touristique.
\item $\square$ Je cite au moins 4 problèmes de l'artisanat camerounais.
\item $\square$ Je propose une solution adaptée à chaque problème (crédit, formation, coopératives, promotion).
\item $\square$ Je sais réaliser un croquis simple avec titre, légende et orientation.
\item $\square$ Je rédige une réponse structurée : définition, développement avec exemples, conclusion.
\end{itemize}
\end{aretenir}

\findecours

\end{document}
